\documentclass[12pt, a4paper]{article}

% ----- PACKAGES -----
\usepackage[utf8]{inputenc}
\usepackage{geometry}
\usepackage{multicol}
\usepackage{xcolor} % For the gray numbers
\newcounter{chaupainum} % Define the counter once here
\geometry{top=1in, bottom=1in, left=1in, right=1in}
\usepackage{fontspec}
\usepackage{polyglossia}
\usepackage{tabularx}
\usepackage{booktabs}

\definecolor{laghuorange}{HTML}{D95B00} % Darker orange for text readability
\definecolor{gurublue}{HTML}{0056B3} % Darker blue for text readability
\newcommand{\laghu}[1]{{\color{laghuorange}\protect\deva{#1}}}
\newcommand{\guru}[1]{{\color{gurublue}\protect\deva{#1}}}
\usepackage[colorlinks=true, linkcolor=blue]{hyperref}
\usepackage{bookmark}
\usepackage{tcolorbox} % For the "Gotchas" and "Memory Anchors"
\usepackage{fontawesome5} % For the Gotcha Lightbulb icon
\newcommand{\gotchaicon}{\textcolor{laghuorange}{\faLightbulb}}
\newcommand{\glossaryicon}{\textcolor{gurublue}{\faBookOpen}}

% ----- LANGUAGE & FONT SETUP FOR XELATEX -----
\setdefaultlanguage{english}

% You may need to change these font names based on your local system fonts.
% If using Overleaf, "Noto Serif Devanagari" and "Noto Serif Telugu" work beautifully.
\newfontfamily\devanagarifont[Script=Devanagari]{Noto Serif Devanagari}
\newfontfamily\telugufont[Script=Telugu]{Noto Sans Telugu}

% Custom commands to easily wrap text in the correct script
\newcommand{\deva}[1]{{\devanagarifont #1}}
\newcommand{\telu}[1]{{\telugufont #1}}

% Custom Box styling for notes
\tcbset{
    gotchabox/.style={colback=red!5!white, colframe=red!75!black, title=Linguistic Gotchas, fonttitle=\bfseries},
    anchorbox/.style={colback=blue!5!white, colframe=blue!75!black, title=Story, fonttitle=\bfseries},
    poetrybox/.style={colback=green!5!white, colframe=green!50!black, title=Poetic Beauty \& Praasa, fonttitle=\bfseries}
}

\setlength{\parindent}{0pt}
\setlength{\parskip}{0.5em}

% --- Custom Commands for Chalisa Recitation ---
\newcounter{dohacount}
\newcounter{chaupaicount}

\newcommand{\resetcounters}{%
  \setcounter{dohacount}{0}%
  \setcounter{chaupaicount}{0}%
}

\newcommand{\doha}[1]{%
  \stepcounter{dohacount}%
  \noindent\textbf{Doha \thedohacount:} \\
  {\large \linespread{1.5}\selectfont #1 \par}%
  \vspace{0.8em}%
}

\newcommand{\chaupai}[1]{%
  \stepcounter{chaupaicount}%
  \noindent\textbf{\thechaupaicount.} {\large \linespread{1.5}\selectfont #1 \par}%
  \vspace{0.8em}%
}

\begin{document}

\title{
    \textbf{ \deva{श्री हनुमान चालीसा}}\\ 
    \telu{శ్రీ హనుమాన చాలీసా} \\  
    \textit{śrī hanumān cālīsā} \\  
    \Large A Linguistic \& Structural Study Guide
}
\author{Compiled for easy learning}
\date{\today}
\maketitle

\tableofcontents
\newpage

% ==========================================
\section*{Introduction: Learning the Hanuman Chalisa}
\addcontentsline{toc}{section}{Introduction: Learning the Hanuman Chalisa}
% ==========================================

\subsection*{Why learn and recite?} 
There are significant cognitive and psychological advantages to memorizing and reciting complex works. Whether it is a song, a story, or a speech, the act of recollection strengthens memory and focus. The \textit{Hanuman Chalisa} is particularly rewarding because it is steeped in history, based on an ancient epic, and specifically designed with a rhythmic structure that makes it accessible and fun to learn. While the learning journey itself is rewarding, daily recitation compounds these benefits.

\subsection*{Why the Hanuman Chalisa?}
Beyond the traditional \deva{फलश्रुति} (Phalashruti), the benefits listed by the poet Tulsidas at the end of the work, the Chalisa is an amazingly beautiful composition. Its sequence of thoughts and concepts is intentionally designed to guide the practitioner into a positive, grounded state of mind.

\subsection*{Strategies for Effective Learning}
\begin{enumerate}
    \item \textbf{Anchor the Meaning:} Start by reviewing the \deva{प्रतिपदार्थ} (word-for-word) tables to understand exactly what you are saying.
    \item \textbf{Listen and Absorb:} Get familiar with the phonetics by listening to the companion audio or other traditional renditions. 
    \item \textbf{Follow the Map:} Practice vocalizing words using the \textbf{Visual Rhythm Maps} to master the timing.
    \item \textbf{Multimodal Practice:} Writing or typing the verses while vocalizing them helps bridge the gap between sight and sound.
    \item \textbf{Spaced Recall:} Allocate the exact same time every day to spend 5 to 10 minutes focused on learning new verses. Later during the day, whenever you have a free minute or two, actively try to recite those verses from memory to the best of your ability. At the end of the day, do one final revision. Active recall builds the strongest neural pathways.
    \item \textbf{Cumulative Learning:} Master one Semantic Group at a time, always practicing from the very beginning to build "muscle memory."
    \item \textbf{Set the Tone:} Always begin with the invocatory \textit{Dohas} to center your mind before the main recitation.
\end{enumerate}

\hr % Optional horizontal rule if defined in your preamble

\subsection*{The Mathematical Heartbeat}
Before chanting, it is beneficial to understand the mathematical heartbeat of the text. Unlike classical Sanskrit or Telugu meters which follow strict syllable counts (\deva{अक्षरछन्दस्}), the \textit{Hanuman Chalisa} is an example of \deva{मात्राछन्दस्}, a system based on \textbf{beats} (matras).

\textbf{Syllable Weights \& Visual Rhythm Maps:}
\begin{itemize}
    \item \textbf{\deva{लघु} (Laghu : Short, light) = 1 Beat (\textcolor{laghuorange}{Orange}):} Short vowels like a, i, u (\deva{अ, इ, उ} / \telu{అ, ఇ, ఉ}).
    \item \textbf{\deva{गुरु} (Guru : Long, heavy) = 2 Beats (\textcolor{gurublue}{Blue}):} Long vowels like aa, ee, oo, e, o (\deva{आ, ई, ऊ, ए, ओ} / \telu{ఆ, ఈ, ఊ, ఏ, ఓ}), or any short vowel followed by a conjunct consonant or nasal dot.
    \item \textbf{Visual Breakdown:} Every line in this guide includes a Rhythm Map where individual syllables are colored (\textcolor{gurublue}{\textbf{Guru}} / \textcolor{laghuorange}{\textbf{Laghu}}) so you can track the melody intuitively without manual counting.
\end{itemize}

\subsection*{The \deva{चौपाई} (Chaupai): The 16-Beat March}
A Chaupai is a \deva{सममात्रिक छन्द} (Even Meter). Every single line must contain exactly \textbf{16 beats}. 
\begin{itemize}
    \item \textbf{Example:} \deva{ज(1) य(1) ह(1) नु(1) मा(2) न(1) ज्ञा(2) न(1) गु(1) न(1) सा(2) ग(1) र(1)} = 16 Beats
\end{itemize}

\subsection*{The \deva{दोहा} (Doha): The 13-11 Sway}
A Doha is an \deva{अर्धसममात्रिक छन्द} (Half-Even Meter). It consists of two lines, each split by a \deva{यति} (Yati / breath pause).
\begin{itemize}
    \item \textbf{First Half:} Exactly \textbf{13 beats}.
    \item \textbf{Second Half:} Exactly \textbf{11 beats}. These must end in a rhyming \textbf{Guru-Laghu} (2-1) "drop" (e.g., \deva{सुधारि} $\rightarrow$ \deva{धा(2) रि(1)}).
\end{itemize}

\begin{tcolorbox}[gotchabox, title=Master Linguistic Rule Sheet]
\textbf{Sanskrit to Awadhi Phonetic Shifts:}
\begin{itemize}
    \item \textbf{The \deva{व} $\rightarrow$ \deva{ब} Shift (v $\rightarrow$ b):} Word-initial 'v' hardens into 'b' (\deva{विद्या} $\rightarrow$ \deva{बिद्या}).
    \item \textbf{The \deva{श/ष} $\rightarrow$ \deva{स} Shift (sh $\rightarrow$ s):} Palatal 'sh' softens to dental 's' (\deva{कपीश} $\rightarrow$ \deva{कपीस}).
    \item \textbf{The \deva{य} $\rightarrow$ \deva{ज} Shift (y $\rightarrow$ j):} Soft 'y' hardens to 'j' (\deva{युग} $\rightarrow$ \deva{जुग}).
    \item \textbf{The \deva{क्ष} Shift (ksh $\rightarrow$ kh / cch):} The harsh Sanskrit \deva{क्ष} simplifies into \deva{ख} (\deva{लक्ष्मण} $\rightarrow$ \deva{लखन}) or \deva{च्छ} (\deva{रक्षक} $\rightarrow$ \deva{रच्छक}).
    \item \textbf{The \deva{ण} $\rightarrow$ \deva{न} Shift (ṇ $\rightarrow$ n):} Retroflex 'n' flattens into dental 'n' (\deva{गुण} $\rightarrow$ \deva{गुन}).
\end{itemize}

\textbf{Rhythm \& Meter Mechanics:}
\begin{itemize}
    \item \textbf{Schwa Holding:} Final consonants are not "clipped" dead as in modern Hindi (\deva{नन्दन्}), nor fully vocalized as in Sanskrit (\deva{नन्द-न}). They are held for one beat to satisfy the 16-beat math.
    \item \textbf{Swara-Bhakti (Vowel Insertion):} Clusters are broken by a neutralizing 'a' to aid flow (\deva{वज्र} $\rightarrow$ \deva{बजर}).
\end{itemize}
\end{tcolorbox}

\subsection*{How to Read the Translation Tables}
\begin{itemize}
    \item \textbf{Column 1 (Awadhi):} The word as it appears in the text (Devanagari and Telugu). 
    \item \textbf{Column 2 (Sanskrit):} The classical root word from which the Awadhi term evolved.
    \item \textbf{Column 3 (English):} The explicit definition.
    \item \textbf{Column 4 (Notes):} Grammatical info. If you see a \textbf{Lightbulb} (\gotchaicon), the word follows a shift from the Master Rule Sheet. A \textbf{Book} (\glossaryicon) indicates a detailed theological entry in the \textbf{Glossary Appendix}.
\end{itemize}

\subsection*{The IAST Transliteration Scheme}
To aid pronunciation for those unfamiliar with the Devanagari script, this guide utilizes the \textbf{International Alphabet of Sanskrit Transliteration} (IAST) with a rigorous phonological approach suited for the Awadhi language.
\begin{itemize}
    \item \textbf{The Anusvara (The Nasal Dot):} In Sanskrit and Awadhi orthography, a dot above a letter (\deva{ं}) implies a nasal sound that assimilates to the consonant immediately following it. Rather than using a generic `ṃ' for every dot, our IAST rendering uses the exact strict linguistic "class nasal" for phonetic precision. For example, before a 'k' (\deva{क}), the dot becomes a velar 'ṅ' (\deva{संकर} $\rightarrow$ \textit{saṅkara}); before a 'c' (\deva{च}), it becomes a palatal 'ñ' (\deva{कंचन} $\rightarrow$ \textit{kañcana}); and before a 'd' (\deva{द}), it becomes a dental 'n' (\deva{नंदन} $\rightarrow$ \textit{nandana}). 
    \item \textbf{The Chandrabindu (The Nasalized Vowel):} A moon-and-dot symbol (\deva{ँ}) represents the nasalization of the vowel itself, rather than a distinct nasal consonant. 
    This is represented by placing a tilde directly over the transliterated vowel (e.g., \deva{तिहुँ} $\rightarrow$ \textit{tihũ} ; \deva{काँपै} $\rightarrow$ \textit{kā̃pai}).  It is read as if that vowel is pronounced normally, but resonating with the help of the nose too. 
\end{itemize}

\newpage

\subsection*{Ramayana Context Boxes}
Beneath the translation of every verse, you will find an anecdote box. These boxes highlight practical, character-building lessons drawn directly from Hanuman's actions. At the bottom right of each box, a specific citation is provided from the authoritative \textit{Valmiki Ramayana}. These citations are formatted as \textbf{Kanda, Sarga:Shloka}. For example, \textit{Sundara Kanda, 1:170} refers to the Sundara Kanda, Sarga (Chapter) 1, Shloka (Verse) 170.

\vspace{1em}

\subsection*{The Semantic Learning Plan}
We have divided the Chalisa into eight semantic groups to facilitate mastery:
\begin{description}
    \item[Group 1: The Invocation] (Dohas 1--2). Cleansing the mind and setting intent.
    \item[Group 2: Appearance] (Chaupais 1--6). Introduction to Hanuman's physical and spiritual form.
    \item[Group 3: Supreme Servant] (Chaupais 7--10). The heroic feats of the Ramayana.
    \item[Group 4: Celestial Praise] (Chaupais 11--15). Rama's embrace and the reaction of the Heavens.
    \item[Group 5: The King-Maker] (Chaupais 16--24). Hanuman's role as a protector and diplomat.
    \item[Group 6: The Healer] (Chaupais 25--31). Practical benefits and the banishing of suffering.
    \item[Group 7: Bestower of Grace] (Chaupais 32--40). The granting of Siddhis and Nidhis.
    \item[Group 8: The Final Prayer] (Concluding Doha). Inviting the Divine into the heart.
\end{description}

\subsection*{Final Tips for Success}
\begin{itemize}
    \item \textbf{Read Before Memorizing:} Understand the English meaning first; it provides a mental "hook" for the sounds.
    \item \textbf{Listen and Shadow:} Use the "Shadowing" technique, i.e., listen to a recording and try to speak exactly in sync with the narrator.
    \item \textbf{Simple Sanskrit:} Use \deva{सरल-संस्कृतम्} as a bridge. If the Awadhi is confusing, the Sanskrit root often clarifies the meaning instantly.
\end{itemize}

\newpage
% --- STRETCH MARGINS FOR THIS SECTION ---
\newgeometry{left=1.5cm, right=1.5cm, top=2cm, bottom=2cm}

% ==========================================
\section*{Complete Hanuman Chalisa}
\addcontentsline{toc}{section}{Complete Hanuman Chalisa}
% ==========================================

This section provides the complete and uninterrupted Hanuman Chalisa for recital in Devanagari, Telugu, and IAST Romanization.

\subsection*{Devanagari (\deva{देवनागरी})}
\addcontentsline{toc}{subsection}{\deva{देवनागरी} — \deva{श्रीगुरु चरन...}}

{ % Start formatting group
\large \linespread{1.5}\selectfont
\setcounter{chaupainum}{0} % Reset to 0 before starting

% --- Opening Dohas (Centered) ---
\begin{center}
    \begin{minipage}{0.8\linewidth}
        \centering
        \makebox[\linewidth][l]{\LARGE \deva{श्रीगुरु चरन सरोज रज}}\\
        \makebox[\linewidth][c]{\LARGE \deva{निज मनु मुकुरु सुधारि।}}\\
        \makebox[\linewidth][l]{\LARGE \deva{बरनउँ रघुबर बिमल जसु}}\\
        \makebox[\linewidth][c]{\LARGE \deva{जो दायकु फल चारि॥}}

    \end{minipage} \\[1em]
    
    \begin{minipage}{0.8\linewidth}
        \centering
        \makebox[\linewidth][l]{\LARGE \deva{बुद्धिहीन तनु जानिके }}\\
        \makebox[\linewidth][c]{\LARGE \deva{सुमिरौं पवन कुमार।}} \\ 
        \makebox[\linewidth][l]{\LARGE \deva{बल बुधि बिद्या देहु मोहिं }}\\
        \makebox[\linewidth][c]{\LARGE \deva{ हरहु कलेस बिकार॥}}
    \end{minipage}
\end{center}

\vspace{1em}

% --- Two Column Listing ---
\setlength{\columnsep}{40pt} % Increase space between the two main columns
\begin{multicols}{2}
\raggedcolumns 

% Define the command INSIDE the group to ensure it uses the correct font
\newcommand{\cp}[2]{%
    \refstepcounter{chaupainum}% Increment counter
    \begin{minipage}{\linewidth}
        \begin{tabular}{@{}p{3ex} p{0.88\linewidth}@{}}
            \small \textcolor{gray}{\arabic{chaupainum}.} & 
            \LARGE \deva{#1} \\
            & \LARGE \deva{#2}
        \end{tabular}
    \end{minipage}\\[0.8em]
}

\begin{flushleft}
\cp{जय हनुमान ज्ञान गुन सागर।}{जय कपीस तिहुँ लोक उजागर॥}
\cp{राम दूत अतुलित बल धामा।}{अंजनि पुत्र पवन सुत नामा॥}
\cp{महाबीर बिक्रम बजरंगी।}{कुमति निवार सुमति के संगी॥}
\cp{कंचन बरन बिराज सुबेसा।}{कानन कुंडल कुंचित केसा॥}
\cp{हाथ बज्र औ ध्वजा बिराजै।}{काँधे मूँज जनेऊ साजै॥}
\cp{संकर सुवन केसरी नंदन।}{तेज प्रताप महा जग बंदन॥}
\cp{बिद्यावान गुनी अति चातुर।}{राम काज करिबे को आतुर॥}
\cp{प्रभु चरित्र सुनिबे को रसया।}{राम लखन सीता मन बसिया॥}
\cp{सूक्ष्म रूप धरि सियहिं दिखावा।}{बिकट रूप धरि लंक जरावा॥}
\cp{भीम रूप धरि असुर संहारे।}{रामचंद्र के काज संवारे॥}
\cp{लाय सजीवन लखन जियाये।}{श्रीरघुबीर हरषि उर लाये॥}
\cp{रघुपति कीन्ही बहुत बड़ाई।}{तुम मम प्रिय भरतहि सम भाई॥}
\cp{सहस बदन तुम्हरो जस गावैं।}{अस कहि श्रीपति कंठ लगावैं॥}
\cp{सनकादिक ब्रह्मादि मुनीसा।}{नारद सारद सहित अहीसा॥}
\cp{जमु कुबेर दिगपाल जहाँ ते।}{कबि कोबिद कहि सके कहाँ ते॥}
\cp{तुम उपकार सुग्रीवहिं कीन्हा।}{राम मिलाय राज पद दीन्हा॥}
\cp{तुम्हरो मंत्र बिभीषन माना।}{लंकेस्वर भए सब जग जाना॥}
\cp{जुग सहस्र जोजन पर भानू।}{लील्यो ताहि मधुर फल जानू॥}
\cp{प्रभु मुद्रिका मेलि मुख माहीं।}{जलधि लांघि गये अचरज नाहीं॥}
\cp{दुर्गम काज जगत के जेते।}{सुगम अनुग्रह तुम्हरे तेते॥}
\cp{राम दुआरे तुम रखवारे।}{होत न आज्ञा बिनु पैसारे॥}
\cp{सब सुख लहै तुम्हारी सरना।}{तुम रच्छक काहू को डर ना॥}
\cp{आपन तेज सम्हारो आपै।}{तीनों लोक हांक तें कांपै॥}
\cp{भूत पिसाच निकट नहिं आवै।}{महाबीर जब नाम सुनावै॥}
\cp{नासै रोग हरै सब पीरा।}{जपत निरन्तर हनुमत बीरा॥}
\cp{संकट तें हनुमान छुड़ावै।}{मन क्रम बचन ध्यान जो लावै॥}
\cp{सब पर राम तपस्वी राजा।}{तिन के काज सकल तुम साजा॥}
\cp{और मनोरथ जो कोई लावै।}{सोइ अमित जीवन फल पावै॥}
\cp{चारों जुग परताप तुम्हारा।}{है परसिद्ध जगत उजियारा॥}
\cp{साधु सन्त के तुम रखवारे।}{असुर निकन्दन राम दुलारे॥}
\cp{अष्ट सिद्धि नौ निधि के दाता।}{अस बर दीन जानकी माता॥}
\cp{राम रसायन तुम्हरे पासा।}{सदा रहो रघुपति के दासा॥}
\cp{तुम्हरे भजन राम को पावै।}{जनम जनम के दुख बिसरावै॥}
\cp{अन्त काल रघुबर पुर जाई।}{जहाँ जन्म हरिभक्त कहाई॥}
\cp{और देवता चित्त न धरई।}{हनुमंत सेइ सर्ब सुख करई॥}
\cp{संकट कटै मिटै सब पीरा।}{जो सुमिरै हनुमत बलबीरा॥}
\cp{जै जै जै हनुमान गोसाईं।}{कृपा करहु गुरु देव की नाईं॥}
\cp{जो सत बार पाठ कर कोई।}{छूटहि बंदि महा सुख होई॥}
\cp{जो यह पढ़ै हनुमान चालीसा।}{होय सिद्धि साखी गौरीसा॥}
\cp{तुलसीदास सदा हरि चेरा।}{कीजै नाथ हृदय महँ डेरा॥}
\end{flushleft}
\end{multicols}

\vspace{1em}

% --- Concluding Doha ---
\begin{center}
    \begin{minipage}{0.8\linewidth}
        \centering
        \makebox[\linewidth][l]{\LARGE \deva{पवन तनय संकट हरन}}\\
        \makebox[\linewidth][c]{\LARGE \deva{मंगल मूरति रूप।}} \\ 
        \makebox[\linewidth][l]{\LARGE \deva{राम लखन सीता सहित}}\\
        \makebox[\linewidth][c]{\LARGE \deva{हृदय बसहु सुर भूप॥}}
    \end{minipage}
\end{center}

} % End formatting group

% --- RESTORE ORIGINAL MARGINS ---

\newpage

\subsection*{Telugu (\telu{తెలుగు})}
\addcontentsline{toc}{subsection}{\telu{తెలుగు}}

{ % Start formatting group
\large \linespread{1.5}\selectfont
\setcounter{chaupainum}{0} % Reset to 0 before starting

% --- Opening Dohas (Centered) ---
\begin{center}
    \begin{minipage}{0.8\linewidth}
        \centering
        \makebox[\linewidth][l]{\LARGE \telu{శ్రీగురు చరన సరోజ రజ}}\
        \makebox[\linewidth][c]{\LARGE \telu{నిజ మను ముకురు సుధారి।}} \\ 
        \makebox[\linewidth][l]{\LARGE \telu{బరనఉఁ రఘుబర బిమల జసు}}\
        \makebox[\linewidth][c]{\LARGE \telu{జో దాయకు ఫల చారి॥}}
    \end{minipage} \\[1em]
    
    \begin{minipage}{0.8\linewidth}
        \centering
        \makebox[\linewidth][l]{\LARGE \telu{బుద్ధిహీన తను జానికే}}\ 
        \makebox[\linewidth][c]{\LARGE \telu{సుమిరౌ పవన కుమార।}} \\ 
        \makebox[\linewidth][l]{\LARGE \telu{బల బుధి బిద్యా దేహు మోహిఁ}}\ 
        \makebox[\linewidth][c]{\LARGE \telu{హరహు కలేస బికార॥}}
    \end{minipage}
\end{center}

\vspace{1em}

% --- Two Column Listing ---
\setlength{\columnsep}{40pt} % Increase space between the two main columns
\begin{multicols}{2}
\raggedcolumns 

% Define the command INSIDE the group to ensure it uses the correct font
\newcommand{\cptelu}[2]{%
    \refstepcounter{chaupainum}% Increment counter
    \begin{minipage}{\linewidth}
        \begin{tabular}{@{}p{3ex} p{0.88\linewidth}@{}}
            \small \textcolor{gray}{\arabic{chaupainum}.} & 
            \large \telu{#1} \\
            & \large \telu{#2}
        \end{tabular}
    \end{minipage}\\[0.8em]
}

\begin{flushleft}
\cptelu{జయ హనుమాన జ్ఞాన గున సాగర।}{జయ కపీస తిహుఁ లోక ఉజాగర॥}
\cptelu{రామ దూత అతులిత బల ధామా।}{అంజని పుత్ర పవన సుత నామా॥}
\cptelu{మహాబీర బిక్రమ బజరంగీ।}{కుమతి నివార సుమతి కే సంగీ॥}
\cptelu{కంచన బరన బిరాజ సుబేసా।}{కానన కుండల కుంచిత కేసా॥}
\cptelu{హాథ బజ్ర ఔ ధ్వజా బిరాజై।}{కాఁధే మూఁజ జనేఊ సాజై॥}
\cptelu{సంకర సువన కేసరీ నందన।}{తేజ ప్రతాప మహా జగ బందన॥}
\cptelu{బిద్యావాన గునీ అతి చాతుర।}{రామ కాజ కరిబే కో ఆతుర॥}
\cptelu{ప్రభు చరిత్ర సునిబే కో రసియా।}{రామ లఖన సీతా మన బసియా॥}
\cptelu{సూక్ష్మ రూప ధరి సియహిఁ దిఖావా।}{బికట రూప ధరి లంక జరావా॥}
\cptelu{భీమ రూప ధరి అసుర సంహారే।}{రామచంద్ర కే కాజ సంవారే॥}
\cptelu{లాయ సజీవన లఖన జియాయే।}{శ్రీరఘుబీర హరషి ఉర లాయే॥}
\cptelu{రఘుపతి కీన్హీ బహుత బడాయీ।}{తుమ మమ ప్రియ భరతహి సమ భాయీ॥}
\cptelu{సహస బదన తుమ్హరో జస గావైఁ।}{అస కహి శ్రీపతి కంఠ లగావైఁ॥}
\cptelu{సనకాదిక బ్రహ్మాది మునీసా।}{నారద సారద సహిత అహీసా॥}
\cptelu{జము కుబేర దిగపాల జహాఁ తే।}{కబి కోబిద కహి సకే కహాఁ తే॥}
\cptelu{తుమ ఉపకార సుగ్రీవహిఁ కీన్హా।}{రామ మిలాయ రాజ పద దీన్హా॥}
\cptelu{తుమ్హరో మంత్ర బిభీషన మానా।}{లంకేశ్వర భఏ సబ జగ జానా॥}
\cptelu{జుగ సహస్ర జోజన పర భానూ।}{లీల్యో తాహి మధుర ఫల జానూ॥}
\cptelu{ప్రభు ముద్రికా మేలి ముఖ మాహీఁ।}{జలధి లాంఘి గయే అచరజ నాహీఁ॥}
\cptelu{దుర్గమ కాజ జగత కే జేతే।}{సుగమ అనుగ్రహ తుమ్హరే తేతే॥}
\cptelu{రామ దుఆరే తుమ రఖవారే।}{హోత న ఆజ్ఞా బిను పైసారే॥}
\cptelu{సబ సుఖ లహై తుమ్హారీ సరనా।}{తుమ రచ్ఛక కాహూ కో డర నా॥}
\cptelu{ఆపన తేజ సమ్హారో ఆపై।}{తీనోం లోక హాంక తేం కాంపై॥}
\cptelu{భూత పిసాచ నికట నహిఁ ఆవై।}{మహాబీర జబ నామ సునావై॥}
\cptelu{నాసై రోగ హరై సబ పీరా।}{జపత నిరంతర హనుమత బీరా॥}
\cptelu{సంకట తేం హనుమాన ఛుడావై।}{మన క్రమ బచన ధ్యాన జో లావై॥}
\cptelu{సబ పర రామ తపస్వీ రాజా।}{తిన కే కాజ సకల తుమ సాజా॥}
\cptelu{ఔర మనోరథ జో కోఈ లావై।}{సోఇ అమిత జీవన ఫల పావై॥}
\cptelu{చారోం జుగ పరతాప తుమ్హారా।}{హై పరసిద్ధ జగత ఉజియారా॥}
\cptelu{సాధు సన్త కే తుమ రఖవారే।}{అసుర నికన్దన రామ దులారే॥}
\cptelu{అష్ట సిద్ధి నౌ నిధి కే దాతా।}{అస బర దీన జానకీ మాతా॥}
\cptelu{రామ రసాయన తుమ్హరే పాసా।}{సదా రహో రఘుపతి కే దాసా॥}
\cptelu{తుమ్హరే భజన రామ కో పావై।}{జనమ జనమ కే దుఖ బిసరావై॥}
\cptelu{అన్త కాల రఘుబర పుర జాయీ।}{జహాఁ జన్మ హరిభక్త కహాయీ॥}
\cptelu{ఔర దేవతా చిత్త న ధరఈ।}{హనుమంత సేఇ సర్బ సుఖ కరఈ॥}
\cptelu{సంకట కటై మిటై సబ పీరా।}{జో సుమిరై హనుమత బలబీరా॥}
\cptelu{జై జై జై హనుమాన గోసాయీఁ।}{కృపా కరహు గురు దేవ కీ నాయీఁ॥}
\cptelu{జో సత బార పాఠ కర కోఈ।}{ఛూటహి బంది మహా సుఖ హోఈ॥}
\cptelu{జో యహ పఢై హనుమాన చాలీసా।}{హోయ సిద్ధి సాఖీ గౌరీసా॥}
\cptelu{తులసీదాస సదా హరి చేరా।}{కీజై నాథ హృదయ మహఁ డేరా॥}
\end{flushleft}
\end{multicols}

\vspace{1em}

% --- Concluding Doha ---
\begin{center}
    \begin{minipage}{0.8\linewidth}
        \centering
        \makebox[\linewidth][l]{\LARGE \telu{పవన తనయ సంకట హరన}}
        \makebox[\linewidth][c]{\LARGE \telu{మంగల మూరతి రూప।}} \\ 
        \makebox[\linewidth][l]{\LARGE \telu{రామ లఖన సీతా సహిత}}
        \makebox[\linewidth][c]{\LARGE \telu{హృదయ బసహు సుర భూప॥}}
    \end{minipage}
\end{center}

} % End formatting group
\newpage

\subsection*{IAST Romanization}
\addcontentsline{toc}{subsection}{IAST Romanization}

{ % Start formatting group
\large \linespread{1.5}\selectfont
\setcounter{chaupainum}{0} % Reset to 0 before starting

% --- Opening Dohas (Centered) ---
\begin{center}
    \begin{minipage}{0.8\linewidth}
        \centering
        \large \textit{śrīguru carana saroja raja, nija manu mukuru sudhāri |} \\ 
        \large \textit{baranaũ raghubara bimala jasu, jo dāyaku phala cāri ||}
    \end{minipage} \\[1em]
    
    \begin{minipage}{0.8\linewidth}
        \centering
        \large \textit{buddhihīna tanu jānike, sumiraũ pavana kumāra |} \\ 
        \large \textit{bala budhi bidyā dehu mohĩ, harahu kalesa bikāra ||}
    \end{minipage}
\end{center}

\vspace{1em}

% --- Two Column Listing ---
\setlength{\columnsep}{40pt} % Increase space between the two main columns
\begin{multicols}{2}
\raggedcolumns 

% Define the command INSIDE the group to ensure it uses the correct font
\newcommand{\cpiast}[2]{%
    \refstepcounter{chaupainum}% Increment counter
    \begin{minipage}{\linewidth}
        \begin{tabular}{@{}p{1ex} p{0.88\linewidth}@{}}
            \small \textcolor{gray}{\arabic{chaupainum}.} & 
            \large \textit{#1} \\
            & \large \textit{#2}
        \end{tabular}
    \end{minipage}\\[0.8em]
}

\begin{flushleft}
\cpiast{jaya hanumāna jñāna guna sāgara |}{jaya kapīsa tihũ loka ujāgara ||}
\cpiast{rāma dūta atulita bala dhāmā |}{añjani-putra pavana suta nāmā ||}
\cpiast{mahābīra bikrama bajaraṅgī |}{kumati nivāra sumati ke saṅgī ||}
\cpiast{kañcana barana birāja subesā |}{kānana kuṇḍala kuñcita kesā ||}
\cpiast{hātha bajra au dhvajā birājai |}{kā̃dhe mū̃ja janeū sājai ||}
\cpiast{saṅkara suvana kesarī nandana |}{teja pratāpa mahā jaga bandana ||}
\cpiast{bidyāvāna gunī ati cātura |}{rāma kāja karibe ko ātura ||}
\cpiast{prabhu caritra sunibe ko rasiyā |}{rāma lakhana sītā mana basiyā ||}
\cpiast{sūkṣma rūpa dhari siyahĩ dikhāvā |}{bikaṭa rūpa dhari laṅka jarāvā ||}
\cpiast{bhīma rūpa dhari asura sãhāre |}{rāmacandra ke kāja sãvāre ||}
\cpiast{lāya sajīvana lakhana jiyāye |}{śrīraghubīra haraṣi ura lāye ||}
\cpiast{raghupati kīnhī bahuta baṛāī |}{tuma mama priya bharatahi sama bhāī ||}
\cpiast{sahasa badana tumharo jasa gāvaĩ |}{asa kahi śrīpati kaṇṭha lagāvaĩ ||}
\cpiast{sanakādika brahmādi munīsā |}{nārada sārada sahita ahīsā ||}
\cpiast{jamu kubera digapāla jahā̃ te |}{kabi kobida kahi sake kahā̃ te ||}
\cpiast{tuma upakāra sugrīvahĩ kīnhā |}{rāma milāya rāja pada dīnhā ||}
\cpiast{tumharo mantra bibhīṣana mānā |}{laṅkesvara bhae saba jaga jānā ||}
\cpiast{juga sahasra jojana para bhānū |}{līlyo tāhi madhura phala jānū ||}
\cpiast{prabhu mudrikā meli mukha māhī̃ |}{jaladhi lāṅghi gaye acaraja nāhīṃ ||}
\cpiast{durgama kāja jagata ke jete |}{sugama anugraha tumhare tete ||}
\cpiast{rāma duāre tuma rakhavāre |}{hota na ājñā binu paisāre ||}
\cpiast{saba sukha lahai tumhārī saranā |}{tuma racchaka kāhū ko ḍara nā ||}
\cpiast{āpana teja samhāro āpai |}{tīnõ loka hā̃ka tẽ kā̃pai ||}
\cpiast{bhūta pisāca nikaṭa nahĩ āvai |}{mahābīra jaba nāma sunāvai ||}
\cpiast{nāsai roga harai saba pīrā |}{japata nirantara hanumata bīrā ||}
\cpiast{saṅkaṭa tẽ hanumāna chuṛāvai |}{mana krama bacana dhyāna jo lāvai ||}
\cpiast{saba para rāma tapasvī rājā |}{tina ke kāja sakala tuma sājā ||}
\cpiast{aura manoratha jo koi lāvai |}{soi amita jīvana phala pāvai ||}
\cpiast{cārõ juga paratāpa tumhārā |}{hai parasiddha jagata ujiyārā ||}
\cpiast{sādhu santa ke tuma rakhavāre |}{asura nikandana rāma dulāre ||}
\cpiast{aṣṭa siddhi nau nidhi ke dātā |}{asa bara dīna jānakī mātā ||}
\cpiast{rāma rasāyana tumhare pāsā |}{sadā raho raghupati ke dāsā ||}
\cpiast{tumhare bhajana rāma ko pāvai |}{janama janama ke dukha bisarāvai ||}
\cpiast{anta kāla raghubara pura jāī |}{jahā̃ janma hari-bhakta kahāī ||}
\cpiast{aura devatā citta na dharaī |}{hanumaṃta sei sarba sukha karaī ||}
\cpiast{saṅkaṭa kaṭai miṭai saba pīrā |}{jo sumirai hanumata balabīrā ||}
\cpiast{jai jai jai hanumāna gosāī̃ |}{kṛpā karahu guru deva kī nāī̃ ||}
\cpiast{jo sata bāra pāṭha kara koī |}{chūṭahi bandi mahā sukha hoī ||}
\cpiast{jo yaha paṛhai hanumāna calīsā |}{hoya siddhi sākhī gaurīsā ||}
\cpiast{tulasīdāsa sadā hari cerā |}{kījai nātha hṛdaya mahã ḍerā ||}
\end{flushleft}
\end{multicols}

\vspace{1em}

% --- Concluding Doha ---
\begin{center}
    \begin{minipage}{0.8\linewidth}
        \centering
        \large \textit{pavana tanaya saṅkaṭa harana, maṅgala mūrati rūpa |} \\ 
        \large \textit{rāma lakhana sītā sahita, hṛdaya basahu sura bhūpa ||}
    \end{minipage}
\end{center}

} % End formatting group
\restoregeometry
\newpage

% ==========================================
\section*{Group 1: The Invocation \& Surrender}
\addcontentsline{toc}{section}{Group 1: The Invocation \& Surrender}
% ==========================================

% ----- DOHA 1 -----
\subsection*{\deva{प्रथम दोहा} (First Doha)}
\addcontentsline{toc}{subsection}{\deva{प्रथम दोहा} (First Doha) — \deva{श्रीगुरु चरन...}}

\begin{minipage}[t]{0.55\textwidth}
\vspace{0pt}

{\large \linespread{1.5}\selectfont
\makebox[\linewidth][l]{\deva{श्रीगुरु चरन सरोज रज}}\\
\makebox[\linewidth][c]{\deva{निज मनु मुकुरु सुधारि।}}\\
\makebox[\linewidth][l]{\deva{बरनउँ रघुबर बिमल जसु}}\\
\makebox[\linewidth][c]{\deva{जो दायकु फल चारि॥}}
\par}

\vspace{0.5em}

{ \linespread{1.5}\selectfont
\makebox[\linewidth][l]{\telu{శ్రీగురు చరన సరోజ రజ }}\\
\makebox[\linewidth][c]{\telu{నిజ మను ముకురు సుధారి।}}\\
\makebox[\linewidth][l]{\telu{బరనఉఁ రఘుబర బిమల జసు}}\\
\makebox[\linewidth][c]{\telu{జో దాయకు ఫల చారి॥}}
\par}

\vspace{0.5em}

{ \linespread{1.3}\selectfont
\textit{śrīguru carana saroja raja, nija manu mukuru sudhāri |}\\
\textit{baranaũ raghubara bimala jasu, jo dāyaku phala cāri ||}
\par}
\end{minipage}%
\hfill
\begin{minipage}[t]{0.42\textwidth}
\vspace{0pt}
\begin{tcolorbox}[anchorbox]
\textbf{The Frame of Mind & Dedication:} Why begin with \textit{Guru Vandana} (respecting the teacher)? Because actively thinking about purifying one's mind (\textit{Mukuru Sudhari}) immediately puts the student in a positive, receptive frame of mind to learn. Hanuman perfectly modeled this respect for education: to learn the scriptures from Surya (the Sun), he flew backward all day across the blazing sky just to keep his eyes respectfully on his teacher. It shows supreme dedication to learning.
\vspace{0.5em}\newline Jambavan acts as a true mentor, wiping away self-doubt by reminding us of our immense, hidden potential.\newline\null\hfill\href{https://www.valmiki.iitk.ac.in/sloka?field_kanda_tid=4&language=dv&field_sarga_value=66}{\textit{Kishkindha Kanda, 66:31-35}}
\end{tcolorbox}
\end{minipage}

\vspace{0.5em}
\subsubsection*{\textbf{Visual Rhythm Map:}}
\begin{table}[h!]
\centering
\renewcommand{\arraystretch}{1.2}
\begin{tabularx}{\textwidth}{|>{\centering\arraybackslash\hsize=1.3333\hsize}X|>{\centering\arraybackslash\hsize=1.0000\hsize}X|>{\centering\arraybackslash\hsize=1.3333\hsize}X|>{\centering\arraybackslash\hsize=0.6667\hsize}X|>{\centering\arraybackslash\hsize=0.6667\hsize}X|>{\centering\arraybackslash\hsize=0.6667\hsize}X|>{\centering\arraybackslash\hsize=1.0000\hsize}X|>{\centering\arraybackslash\hsize=1.3333\hsize}X|}
\hline
\textbf{\guru{श्री}\laghu{गु}\laghu{रु}} \par (4) & \textbf{\laghu{च}\laghu{र}\laghu{न}} \par (3) & \textbf{\laghu{स}\guru{रो}\laghu{ज}} \par (4) & \textbf{\laghu{र}\laghu{ज}} \par (2) & \textbf{\laghu{नि}\laghu{ज}} \par (2) & \textbf{\laghu{म}\laghu{नु}} \par (2) & \textbf{\laghu{मु}\laghu{कु}\laghu{रु}} \par (3) & \textbf{\laghu{सु}\guru{धा}\laghu{रि}} \par (4) \\
\hline
\end{tabularx}
\vspace{-1.5em}
\end{table}

\begin{table}[h!]
\centering
\renewcommand{\arraystretch}{1.2}
\begin{tabularx}{\textwidth}{|>{\centering\arraybackslash\hsize=1.3333\hsize}X|>{\centering\arraybackslash\hsize=1.3333\hsize}X|>{\centering\arraybackslash\hsize=1.0000\hsize}X|>{\centering\arraybackslash\hsize=0.6667\hsize}X|>{\centering\arraybackslash\hsize=0.6667\hsize}X|>{\centering\arraybackslash\hsize=1.3333\hsize}X|>{\centering\arraybackslash\hsize=0.6667\hsize}X|>{\centering\arraybackslash\hsize=1.0000\hsize}X|}
\hline
\textbf{\laghu{ब}\laghu{र}\laghu{न}\laghu{उँ}} \par (4) & \textbf{\laghu{र}\laghu{घु}\laghu{ब}\laghu{र}} \par (4) & \textbf{\laghu{बि}\laghu{म}\laghu{ल}} \par (3) & \textbf{\laghu{ज}\laghu{सु}} \par (2) & \textbf{\guru{जो}} \par (2) & \textbf{\guru{दा}\laghu{य}\laghu{कु}} \par (4) & \textbf{\laghu{फ}\laghu{ल}} \par (2) & \textbf{\guru{चा}\laghu{रि}} \par (3) \\
\hline
\end{tabularx}
\vspace{-1.5em}
\end{table}

\vspace{0.5em}
\textbf{ \deva{सरल-संस्कृतम्}:}\\
 \deva{श्रीगुरोः चरण-कमल-धूल्या स्व-मनः-रूपी-दर्पणं स्वच्छं कृत्वा}\\
 \deva{अहं श्रीरघुनाथस्य निर्मलं यशः वर्णयामि, यत् यशः चत्वारि फलानि (धर्म-अर्थ-काम-मोक्षान्) ददाति।}\\

Having cleansed the mirror of my mind with the dust of my Guru's lotus feet\\
I narrate the pure glory of Lord Rama, which bestows the four fruits of life.


\vspace{1em}
\newpage
\textbf{\deva{प्रतिपदार्थः} (Meaning of each word):}
\begin{table}[h!]
\centering
\renewcommand{\arraystretch}{1.2}
\begin{tabularx}{\textwidth}{@{} l l X X @{}}
\toprule
\textbf{Awadhi} & \textbf{Sanskrit} & \textbf{English} & \textbf{Grammar \& Notes} \\
\midrule
\deva{श्रीगुरु} / \telu{శ్రీగురు} / \textit{śrīguru}  &  \deva{श्रीगुरोः}  &  Venerable Guru's  & Genitive (-ḥ) \\
\deva{चरन} / \telu{చరన} / \textit{carana}  &  \deva{चरण}  &  Feet  & \\ 
\deva{सरोज} / \telu{సరోజ} / \textit{saroja}  &  \deva{कमल}  &  Lotus  &   \\
\deva{रज} / \telu{రజ} / \textit{raja}  &  \deva{धूलिः / रजः}  &  Dust  &   \\
\deva{निज} / \telu{నిజ} / \textit{nija}  &  \deva{स्व / आत्मीय}  &  Own  &   \\
\deva{मनु} / \telu{మను} / \textit{manu}  &  \deva{मनः}  &  Mind  &   \\
\deva{मुकुरु} \gotchaicon / \telu{ముకురు} / \textit{mukuru}  &  \deva{दर्पणम्}  &  Mirror  & Nominative -u ending \\
\deva{सुधारि} \gotchaicon / \telu{సుధారి} / \textit{sudhāri}  &  \deva{स्वच्छं कृत्वा (संशोध्य)}  &  Having cleaned/purified  &   \\
\deva{बरनउँ} / \telu{బరనఉఁ} / \textit{baranaũ}  &  \deva{वर्णयामि}  &  I describe / narrate  & \\ 
\deva{रघुबर} / \telu{రఘుబర} / \textit{raghubara}  &  \deva{श्रीरघुनाथस्य}  &  Lord Rama's  & \\ 
\deva{बिमल} / \telu{బిమల} / \textit{bimala}  &  \deva{निर्मलम्}  &  Pure  & \\ 
\deva{जसु} \gotchaicon / \telu{జసు} / \textit{jasu}  &  \deva{यशः}  &  Glory  & \\ 
\deva{जो} / \telu{జో} / \textit{jo}  &  \deva{यत्}  &  Which  &   \\
\deva{दायकु} \gotchaicon / \telu{దాయకు} / \textit{dāyaku}  &  \deva{दायकम् (फलप्रदम्)}  &  Bestower / Giver  & Nominative -u ending \\
\deva{फल चारि} \glossaryicon / \telu{ఫల చారి} / \textit{phala cāri}  &  \deva{चत्वारि फलानि}  &  Four fruits  &   \\
\bottomrule
\end{tabularx}
\end{table}

\begin{tcolorbox}[gotchabox]
\begin{itemize}
    \item \textbf{The -u Ending (Nominative Case):} The short 'u' on \deva{मुकुरु} (Mukuru), \deva{जसु} (Jasu), and \deva{दायकु} (Dayaku) doesn't mean it's a verb---it is the classical Apabhramsha nominative case marker derived from the Sanskrit '-aḥ' ending.
    \item \textbf{Meter Alert:} Keep the 'i' short in \deva{सुधारि} (Sudhaari) and \deva{चारि} (Chaari) to maintain the exact 11-beat math of the Doha's second half! Both end in the classic Guru-Laghu (2-1) drop.
\end{itemize}
\end{tcolorbox}

\begin{tcolorbox}[poetrybox]
Notice the beautiful structural pairing of \deva{सुधारि} (\textit{sudhāri}) and \deva{चारि} (\textit{cāri}) at the end of the half-verses, establishing the melodic pace of the Doha meter.
\end{tcolorbox}


\newpage


% ----- DOHA 2 -----

\newpage

\subsection*{\deva{द्वितीय दोहा} (Second Doha)}
\addcontentsline{toc}{subsection}{\deva{द्वितीय दोहा} (Second Doha) — \deva{बुद्धिहीन तनु...}}

\begin{minipage}[t]{0.55\textwidth}
\vspace{0pt}

{\large \linespread{1.5}\selectfont
\makebox[\linewidth][l]{\deva{बुद्धिहीन तनु जानिके}}\\
\makebox[\linewidth][c]{\deva{सुमिरौं पवन कुमार।}}\\
\makebox[\linewidth][l]{\deva{बल बुधि बिद्या देहु मोहिं }}\\
\makebox[\linewidth][c]{\deva{हरहु कलेस बिकार॥}}
\par}

\vspace{0.5em}

{\large \linespread{1.5}\selectfont
\makebox[\linewidth][l]{\telu{బుద్ధిహీన తను జానికే}}\\
\makebox[\linewidth][c]{\telu{సుమిరౌఁ పవన కుమార।}}\\
\makebox[\linewidth][l]{\telu{బల బుధి బిద్యా దేహు మోహిఁ}}\\
\makebox[\linewidth][c]{\telu{హరహు కలేస బికార॥}}
\par}

\vspace{0.5em}

{\large \linespread{1.3}\selectfont
\textit{buddhihīna tanu jānike, sumiraũ pavana kumāra |}\\
\textit{bala budhi bidyā dehu mohĩ, harahu kalesa bikāra ||}
\par}
\end{minipage}%
\hfill
\begin{minipage}[t]{0.42\textwidth}
\vspace{0pt}
\begin{tcolorbox}[anchorbox]
\textbf{The Qualities of Success:} We ask for strength (\textit{Bala}), intellect (\textit{Budhi}), and wisdom (\textit{Bidya}). In the *Valmiki Ramayana* (*Sundara Kanda, 1:198*), the celestials praise Hanuman, declaring that he who has four key leadership traits—courage (\textit{Dhriti}), vision (\textit{Drishti}), intellect (\textit{Mati}), and skill (\textit{Dakshya})—like him, will never fail in his undertakings.
\vspace{0.5em}\newline Hanuman's four timeless qualities of courage, vision, intellect, and skill represent the ultimate blueprint for overcoming life's challenges.\newline\null\hfill\href{https://www.valmiki.iitk.ac.in/sloka?field_kanda_tid=5&language=dv&field_sarga_value=1}{\textit{Sundara Kanda, 1:198}}
\end{tcolorbox}
\end{minipage}

\vspace{0.5em}
\subsubsection*{\textbf{Visual Rhythm Map:}}
\begin{table}[h!]
\centering
\renewcommand{\arraystretch}{1.2}
\begin{tabularx}{\textwidth}{|>{\centering\arraybackslash\hsize=1.5000\hsize}X|>{\centering\arraybackslash\hsize=0.5000\hsize}X|>{\centering\arraybackslash\hsize=1.2500\hsize}X|>{\centering\arraybackslash\hsize=1.0000\hsize}X|>{\centering\arraybackslash\hsize=0.7500\hsize}X|>{\centering\arraybackslash\hsize=1.0000\hsize}X|}
\hline
\textbf{\guru{बु}\laghu{द्धि}\guru{ही}\laghu{न}} \par (6) & \textbf{\laghu{त}\laghu{नु}} \par (2) & \textbf{\guru{जा}\laghu{नि}\guru{के}} \par (5) & \textbf{\laghu{सु}\laghu{मि}\guru{रौ}} \par (4) & \textbf{\laghu{प}\laghu{व}\laghu{न}} \par (3) & \textbf{\laghu{कु}\guru{मा}\laghu{र}} \par (4) \\
\hline
\end{tabularx}
\vspace{-1.5em}
\end{table}

\begin{table}[h!]
\centering
\renewcommand{\arraystretch}{1.2}
\begin{tabularx}{\textwidth}{|>{\centering\arraybackslash\hsize=0.6667\hsize}X|>{\centering\arraybackslash\hsize=0.6667\hsize}X|>{\centering\arraybackslash\hsize=1.0000\hsize}X|>{\centering\arraybackslash\hsize=1.0000\hsize}X|>{\centering\arraybackslash\hsize=1.0000\hsize}X|>{\centering\arraybackslash\hsize=1.0000\hsize}X|>{\centering\arraybackslash\hsize=1.3333\hsize}X|>{\centering\arraybackslash\hsize=1.3333\hsize}X|}
\hline
\textbf{\laghu{ब}\laghu{ल}} \par (2) & \textbf{\laghu{बु}\laghu{धि}} \par (2) & \textbf{\laghu{बि}\guru{द्या}} \par (3) & \textbf{\guru{दे}\laghu{हु}} \par (3) & \textbf{\guru{मो}\laghu{हिं}} \par (3) & \textbf{\laghu{ह}\laghu{र}\laghu{हु}} \par (3) & \textbf{\laghu{क}\guru{ले}\laghu{स}} \par (4) & \textbf{\laghu{बि}\guru{का}\laghu{र}} \par (4) \\
\hline
\end{tabularx}
\vspace{-1.5em}
\end{table}

\vspace{0.5em}
\textbf{ \deva{सरल-संस्कृतम्}:}\\
 \deva{हे पवनकुमार! मां बुद्धिहीनं ज्ञात्वा (मत्वा) अहं त्वां स्मरामि।}\\
 \deva{मह्यं बलं, बुद्धिं, विद्यां च यच्छ (देहि), मम सर्वान् क्लेशान् विकारान् (दोषान्) च दूरीकुरु।}\\

Knowing myself to be devoid of intelligence, O Son of the Wind, I remember you.\\
Grant me strength, wisdom, and knowledge, and remove all my afflictions and impurities.


\vspace{1em}
\newpage
\textbf{\deva{प्रतिपदार्थः} (Meaning of each word):}
\begin{table}[h!]
\centering
\renewcommand{\arraystretch}{1.2}
\begin{tabularx}{\textwidth}{@{} l l X X @{}}
\toprule
\textbf{Awadhi} & \textbf{Sanskrit} & \textbf{English} & \textbf{Grammar \& Notes} \\
\midrule
\deva{बुद्धिहीन} / \telu{బుద్ధిహీన} / \textit{buddhihīna}  &  \deva{बुद्धि-रहितं}  &  Devoid of intelligence  &   \\
\deva{तनु} / \telu{తను} / \textit{tanu}  &  \deva{शरीरं / आत्मानम्}  &  Body / Myself  &   \\
\deva{जानिके} \gotchaicon / \telu{జానికే} / \textit{jānike}  &  \deva{ज्ञात्वा}  &  Having known  & Conjunctive Participle (-ike) \\
\deva{सुमिरौं} / \telu{సుమిరౌఁ} / \textit{sumiraũ}  &  \deva{स्मरामि}  &  I remember/meditate  &   \\
\deva{पवन-कुमार} / \telu{పవన-కుమార} / \textit{pavana kumāra}  &  \deva{हे वायुपुत्र}  &  O Son of the Wind  &   \\
\deva{बल} / \telu{బల} / \textit{bala}  &  \deva{बलम्}  &  Strength  &   \\
\deva{बुधि} \gotchaicon / \telu{బుధి} / \textit{budhi}  &  \deva{बुद्धिम्}  &  Wisdom  & Meter tetris (1+1=2) \\
\deva{बिद्या} / \telu{బిద్యా} / \textit{bidyā}  &  \deva{विद्याम्}  &  Knowledge  & \\ 
\deva{देहु} \gotchaicon / \telu{దేహు} / \textit{dehu}  &  \deva{देहि (यच्छ)}  &  Give / Grant  & Imperative -hu ending \\
\deva{मोहिं} \gotchaicon / \telu{మోహిఁ} / \textit{mohĩ}  &  \deva{मह्यम्}  &  To me  & Dative marker (-hiṃ) \\
\deva{हरहु} \gotchaicon / \telu{హరహు} / \textit{harahu}  &  \deva{हर (दूरीकुरु)}  &  Remove / Destroy  & Imperative -hu ending \\
\deva{कलेस} / \telu{కలేస} / \textit{kalesa}  &  \deva{क्लेशान् / दुःखानि}  &  Sufferings  & \\ 
\deva{बिकार} / \telu{బికార} / \textit{bikāra}  &  \deva{विकारान् / दोषान्}  &  Impurities / Flaws  & \\ 
\bottomrule
\end{tabularx}
\end{table}

\begin{tcolorbox}[gotchabox]
\begin{itemize}
    \item \textbf{Verbs ending in -hu (\deva{देहु, हरहु}):} Adding -hu (-ahu) to a root is the Awadhi equivalent of the Sanskrit imperative mood (making a request/command).
    \item \textbf{Dative Suffix -hiṃ (\deva{मोहिं}):} Maps to Sanskrit Mahyam ("to me"). "Dehu mohĩ" = Give to me.
    \item \textbf{Anusvara vs Chandrabindu (\deva{मोहिं} / \telu{మోహిఁ}):} Why does Devanagari use a dot (\deva{ं}) while Telugu uses a half-circle (\telu{ఁ}) for the exact same nasal sound? In Devanagari orthography, when a vowel matra extends \textit{above} the top line (like \deva{ि, ी, े, ो}), there is physically no vertical space to draw the full Chandrabindu (\deva{ँ}). So, purely for typographical convenience, writers use the Anusvara dot (\deva{ं}) even though it is pronounced as a nasalized vowel. However, if the vowel matra goes \textit{below} or beside the line (like the 'u' in \deva{तिहुँ}), the space above is empty, so the full Chandrabindu \deva{ँ} is correctly written! Telugu script has no top-line constraint, so it always correctly uses the Ara-sunna (\telu{ఁ}).
    \item \textbf{Meter Tetris (\deva{बुद्धि} vs \deva{बुधि}):} Line 1 heavily weights \deva{बुद्धि} (bu + conjunct ddhi) to make it 2+1+2+1=6 beats. Line 2 uses light Awadhi \deva{बुधि} (1+1=2 beats) to save space!
    \item \textbf{The 'e' vowel weight (\deva{जानिके}):} The sound 'e' is historically a compound vowel (a + i = e) and physically requires twice as much breath to produce, making it an automatic 2-beat Guru.
\end{itemize}
\end{tcolorbox}

\newpage



% ==========================================
\section*{Group 2: Appearance \& Attributes}
\addcontentsline{toc}{section}{Group 2: Appearance \& Attributes}
% ==========================================

% ----- CHAUPAI 1 -----
\subsection*{\deva{प्रथम चौपाई} (First Chaupai)}
\addcontentsline{toc}{subsection}{\deva{प्रथम चौपाई} (First Chaupai) — \deva{जय हनुमान...}}

\begin{minipage}[t]{0.55\textwidth}
\vspace{0pt}

{\large \linespread{1.5}\selectfont
\deva{जय हनुमान ज्ञान गुन सागर।}\\
\deva{जय कपीस तिहुँ लोक उजागर॥}
\par}

\vspace{0.5em}

{\large \linespread{1.5}\selectfont
\telu{జయ హనుమాన జ్ఞాన గున సాగర।}\\
\telu{జయ కపీస తిహుఁ లోక ఉజాగర॥}
\par}

\vspace{0.5em}

{\large \linespread{1.3}\selectfont
\textit{jaya hanumāna jñāna guna sāgara |}\\
\textit{jaya kapīsa tihũ loka ujāgara ||}
\par}
\end{minipage}%
\hfill
\begin{minipage}[t]{0.42\textwidth}
\vspace{0pt}
\begin{tcolorbox}[anchorbox]
\textbf{The Power of Good Communication:} Hanuman is praised here as an "Ocean of Knowledge." When Lord Rama first met him in the Kishkindha forest, Rama was immediately blown away! Not by his physical strength, but by his flawless speech, clear grammar, and humble posture. True intelligence (\textit{Jnana}) is best expressed through respectful, perfectly articulated communication.
\vspace{0.5em}\newline Lord Rama instantly respects Hanuman's flawless communication, proving that meticulous education commands the highest respect.\newline\null\hfill\href{https://www.valmiki.iitk.ac.in/sloka?field_kanda_tid=4&language=dv&field_sarga_value=3}{\textit{Kishkindha Kanda, 3:28-33}}
\end{tcolorbox}
\end{minipage}

\vspace{0.5em}
\subsubsection*{\textbf{Visual Rhythm Map:}}
\begin{table}[h!]
\centering
\renewcommand{\arraystretch}{1.2}
\setlength{\tabcolsep}{0pt}
\begin{tabularx}{\textwidth}{|*{16}{>{\centering\arraybackslash\hsize=1.0\hsize}X|}}
\hline
\multicolumn{2}{|>{\centering\arraybackslash\hsize=2.0\hsize}X|}{\textbf{\laghu{ज}\laghu{य}} \par (2)} &
\multicolumn{5}{>{\centering\arraybackslash\hsize=5.0\hsize}X|}{\textbf{\laghu{ह}\laghu{नु}\guru{मा}\laghu{न}} \par (5)} &
\multicolumn{3}{>{\centering\arraybackslash\hsize=3.0\hsize}X|}{\textbf{\guru{ज्ञा}\laghu{न}} \par (3)} &
\multicolumn{2}{>{\centering\arraybackslash\hsize=2.0\hsize}X|}{\textbf{\laghu{गु}\laghu{न}} \par (2)} &
\multicolumn{4}{>{\centering\arraybackslash\hsize=4.0\hsize}X|}{\textbf{\guru{सा}\laghu{ग}\laghu{र}} \par (4)} \\
\hline
\end{tabularx}
\vspace{-1.5em}
\end{table}
\begin{table}[h!]
\centering
\renewcommand{\arraystretch}{1.2}
\setlength{\tabcolsep}{0pt}
\begin{tabularx}{\textwidth}{|*{16}{>{\centering\arraybackslash\hsize=1.0\hsize}X|}}
\hline
\multicolumn{2}{|>{\centering\arraybackslash\hsize=2.0\hsize}X|}{\textbf{\laghu{ज}\laghu{य}} \par (2)} &
\multicolumn{4}{>{\centering\arraybackslash\hsize=4.0\hsize}X|}{\textbf{\laghu{क}\guru{पी}\laghu{स}} \par (4)} &
\multicolumn{2}{>{\centering\arraybackslash\hsize=2.0\hsize}X|}{\textbf{\laghu{ति}\laghu{हुँ}} \par (2)} &
\multicolumn{3}{>{\centering\arraybackslash\hsize=3.0\hsize}X|}{\textbf{\guru{लो}\laghu{क}} \par (3)} &
\multicolumn{5}{>{\centering\arraybackslash\hsize=5.0\hsize}X|}{\textbf{\laghu{उ}\guru{जा}\laghu{ग}\laghu{र}} \par (5)} \\
\hline
\end{tabularx}
\vspace{-1.5em}
\end{table}

\vspace{1em}

 \textbf{ \deva{सरल-संस्कृतम्}:}\\
 \deva{हे हनुमान्! भवतः जयः भवतु, यः ज्ञानस्य गुणानां च सागरः अस्ति।}\\
 \deva{हे कपीश (वानराणां स्वामिन्)! भवतः जयः भवतु, यः त्रिषु लोकेषु (स्वर्गे, भूमौ, पाताले) प्रकाशमानः (प्रसिद्धः) अस्ति।}\\

 Victory to you, O Hanuman, who is the ocean of knowledge and virtues.\\
 Victory to the Lord of the Monkeys, who illuminates the three worlds.

 \vspace{1em}
\newpage
\textbf{\deva{प्रतिपदार्थः} (Meaning of each word):}
\begin{table}[h!]
\centering
\renewcommand{\arraystretch}{1.2}
\begin{tabularx}{\textwidth}{@{} l l X X @{}}
\toprule
\textbf{Awadhi} & \textbf{Sanskrit} & \textbf{English} & \textbf{Grammar \& Notes} \\
\midrule
\deva{ज्ञान} / \telu{జ్ఞాన}  &  \deva{ज्ञानस्य}  &  Of knowledge  &   \\
\deva{गुन} / \telu{గున}  &  \deva{गुणानाम्}  &  Of virtues/qualities  & \\ 
\deva{सागर} / \telu{సాగర}  &  \deva{सागरः}  &  Ocean  &   \\
\deva{कपीस} \gotchaicon / \telu{కపీస}  &  \deva{कपीशः (कपि+ईश)}  &  Lord of Monkeys  & \\ 
\deva{तिहुँ} / \telu{తిహుఁ}  &  \deva{त्रिषु}  &  In the three  & $tri \rightarrow tihu$ \\ 
\deva{लोक} / \telu{లోక}  &  \deva{लोकेषु}  &  Worlds  &   \\
\deva{उजागर} / \telu{ఉజాగర}  &  \deva{उजागरः (प्रसिद्धः)}  &  Illuminator / Brightness  &   \\
\bottomrule
\end{tabularx}
\end{table}

\begin{tcolorbox}[gotchabox]
\begin{itemize}
    \item \textbf{The \deva{श/ष} $\rightarrow$ \deva{स} Shift:} \deva{कपीश} (Kapeesha) becomes \deva{कपीस} (Kapeesa).
    \item \textbf{The \deva{ण} $\rightarrow$ \deva{न} Shift:} Sanskrit \deva{गुण} with a retroflex N becomes  \deva{गुन} with a dental 'n'. 
    \item \textbf{Schwa Holding:} To maintain the 16-matra beat, do not cut the final consonants dead. Hold the "r" in Saaga-r and Ujaaga-r for exactly one beat.
\end{itemize}
\end{tcolorbox}

\newpage


% ----- CHAUPAI 2 -----

\newpage

\subsection*{\deva{द्वितीय चौपाई} (Second Chaupai)}
\addcontentsline{toc}{subsection}{\deva{द्वितीय चौपाई} (Second Chaupai) — \deva{राम दूत...}}

\begin{minipage}[t]{0.55\textwidth}
\vspace{0pt}

{\large \linespread{1.5}\selectfont
\deva{राम दूत अतुलित बल धामा।}\\
\deva{अंजनि-पुत्र पवन सुत नामा॥}
\par}

\vspace{0.5em}

{\large \linespread{1.5}\selectfont
\telu{రామ దూత అతులిత బల ధామా।}\\
\telu{అంజని పుత్ర పవన సుత నామా॥}
\par}

\vspace{0.5em}

{\large \linespread{1.3}\selectfont
\textit{rāma dūta atulita bala dhāmā |}\\
\textit{añjani-putra pavana suta nāmā ||}
\par}
\end{minipage}%
\hfill
\begin{minipage}[t]{0.42\textwidth}
\vspace{0pt}
\begin{tcolorbox}[anchorbox]
\textbf{Functional Leadership:} True leadership is about capability and trust, not titles. When millions of soldiers were dispatched to find Sita, Prince Angada was the official commander of the party, yet Lord Rama handed his personal signet ring to Hanuman alone. Rama recognized Hanuman as the functional leader whose competence, intellect, and empathy would guide the team to success.
\vspace{0.5em}\newline Out of millions, Lord Rama entrusted his signet ring uniquely to Hanuman because of his unmatched leadership and reliability.\newline\null\hfill\href{https://www.valmiki.iitk.ac.in/sloka?field_kanda_tid=4&language=dv&field_sarga_value=44}{\textit{Kishkindha Kanda, 44:11-13}}
\end{tcolorbox}
\end{minipage}

\vspace{0.5em}
\subsubsection*{\textbf{Visual Rhythm Map:}}
\begin{table}[h!]
\centering
\renewcommand{\arraystretch}{1.2}
\setlength{\tabcolsep}{0pt}
\begin{tabularx}{\textwidth}{|*{16}{>{\centering\arraybackslash\hsize=1.0\hsize}X|}}
\hline
\multicolumn{3}{|>{\centering\arraybackslash\hsize=3.0\hsize}X|}{\textbf{\guru{रा}\laghu{म}} \par (3)} &
\multicolumn{3}{>{\centering\arraybackslash\hsize=3.0\hsize}X|}{\textbf{\guru{दू}\laghu{त}} \par (3)} &
\multicolumn{4}{>{\centering\arraybackslash\hsize=4.0\hsize}X|}{\textbf{\laghu{अ}\laghu{तु}\laghu{लि}\laghu{त}} \par (4)} &
\multicolumn{2}{>{\centering\arraybackslash\hsize=2.0\hsize}X|}{\textbf{\laghu{ब}\laghu{ल}} \par (2)} &
\multicolumn{4}{>{\centering\arraybackslash\hsize=4.0\hsize}X|}{\textbf{\guru{धा}\guru{मा}} \par (4)} \\
\hline
\end{tabularx}
\vspace{-1.5em}
\end{table}

\begin{table}[h!]
\centering
\renewcommand{\arraystretch}{1.2}
\setlength{\tabcolsep}{0pt}
\begin{tabularx}{\textwidth}{|*{16}{>{\centering\arraybackslash\hsize=1.0\hsize}X|}}
\hline
\multicolumn{4}{|>{\centering\arraybackslash\hsize=4.0\hsize}X|}{\textbf{\guru{अं}\laghu{ज}\laghu{नि}} \par (4)} &
\multicolumn{3}{>{\centering\arraybackslash\hsize=3.0\hsize}X|}{\textbf{\guru{पु}\laghu{त्र}} \par (3)} &
\multicolumn{3}{>{\centering\arraybackslash\hsize=3.0\hsize}X|}{\textbf{\laghu{प}\laghu{व}\laghu{न}} \par (3)} &
\multicolumn{2}{>{\centering\arraybackslash\hsize=2.0\hsize}X|}{\textbf{\laghu{सु}\laghu{त}} \par (2)} &
\multicolumn{4}{>{\centering\arraybackslash\hsize=4.0\hsize}X|}{\textbf{\guru{ना}\guru{मा}} \par (4)} \\
\hline
\end{tabularx}
\vspace{-1.5em}
\end{table}

\vspace{1em}

    \textbf{ \deva{सरल-संस्कृतम्}:}\\
 \deva{हे रामदूत! त्वम् अतुलित-बलस्य धाम (आश्रयः) असि।}\\
 \deva{हे अञ्जनिपुत्र! त्वं पवनसुत-नाम्ना प्रसिद्धः असि॥}\\

O Emissary of Rama! You are the abode of incomparable strength.\\
O Son of Anjani! You are famous by the name 'Son of the Wind'.


\vspace{1em}
\newpage
\textbf{\deva{प्रतिपदार्थः} (Meaning of each word):}
\begin{table}[h!]
\centering
\renewcommand{\arraystretch}{1.2}
\begin{tabularx}{\textwidth}{@{} l l X X @{}}
\toprule
\textbf{Awadhi} & \textbf{Sanskrit} & \textbf{English} & \textbf{Grammar \& Notes} \\
\midrule
\deva{राम दूत} / \telu{రామ దూత} / \textit{rāma dūta}  &  \deva{श्रीरामस्य दूतः}  &  Emissary of Rama  &   \\
\deva{अतुलित} / \telu{అతులిత} / \textit{atulita}  &  \deva{अतुलितम्}  &  Incomparable  &   \\
\deva{बल} / \telu{బల} / \textit{bala}  &  \deva{बलम्}  &  Strength  &   \\
\deva{धामा} / \telu{ధామా} / \textit{dhāmā}  &  \deva{धाम / आश्रयः}  &  Abode / Repository  & Meter: Added 'aa' for rhyme \\
\deva{अंजनि पुत्र} / \telu{అంజని పుత్ర} / \textit{añjani putra}  &  \deva{अञ्जनायाः पुत्रः}  &  Son of Anjani  &   \\
\deva{पवन सुत} / \telu{పవన సుత} / \textit{pavana suta}  &  \deva{पवनस्य सुतः}  &  Son of the Wind  &   \\
\deva{नामा} / \telu{నామా} / \textit{nāmā}  &  \deva{नाम्ना}  &  By the name  & Meter: Added 'aa' for rhyme \\
\bottomrule
\end{tabularx}
\end{table}

\begin{tcolorbox}[poetrybox]
\textbf{Praasa (The Rhyme Engine):}\\
Note the elongated \deva{आ} (aa) at the end of \deva{धामा} (\textit{dhāmā}) and \deva{नामा} (\textit{nāmā}). Sanskrit/Hindi would use \deva{धाम} (Dhaam) and \deva{नाम} (Naam), but Awadhi uses the heavy 2-beat suffix to complete the 16-beat line and create the musical bounce.
\end{tcolorbox}

\begin{tcolorbox}[gotchabox]
\begin{itemize}
    \item \textbf{Conjunct Consonant Weight (\deva{पुत्र}):} The short 'u' in \deva{पु} gets weighed down and becomes a 2-beat Guru because it crashes into the conjunct consonant \deva{त्र} (t+ra).
\end{itemize}
\end{tcolorbox}



\newpage


% ----- CHAUPAI 3 -----

\newpage

\subsection*{\deva{तृतीय चौपाई} (Third Chaupai)}
\addcontentsline{toc}{subsection}{\deva{तृतीय चौपाई} (Third Chaupai) — \deva{महाबीर बिक्रम...}}

\begin{minipage}[t]{0.55\textwidth}
\vspace{0pt}

{\large \linespread{1.5}\selectfont
\deva{महाबीर बिक्रम बजरंगी।}\\
\deva{कुमति निवार सुमति के संगी॥}
\par}

\vspace{0.5em}

{\large \linespread{1.5}\selectfont
\telu{మహాబీర బిక్రమ బజరంగీ।}\\
\telu{కుమతి నివార సుమతి కే సంగీ॥}
\par}

\vspace{0.5em}

{\large \linespread{1.3}\selectfont
\textit{mahābīra bikrama bajaraṅgī |}\\
\textit{kumati nivāra sumati ke saṅgī ||}
\par}
\end{minipage}%
\hfill
\begin{minipage}[t]{0.42\textwidth}
\vspace{0pt}
\begin{tcolorbox}[anchorbox]
\textbf{Choosing Good Company:} Physical strength (a \textit{Vajra} body) is dangerous without a clear mind. Hanuman's greatest power is that he dispels toxic, negative thoughts (\textit{Kumati}) and actively brings us into the company of good thoughts and wise people (\textit{Sumati ke Sangi}). True strength is maintaining a disciplined mind.
\vspace{0.5em}\newline When everyone else panicked, Hanuman delivered brilliant, clear-headed counsel to accept Vibhishana.\newline\null\hfill\href{https://www.valmiki.iitk.ac.in/sloka?field_kanda_tid=6&language=dv&field_sarga_value=17}{\textit{Yuddha Kanda, 17:50-68}}
\end{tcolorbox}
\end{minipage}

\vspace{0.5em}
\subsubsection*{\textbf{Visual Rhythm Map:}}
\begin{table}[h!]
\centering
\renewcommand{\arraystretch}{1.2}
\setlength{\tabcolsep}{0pt}
\begin{tabularx}{\textwidth}{|*{16}{>{\centering\arraybackslash\hsize=1.0\hsize}X|}}
\hline
\multicolumn{6}{|>{\centering\arraybackslash\hsize=6.0\hsize}X|}{\textbf{\laghu{म}\guru{हा}\guru{बी}\laghu{र}} \par (6)} &
\multicolumn{4}{>{\centering\arraybackslash\hsize=4.0\hsize}X|}{\textbf{\guru{बि}\laghu{क्र}\laghu{म}} \par (4)} &
\multicolumn{6}{>{\centering\arraybackslash\hsize=6.0\hsize}X|}{\textbf{\laghu{ब}\laghu{ज}\guru{रं}\guru{गी}} \par (6)} \\
\hline
\end{tabularx}
\vspace{-1.5em}
\end{table}

\begin{table}[h!]
\centering
\renewcommand{\arraystretch}{1.2}
\setlength{\tabcolsep}{0pt}
\begin{tabularx}{\textwidth}{|*{16}{>{\centering\arraybackslash\hsize=1.0\hsize}X|}}
\hline
\multicolumn{3}{|>{\centering\arraybackslash\hsize=3.0\hsize}X|}{\textbf{\laghu{कु}\laghu{म}\laghu{ति}} \par (3)} &
\multicolumn{4}{>{\centering\arraybackslash\hsize=4.0\hsize}X|}{\textbf{\laghu{नि}\guru{वा}\laghu{र}} \par (4)} &
\multicolumn{3}{>{\centering\arraybackslash\hsize=3.0\hsize}X|}{\textbf{\laghu{सु}\laghu{म}\laghu{ति}} \par (3)} &
\multicolumn{2}{>{\centering\arraybackslash\hsize=2.0\hsize}X|}{\textbf{\guru{के}} \par (2)} &
\multicolumn{4}{>{\centering\arraybackslash\hsize=4.0\hsize}X|}{\textbf{\guru{सं}\guru{गी}} \par (4)} \\
\hline
\end{tabularx}
\vspace{-1.5em}
\end{table}

\vspace{1em}
\textbf{ \deva{सरल-संस्कृतम्}:}\\
 \deva{हे महावीर! त्वं पराक्रमी वज्राङ्गी (वज्र-समान-अङ्गवान्) च असि।}\\
 \deva{त्वं कुमतिं निवारयसि, सुमतीनां च सङ्गी (रक्षकः) असि॥}\\

O Great Hero! You are valiant and possess a body as strong as a thunderbolt.\\
You dispel evil thoughts, and you are the companion of good thoughts.

\vspace{1em}
\newpage
\textbf{\deva{प्रतिपदार्थः} (Meaning of each word):}
\begin{table}[h!]
\centering
\renewcommand{\arraystretch}{1.2}
\begin{tabularx}{\textwidth}{@{} l l X X @{}}
\toprule
\textbf{Awadhi} & \textbf{Sanskrit} & \textbf{English} & \textbf{Grammar \& Notes} \\
\midrule
\deva{महाबीर} / \telu{మహాబీర} / \textit{mahābīra}  &  \deva{महावीरः}  &  Great Hero  & \\ 
\deva{बिक्रम} \gotchaicon / \telu{బిక్రమ} / \textit{bikrama}  &  \deva{विक्रमः}  &  Valiant/Courageous  & \\ 
\deva{बजरंगी} / \telu{బజరంగీ} / \textit{bajaraṅgī}  &  \deva{वज्राङ्गी}  &  Thunderbolt body  & Swara-Bhakti splitting \\
\deva{कुमति} / \telu{కుమతి} / \textit{kumati}  &  \deva{कुमतिम्}  &  Evil intellect  &   \\
\deva{निवार} \gotchaicon / \telu{నివార} / \textit{nivāra}  &  \deva{निवारकः}  &  Dispeller / Remover  & Meter trap: short 'i' \\
\deva{सुमति} / \telu{సుమతి} / \textit{sumati}  &  \deva{सुमतीनाम्}  &  Good intellect  &   \\
\deva{संगी} / \telu{సంగీ} / \textit{saṅgī}  &  \deva{सङ्गी / सखा}  &  Companion  &   \\
\bottomrule
\end{tabularx}
\end{table}

\begin{tcolorbox}[gotchabox]
\begin{itemize}
    \item \textbf{The Triple \deva{व} $\rightarrow$ \deva{ब} Shift:} This line is the ultimate test of the Awadhi rule! \deva{वीर} (Veera) $\rightarrow$ \deva{बीर} (Beera). \deva{विक्रम} (Vikrama) $\rightarrow$ \deva{बिक्रम} (Bikrama). 
    \item \textbf{Swara-Bhakti (\deva{वज्र $\rightarrow$ बजर}):} The harsh conjunct 'jra' in \deva{वज्र} (Vajra) is broken apart by inserting a neutralizing 'a' vowel, stretching it from Vaj-ra to Ba-ja-ra. This preserves the mouth's flow.
    \item \textbf{Meter Trap:} Keep the 'i' short in \deva{निवार} (Nivaara). If you elongate it to \deva{नीवार} (Neevaara), you add a beat and break the 16-matra math!
\end{itemize}
\end{tcolorbox}

\newpage


% ----- CHAUPAI 4 -----

\newpage

\subsection*{\deva{चतुर्थ चौपाई} (Fourth Chaupai)}
\addcontentsline{toc}{subsection}{\deva{चतुर्थ चौपाई} (Fourth Chaupai) — \deva{कंचन बरन...}}

\begin{minipage}[t]{0.55\textwidth}
\vspace{0pt}

{\large \linespread{1.5}\selectfont
\deva{कंचन बरन बिराज सुबेसा।}\\
\deva{कानन कुंडल कुंचित केसा॥}
\par}

\vspace{0.5em}

{\large \linespread{1.5}\selectfont
\telu{కంచన బరన బిరాజ సుబేసా।}\\
\telu{కానన కుండల కుంచిత కేసా॥}
\par}

\vspace{0.5em}

{\large \linespread{1.3}\selectfont
\textit{kañcana barana birāja subesā |}\\
\textit{kānana kuṇḍala kuñcita kesā ||}
\par}
\end{minipage}%
\hfill
\begin{minipage}[t]{0.42\textwidth}
\vspace{0pt}
\begin{tcolorbox}[anchorbox]
\textbf{The Reward of Education:} Why the golden glow and beautiful earrings (\textit{Kanan Kundal})? According to tradition, when Hanuman completed his immense, relentless education under the Sun God, Surya was so impressed by his student that he gifted him a fraction of his own blinding brilliance and divine earrings as a graduation gift. True radiance is not just physical, it is the result of dedication, hard work and earned through good education.
\vspace{0.5em}\newline His brilliant, golden radiance was not inherited, but earned through the disciplined fire of his actions.\newline\null\hfill\textit{Kishkindha Kanda, 66}
\end{tcolorbox}
\end{minipage}

\vspace{0.5em}
\subsubsection*{\textbf{Visual Rhythm Map:}}
\begin{table}[h!]
\centering
\renewcommand{\arraystretch}{1.2}
\setlength{\tabcolsep}{0pt}
\begin{tabularx}{\textwidth}{|*{16}{>{\centering\arraybackslash\hsize=1.0\hsize}X|}}
\hline
\multicolumn{4}{|>{\centering\arraybackslash\hsize=4.0\hsize}X|}{\textbf{\guru{कं}\laghu{च}\laghu{न}} \par (4)} &
\multicolumn{3}{>{\centering\arraybackslash\hsize=3.0\hsize}X|}{\textbf{\laghu{ब}\laghu{र}\laghu{न}} \par (3)} &
\multicolumn{4}{>{\centering\arraybackslash\hsize=4.0\hsize}X|}{\textbf{\laghu{बि}\guru{रा}\laghu{ज}} \par (4)} &
\multicolumn{5}{>{\centering\arraybackslash\hsize=5.0\hsize}X|}{\textbf{\laghu{सु}\guru{बे}\guru{सा}} \par (5)} \\
\hline
\end{tabularx}
\vspace{-1.5em}
\end{table}

\begin{table}[h!]
\centering
\renewcommand{\arraystretch}{1.2}
\setlength{\tabcolsep}{0pt}
\begin{tabularx}{\textwidth}{|*{16}{>{\centering\arraybackslash\hsize=1.0\hsize}X|}}
\hline
\multicolumn{4}{|>{\centering\arraybackslash\hsize=4.0\hsize}X|}{\textbf{\guru{का}\laghu{न}\laghu{न}} \par (4)} &
\multicolumn{4}{>{\centering\arraybackslash\hsize=4.0\hsize}X|}{\textbf{\guru{कुं}\laghu{ड}\laghu{ल}} \par (4)} &
\multicolumn{4}{>{\centering\arraybackslash\hsize=4.0\hsize}X|}{\textbf{\guru{कुं}\laghu{चि}\laghu{त}} \par (4)} &
\multicolumn{4}{>{\centering\arraybackslash\hsize=4.0\hsize}X|}{\textbf{\guru{के}\guru{सा}} \par (4)} \\
\hline
\end{tabularx}
\vspace{-1.5em}
\end{table}

\vspace{1em}
\textbf{ \deva{सरल-संस्कृतम्}:}\\
 \deva{भवतः कञ्चन-वर्णः (सुवर्ण-वर्णः) अस्ति, भवान् सुन्दर-वेषेण च विराजते।}\\
 \deva{भवतः कर्णयोः कुण्डले स्तः, केशाः च कुञ्चिताः सन्ति॥}\\

Your complexion is golden, and you are adorned with beautiful attire.\\
You wear earrings on your ears, and your hair is curly.


\vspace{1em}
\newpage
\textbf{\deva{प्रतिपदार्थः} (Meaning of each word):}
\begin{table}[h!]
\centering
\renewcommand{\arraystretch}{1.2}
\begin{tabularx}{\textwidth}{@{} l l X X @{}}
\toprule
\textbf{Awadhi} & \textbf{Sanskrit} & \textbf{English} & \textbf{Grammar \& Notes} \\
\midrule
\deva{कंचन} / \telu{కంచన}  &  \deva{कञ्चन / सुवर्णम्}  &  Gold  &   \\
\deva{बरन} / \telu{బరన}  &  \deva{वर्णः}  &  Complexion / Color  & \\ 
\deva{बिराज} / \telu{బిరాజ}  &  \deva{विराजते}  &  Shines / Is adorned  & \\ 
\deva{सुबेसा} \gotchaicon / \telu{సుబేసా}  &  \deva{सुवेषः}  &  Beautiful attire  & \\ 
\deva{कानन} / \telu{కానన}  &  \deva{कर्णयोः}  &  In the ears  & \\ 
\deva{कुंडल} / \telu{కుండల}  &  \deva{कुण्डले}  &  Earrings  &   \\
\deva{कुंचित} / \telu{కుంచిత}  &  \deva{कुञ्चिताः}  &  Curly  &   \\
\deva{केसा} / \telu{కేసా}  &  \deva{केशाः}  &  Hair  & \\ 
\bottomrule
\end{tabularx}
\end{table}

\begin{tcolorbox}[gotchabox]
\begin{itemize}
    \item \textbf{Phonetic Shifts:} \deva{सुवेष} (Suvesha) becomes \deva{सुबेसा} (Subesaa) — exhibiting both the \deva{व} $\rightarrow$ \deva{ब} and \deva{श/ष} $\rightarrow$ \deva{स} shifts simultaneously! Likewise, \deva{केश} (Kesha) becomes \deva{केस} (Kesa).
\end{itemize}
\end{tcolorbox}



\newpage


% ----- CHAUPAI 5 -----

\newpage

\subsection*{\deva{पञ्चम चौपाई} (Fifth Chaupai)}
\addcontentsline{toc}{subsection}{\deva{पञ्चम चौपाई} (Fifth Chaupai) — \deva{हाथ बज्र...}}

\begin{minipage}[t]{0.55\textwidth}
\vspace{0pt}

{\large \linespread{1.5}\selectfont
\deva{हाथ बज्र औ ध्वजा बिराजै।}\\
\deva{काँधे मूँज जनेऊ साजै॥}
\par}

\vspace{0.5em}

{\large \linespread{1.5}\selectfont
\telu{హాథ బజ్ర ఔ ధ్వజా బిరాజై।}\\
\telu{కాఁధే మూఁజ జనేఊ సాజై॥}
\par}

\vspace{0.5em}

{\large \linespread{1.3}\selectfont
\textit{hātha bajra au dhvajā birājai |}\\
\textit{kāṃdhe mūṃja janeū sājai ||}
\par}
\end{minipage}%
\hfill
\begin{minipage}[t]{0.42\textwidth}
\vspace{0pt}
\begin{tcolorbox}[anchorbox]
\textbf{The Scholar and the Protector:} He holds a devastating weapon (\textit{Bajra}) in his hand, but wears the sacred thread (\textit{Janeu}) of a peaceful scholar on his shoulder. This is the ultimate balance we should all strive for: possessing a highly educated, peaceful mind, while remaining strong enough to fiercely protect others from injustice.
\vspace{0.5em}\newline He carried the absolute power of a warrior beneath the humble, unthreatening garb of a dedicated scholar.\newline\null\hfill\href{https://www.valmiki.iitk.ac.in/sloka?field_kanda_tid=4&language=dv&field_sarga_value=3}{\textit{Kishkindha Kanda, 3:2}}
\end{tcolorbox}
\end{minipage}

\vspace{0.5em}
\subsubsection*{\textbf{Visual Rhythm Map:}}
\begin{table}[h!]
\centering
\renewcommand{\arraystretch}{1.2}
\setlength{\tabcolsep}{0pt}
\begin{tabularx}{\textwidth}{|*{16}{>{\centering\arraybackslash\hsize=1.0\hsize}X|}}
\hline
\multicolumn{3}{|>{\centering\arraybackslash\hsize=3.0\hsize}X|}{\textbf{\guru{हा}\laghu{थ}} \par (3)} &
\multicolumn{3}{>{\centering\arraybackslash\hsize=3.0\hsize}X|}{\textbf{\guru{ब}\laghu{ज्र}} \par (3)} &
\multicolumn{2}{>{\centering\arraybackslash\hsize=2.0\hsize}X|}{\textbf{\guru{औ}} \par (2)} &
\multicolumn{3}{>{\centering\arraybackslash\hsize=3.0\hsize}X|}{\textbf{\laghu{ध्व}\guru{जा}} \par (3)} &
\multicolumn{5}{>{\centering\arraybackslash\hsize=5.0\hsize}X|}{\textbf{\laghu{बि}\guru{रा}\guru{जै}} \par (5)} \\
\hline
\end{tabularx}
\vspace{-1.5em}
\end{table}

\begin{table}[h!]
\centering
\renewcommand{\arraystretch}{1.2}
\setlength{\tabcolsep}{0pt}
\begin{tabularx}{\textwidth}{|*{16}{>{\centering\arraybackslash\hsize=1.0\hsize}X|}}
\hline
\multicolumn{4}{|>{\centering\arraybackslash\hsize=4.0\hsize}X|}{\textbf{\guru{काँ}\guru{धे}} \par (4)} &
\multicolumn{3}{>{\centering\arraybackslash\hsize=3.0\hsize}X|}{\textbf{\guru{मूँ}\laghu{ज}} \par (3)} &
\multicolumn{5}{>{\centering\arraybackslash\hsize=5.0\hsize}X|}{\textbf{\laghu{ज}\guru{ने}\guru{ऊ}} \par (5)} &
\multicolumn{4}{>{\centering\arraybackslash\hsize=4.0\hsize}X|}{\textbf{\guru{सा}\guru{जै}} \par (4)} \\
\hline
\end{tabularx}
\vspace{-1.5em}
\end{table}

\vspace{1em}
\textbf{ \deva{सरल-संस्कृतम्}:}\\
 \deva{भवतः हस्ते वज्रं ध्वजः च विराजते।}\\
 \deva{भवतः स्कन्धे मौञ्जी-यज्ञोपवीतं शोभते॥}\\

In your hand shines the thunderbolt and the banner.\\
On your shoulder, the sacred thread of Munja grass looks beautiful.

\vspace{1em}
\newpage
\textbf{\deva{प्रतिपदार्थः} (Meaning of each word):}
\begin{table}[h!]
\centering
\renewcommand{\arraystretch}{1.2}
\begin{tabularx}{\textwidth}{@{} l l X X @{}}
\toprule
\textbf{Awadhi} & \textbf{Sanskrit} & \textbf{English} & \textbf{Grammar \& Notes} \\
\midrule
\deva{हाथ} / \telu{హాథ}  &  \deva{हस्ते}  &  In hands  &   \\
\deva{बज्र} \gotchaicon / \telu{బజ్ర}  &  \deva{वज्रम्}  &  Thunderbolt  & \\ 
\deva{औ} / \telu{ఔ}  &  \deva{और / च}  &  And  &   \\
\deva{ध्वजा} / \telu{ధ్వజా}  &  \deva{ध्वजः}  &  Flag / Banner  &   \\
\deva{बिराजै} / \telu{బిరాజై}  &  \deva{विराजते}  &  Shines  &   \\
\deva{काँधे} \gotchaicon / \telu{కాఁధే} & \deva{स्कन्धे} & On the shoulder & \\
\deva{मूँज} \gotchaicon / \telu{మూఁజ} & \deva{मौञ्जी} & Munja grass & Moon-dot (Chandrabindu) \\
\deva{जनेऊ} / \telu{జనేఊ}  &  \deva{यज्ञोपवीतम्}  &  Sacred thread  &   \\
\deva{साजै} / \telu{సాజై}  &  \deva{सज्जते / शोभते}  &  Adorns / Looks beautiful  &   \\
\bottomrule
\end{tabularx}
\end{table}

\begin{tcolorbox}[gotchabox]
\begin{itemize}
    \item \textbf{Nasalization:} Pay close attention to the \deva{चन्द्रबिन्दु} (Chandrabindu / Moon-dot) on \deva{काँधे} (Kaandhe) and \deva{मूँज} (Moonja). This nasalizes the vowel without adding the extra weight of a full consonant, preserving the 16-beat meter.
    \item \textbf{Conjunct Consonant Weight (\deva{बज्र}):} The `a` in `ba` becomes a 2-beat Guru because it precedes the conjunct `jra`!
\end{itemize}
\end{tcolorbox}

\newpage


% ----- CHAUPAI 6 -----

\newpage

\subsection*{\deva{षष्ठ चौपाई} (Sixth Chaupai)}
\addcontentsline{toc}{subsection}{\deva{षष्ठ चौपाई} (Sixth Chaupai) — \deva{संकर सुवन...}}

\begin{minipage}[t]{0.55\textwidth}
\vspace{0pt}

{\large \linespread{1.5}\selectfont
\deva{संकर सुवन केसरी नंदन।}\\
\deva{तेज प्रताप महा जग बंदन॥}
\par}

\vspace{0.5em}

{\large \linespread{1.5}\selectfont
\telu{సంకర సువన కేసరీ నందన।}\\
\telu{తేజ ప్రతాప మహా జగ బందన॥}
\par}

\vspace{0.5em}

{\large \linespread{1.3}\selectfont
\textit{saṅkara suvana kesarī nandana |}\\
\textit{teja pratāpa mahā jaga bandana ||}
\par}
\end{minipage}%
\hfill
\begin{minipage}[t]{0.42\textwidth}
\vspace{0pt}
\begin{tcolorbox}[anchorbox]
\textbf{The Radiance of Grounded Service:} \textit{Shankara} represents the highest spiritual intellect and ego-less discipline, while \textit{Kesari} represents a highly practical, grounded, forest-dwelling leader. Hanuman embodies the perfect balance of both: he possesses the brilliant mind of a supreme thinker, yet uses his fierce protective energy (\textit{Teja Pratapa}) entirely for the selfless service of others rather than personal gain. The world reveres him (\textit{Jaga Bandana}) because he proves that true majestic radiance comes from balancing high-minded vision with practical, ego-less dedication to your community.
\vspace{0.5em}\newline Jambavan honors his earthly, grounded origins, proving true radiance stems from one's actions, not just their birth.\newline\null\hfill\href{https://www.valmiki.iitk.ac.in/sloka?field_kanda_tid=4&language=dv&field_sarga_value=66}{\textit{Kishkindha Kanda, 66:29}}
\end{tcolorbox}
\end{minipage}

\vspace{0.5em}
\subsubsection*{\textbf{Visual Rhythm Map:}}
\begin{table}[h!]
\centering
\renewcommand{\arraystretch}{1.2}
\setlength{\tabcolsep}{0pt}
\begin{tabularx}{\textwidth}{|*{16}{>{\centering\arraybackslash\hsize=1.0\hsize}X|}}
\hline
\multicolumn{4}{|>{\centering\arraybackslash\hsize=4.0\hsize}X|}{\textbf{\guru{सं}\laghu{क}\laghu{र}} \par (4)} &
\multicolumn{3}{>{\centering\arraybackslash\hsize=3.0\hsize}X|}{\textbf{\laghu{सु}\laghu{व}\laghu{न}} \par (3)} &
\multicolumn{5}{>{\centering\arraybackslash\hsize=5.0\hsize}X|}{\textbf{\guru{के}\laghu{स}\guru{री}} \par (5)} &
\multicolumn{4}{>{\centering\arraybackslash\hsize=4.0\hsize}X|}{\textbf{\guru{नं}\laghu{द}\laghu{न}} \par (4)} \\
\hline
\end{tabularx}
\vspace{-1.5em}
\end{table}

\begin{table}[h!]
\centering
\renewcommand{\arraystretch}{1.2}
\setlength{\tabcolsep}{0pt}
\begin{tabularx}{\textwidth}{|*{16}{>{\centering\arraybackslash\hsize=1.0\hsize}X|}}
\hline
\multicolumn{3}{|>{\centering\arraybackslash\hsize=3.0\hsize}X|}{\textbf{\guru{ते}\laghu{ज}} \par (3)} &
\multicolumn{4}{>{\centering\arraybackslash\hsize=4.0\hsize}X|}{\textbf{\laghu{प्र}\guru{ता}\laghu{प}} \par (4)} &
\multicolumn{3}{>{\centering\arraybackslash\hsize=3.0\hsize}X|}{\textbf{\laghu{म}\guru{हा}} \par (3)} &
\multicolumn{2}{>{\centering\arraybackslash\hsize=2.0\hsize}X|}{\textbf{\laghu{ज}\laghu{ग}} \par (2)} &
\multicolumn{4}{>{\centering\arraybackslash\hsize=4.0\hsize}X|}{\textbf{\guru{बं}\laghu{द}\laghu{न}} \par (4)} \\
\hline
\end{tabularx}
\vspace{-1.5em}
\end{table}

\vspace{1em}
\textbf{ \deva{सरल-संस्कृतम्}:}\\
 \deva{भवान् शङ्करस्य अवतारः, केसरी-पुत्रः च असि।}\\
 \deva{भवतः तेजः प्रतापः च महान् अस्ति, सम्पूर्णं जगत् भवन्तं वन्दते (नमति)॥}\\


You are the incarnation of Shiva and the son of Kesari.\\
Your radiance and valor are immense, and you are revered by the whole universe.


\vspace{1em}
\newpage
\textbf{\deva{प्रतिपदार्थः} (Meaning of each word):}
\begin{table}[h!]
\centering
\renewcommand{\arraystretch}{1.2}
\begin{tabularx}{\textwidth}{@{} l l X X @{}}
\toprule
\textbf{Awadhi} & \textbf{Sanskrit} & \textbf{English} & \textbf{Grammar \& Notes} \\
\midrule
\deva{संकर} \gotchaicon / \telu{సంకర} / \textit{saṅkara}  &  \deva{शङ्करस्य}  &  Of Shiva  & \\ 
\deva{संकर सुवन} / \telu{సంకర సువన} / \textit{saṅkara suvana}  &  \deva{शङ्करस्य पुत्रः (हनुमान्)}  &  Son of Shiva  &   \\
\deva{केसरी नंदन} / \telu{కేసరీ నందన} / \textit{kesarī nandana}  &  \deva{केसरिणः नन्दनः}  &  Son of Kesari  & \\ 
\deva{तेज} / \telu{తేజ} / \textit{teja}  &  \deva{तेजस्वित्वम्}  &  Radiance / Energy  &   \\
\deva{प्रताप} / \telu{ప్రతాప} / \textit{pratāpa}  &  \deva{प्रतापः}  &  Valor / Prowess  &   \\
\deva{महा} / \telu{మహా} / \textit{mahā}  &  \deva{महान्}  &  Great  &   \\
\deva{जग} / \telu{జగ} / \textit{jaga} & \deva{जगति} & In the world & \\
\deva{बंदन} / \telu{బందన} / \textit{bandana} & \deva{वन्दनीयः} & Adored / Worshiped & \\
 
\bottomrule
\end{tabularx}
\end{table}

\begin{tcolorbox}[gotchabox]
\begin{itemize}
    \item \textbf{Phonetic Shifts:} Sanskrit \deva{शङ्कर} (Shankara) seamlessly shifts into the Awadhi \deva{संकर} (Sankara) losing the palatal 'sh'. Likewise, \deva{वन्दन} (Vandana) hardens into \deva{बंदन} (Bandana) with the \deva{व} $\rightarrow$ \deva{ब} shift.
    \item \textbf{The Secret of \deva{सुवन} (The Glide / \deva{श्रुति}):} If \deva{वन्दन} becomes \deva{बंदन}, why didn't \deva{सुवन} become \deva{सुबन}? The word \deva{सुवन} comes from the Sanskrit root \deva{सू} (to beget) plus the noun suffix \deva{अन}. Saying \deva{सू} and \deva{अन} rapidly together naturally creates a soft 'w/v' sound as a breath-bridge (a \textit{Glide} or \deva{व-श्रुति}). Because this \deva{व} is just a fluid phonetic bridge between two vowels—and not a heavy structural consonant—the Awadhi speakers kept it soft instead of hammering it into a hard \deva{ब}.
\end{itemize}
\end{tcolorbox}

\newpage



% ==========================================
\section*{Group 3: The Supreme Servant of Rama}
\addcontentsline{toc}{section}{Group 3: The Supreme Servant of Rama}
% ==========================================

% ----- CHAUPAI 7 -----
\subsection*{\deva{सप्तम चौपाई} (Seventh Chaupai)}
\addcontentsline{toc}{subsection}{\deva{सप्तम चौपाई} (Seventh Chaupai) — \deva{बिद्यावान गुनी...}}

\begin{minipage}[t]{0.55\textwidth}
\vspace{0pt}

{\large \linespread{1.5}\selectfont
\deva{बिद्यावान गुनी अति चातुर।}\\
\deva{राम काज करिबे को आतुर॥}
\par}

\vspace{0.5em}

{\large \linespread{1.5}\selectfont
\telu{బిద్యావాన గునీ అతి చాతుర।}\\
\telu{రామ కాజ కరిబే కో ఆతుర॥}
\par}

\vspace{0.5em}

{\large \linespread{1.3}\selectfont
\textit{bidyāvāna gunī ati cātura |}\\
\textit{rāma kāja karibe ko ātura ||}
\par}
\end{minipage}%
\hfill
\begin{minipage}[t]{0.42\textwidth}
\vspace{0pt}

\end{minipage}

\vspace{0.5em}
\subsubsection*{\textbf{Visual Rhythm Map:}}
\begin{table}[h!]
\centering
\renewcommand{\arraystretch}{1.2}
\setlength{\tabcolsep}{0pt}
\begin{tabularx}{\textwidth}{|*{16}{>{\centering\arraybackslash\hsize=1.0\hsize}X|}}
\hline
\multicolumn{7}{|>{\centering\arraybackslash\hsize=7.0\hsize}X|}{\textbf{\guru{बि}\guru{द्या}\guru{वा}\laghu{न}} \par (7)} &
\multicolumn{3}{>{\centering\arraybackslash\hsize=3.0\hsize}X|}{\textbf{\laghu{गु}\guru{नी}} \par (3)} &
\multicolumn{2}{>{\centering\arraybackslash\hsize=2.0\hsize}X|}{\textbf{\laghu{अ}\laghu{ति}} \par (2)} &
\multicolumn{4}{>{\centering\arraybackslash\hsize=4.0\hsize}X|}{\textbf{\guru{चा}\laghu{तु}\laghu{र}} \par (4)} \\
\hline
\end{tabularx}
\vspace{-1.5em}
\end{table}

\begin{table}[h!]
\centering
\renewcommand{\arraystretch}{1.2}
\setlength{\tabcolsep}{0pt}
\begin{tabularx}{\textwidth}{|*{16}{>{\centering\arraybackslash\hsize=1.0\hsize}X|}}
\hline
\multicolumn{3}{|>{\centering\arraybackslash\hsize=3.0\hsize}X|}{\textbf{\guru{रा}\laghu{म}} \par (3)} &
\multicolumn{3}{>{\centering\arraybackslash\hsize=3.0\hsize}X|}{\textbf{\guru{का}\laghu{ज}} \par (3)} &
\multicolumn{4}{>{\centering\arraybackslash\hsize=4.0\hsize}X|}{\textbf{\laghu{क}\laghu{रि}\guru{बे}} \par (4)} &
\multicolumn{2}{>{\centering\arraybackslash\hsize=2.0\hsize}X|}{\textbf{\guru{को}} \par (2)} &
\multicolumn{4}{>{\centering\arraybackslash\hsize=4.0\hsize}X|}{\textbf{\guru{आ}\laghu{तु}\laghu{र}} \par (4)} \\
\hline
\end{tabularx}
\vspace{-1.5em}
\end{table}

\vspace{1em}
\textbf{ \deva{सरल-संस्कृतम्}:}\\
 \deva{भवान् विद्यावान्, गुणवान्, अतीव चतुरः च अस्ति।}\\
 \deva{श्रीरामस्य कार्यं कर्तुं भवान् सदैव आतुरः (उत्सुकः) अस्ति।}\\


You are highly learned, virtuous, and extremely clever/skillful.\\
You are always eager and ready to execute Lord Rama's work.


\vspace{1em}
\newpage
\textbf{\deva{प्रतिपदार्थः} (Meaning of each word):}
\begin{table}[h!]
\centering
\renewcommand{\arraystretch}{1.2}
\begin{tabularx}{\textwidth}{@{} l l X X @{}}
\toprule
\textbf{Awadhi} & \textbf{Sanskrit} & \textbf{English} & \textbf{Grammar \& Notes} \\
\midrule
\deva{बिद्यावान} \gotchaicon / \telu{బిద్యావాన}  &  \deva{विद्यावान्}  &  Learned / Scholar  & \\ 
\deva{गुनी} / \telu{గునీ}  &  \deva{गुणी (गुणवान्)}  &  Virtuous  &   \\
\deva{अति} / \telu{అతి}  &  \deva{अतीव}  &  Extremely  &   \\
\deva{चातुर} / \telu{చాతుర}  &  \deva{चतुरः}  &  Clever  &   \\
\deva{राम काज} / \telu{రామ కాజ}  &  \deva{श्रीरामस्य कार्यम्}  &  Rama's work  & Prakrit shift () \\ 
\deva{करिबे को} / \telu{కరిబే కో}  &  \deva{कर्तुम्}  &  To do (Infinitive)  & Awadhi infinitive \\
\deva{आतुर} / \telu{ఆతుర}  &  \deva{आतुरः (उत्सुकः)}  &  Eager / Ready  &   \\
\bottomrule
\end{tabularx}
\end{table}

\begin{tcolorbox}[gotchabox]
\begin{itemize}
    \item \textbf{The "Word-Initial" vs "Suffix" Rule (\deva{बिद्यावान}):} Why \deva{बिद्यावान} and not \deva{बिद्याबान}? Because the first $v$ carries the heavy front door weight of the word, it requires maximum articulatory force and hardens into the plosive $b$. The $v$ in $-vaan$ is a possession suffix (\deva{वतुँप् प्रत्यय}). Because it occurs at the end of the word where the mouth naturally relaxes (Mukha-Sukha / Lenition), it stays as the gliding breathy $v$.
\end{itemize}
\end{tcolorbox}

\begin{tcolorbox}[poetrybox]
\textbf{Praasa (Internal/External Rhyme):}\\
Line 1 focuses on his internal state (Smart / \deva{चातुर}). Line 2 focuses on his external action (Eager / \deva{आतुर}). Notice the perfect rhythm and rhyme between Chaatura and Aatura!
\end{tcolorbox}


\newpage


% ----- CHAUPAI 8 -----

\newpage

\subsection*{\deva{अष्टम चौपाई} (Eighth Chaupai)}
\addcontentsline{toc}{subsection}{\deva{अष्टम चौपाई} (Eighth Chaupai) — \deva{प्रभु चरित्र...}}

\begin{minipage}[t]{0.55\textwidth}
\vspace{0pt}

{\large \linespread{1.5}\selectfont
\deva{प्रभु चरित्र सुनिबे को रसिया।}\\
\deva{राम लखन सीता मन बसिया॥}
\par}

\vspace{0.5em}

{\large \linespread{1.5}\selectfont
\telu{ప్రభు చరిత్ర సునిబే కో రసియా।}\\
\telu{రామ లఖన సీతా మన బసియా॥}
\par}

\vspace{0.5em}

{\large \linespread{1.3}\selectfont
\textit{prabhu caritra sunibe ko rasiyā |}\\
\textit{rāma lakhana sītā mana basiyā ||}
\par}
\end{minipage}%
\hfill
\begin{minipage}[t]{0.42\textwidth}
\vspace{0pt}
\begin{tcolorbox}[anchorbox]
\textbf{The Art of Listening:} He is praised for being exceptionally eager to listen (\textit{Sunibe ko Rasiya}). In the *Valmiki Ramayana* (*Kishkindha Kanda, 3:28-33*), Lord Rama notes that Hanuman's flawless speech is proof of his highly disciplined, attentive listening as a student of the Vedas. True mastery in communication and leadership begins with the patient, active absorption of knowledge.
\vspace{0.5em}\newline Hanuman's legendary eloquence was the direct result of his deep, attentive listening and study under his mentors.\newline\null\hfill\href{https://www.valmiki.iitk.ac.in/sloka?field_kanda_tid=4&language=dv&field_sarga_value=3}{\textit{Kishkindha Kanda, 3:28-33}}
\end{tcolorbox}
\end{minipage}

\vspace{0.5em}
\subsubsection*{\textbf{Visual Rhythm Map:}}
\begin{table}[h!]
\centering
\renewcommand{\arraystretch}{1.2}
\setlength{\tabcolsep}{0pt}
\begin{tabularx}{\textwidth}{|*{16}{>{\centering\arraybackslash\hsize=1.0\hsize}X|}}
\hline
\multicolumn{2}{|>{\centering\arraybackslash\hsize=2.0\hsize}X|}{\textbf{\laghu{प्र}\laghu{भु}} \par (2)} &
\multicolumn{4}{>{\centering\arraybackslash\hsize=4.0\hsize}X|}{\textbf{\laghu{च}\guru{रि}\laghu{त्र}} \par (4)} &
\multicolumn{4}{>{\centering\arraybackslash\hsize=4.0\hsize}X|}{\textbf{\laghu{सु}\laghu{नि}\guru{बे}} \par (4)} &
\multicolumn{2}{>{\centering\arraybackslash\hsize=2.0\hsize}X|}{\textbf{\guru{को}} \par (2)} &
\multicolumn{4}{>{\centering\arraybackslash\hsize=4.0\hsize}X|}{\textbf{\laghu{र}\laghu{सि}\guru{या}} \par (4)} \\
\hline
\end{tabularx}
\vspace{-1.5em}
\end{table}

\begin{table}[h!]
\centering
\renewcommand{\arraystretch}{1.2}
\setlength{\tabcolsep}{0pt}
\begin{tabularx}{\textwidth}{|*{16}{>{\centering\arraybackslash\hsize=1.0\hsize}X|}}
\hline
\multicolumn{3}{|>{\centering\arraybackslash\hsize=3.0\hsize}X|}{\textbf{\guru{रा}\laghu{म}} \par (3)} &
\multicolumn{3}{>{\centering\arraybackslash\hsize=3.0\hsize}X|}{\textbf{\laghu{ल}\laghu{ख}\laghu{न}} \par (3)} &
\multicolumn{4}{>{\centering\arraybackslash\hsize=4.0\hsize}X|}{\textbf{\guru{सी}\guru{ता}} \par (4)} &
\multicolumn{2}{>{\centering\arraybackslash\hsize=2.0\hsize}X|}{\textbf{\laghu{म}\laghu{न}} \par (2)} &
\multicolumn{4}{>{\centering\arraybackslash\hsize=4.0\hsize}X|}{\textbf{\laghu{ब}\laghu{सि}\guru{या}} \par (4)} \\
\hline
\end{tabularx}
\vspace{-1.5em}
\end{table}

\vspace{1em}
\textbf{ \deva{सरल-संस्कृतम्}:}\\
 \deva{भवान् प्रभोः (श्रीरामस्य) चरित्रं श्रोतुं रसिकः अस्ति।}\\
 \deva{श्रीरामः, लक्ष्मणः, सीता च भवतः मनसि वसन्ति।}\\


You delight in listening to the stories and glories of the Lord.\\
Rama, Lakshmana, and Sita reside in your mind.


\vspace{1em}
\newpage
\textbf{\deva{प्रतिपदार्थः} (Meaning of each word):}
\begin{table}[h!]
\centering
\renewcommand{\arraystretch}{1.2}
\begin{tabularx}{\textwidth}{@{} l l X X @{}}
\toprule
\textbf{Awadhi} & \textbf{Sanskrit} & \textbf{English} & \textbf{Grammar \& Notes} \\
\midrule
\deva{प्रभु} / \telu{ప్రభు} / \textit{prabhu}  &  \deva{प्रभोः}  &  Lord's  &   \\
\deva{चरित्र} / \telu{చరిత్ర} / \textit{caritra}  &  \deva{चरित्रम् / कथाम्}  &  Character / Stories  &   \\
\deva{सुनिबे को} / \telu{సునిబే కో} / \textit{sunibe ko}  &  \deva{श्रोतुम्}  &  To listen  & Awadhi infinitive verb \\
\deva{रसिया} / \telu{రసియా} / \textit{rasiyā}  &  \deva{रसिकः}  &  One who relishes  &   \\
\deva{राम} / \telu{రామ} / \textit{rāma}  &  \deva{श्रीरामः}  &  Rama  &   \\
\deva{लखन} \gotchaicon / \telu{లఖన} / \textit{lakhana}  &  \deva{लक्ष्मणः}  &  Lakshmana  & cluster simplification \\ 
\deva{सीता} \gotchaicon / \telu{సీతా} / \textit{sītā}  &  \deva{सीता} \gotchaicon  &  Sita  &   \\
\deva{मन} / \telu{మన} / \textit{mana}  &  \deva{मनसि}  &  In the mind  &   \\
\deva{बसिया} \gotchaicon / \telu{బసియా} / \textit{basiyā}  &  \deva{వसन्ति / వాసీ}  &  Reside / Dweller  & \\ 
\bottomrule
\end{tabularx}
\end{table}

\begin{tcolorbox}[gotchabox]
\begin{itemize}
    \item \textbf{Consonant Melt (\deva{लक्ष्मण $\rightarrow$ लखन}):} The harsh consonant cluster `kshma` (\deva{क्ष्म}) in Sanskrit melts down into the softer, breathy `kha` (\deva{ख}) in Awadhi, making it much easier to sing! The Gita Press version uses \deva{लषन} instead of \deva{लखन}. Both carry the same metrical value.
    \item \textbf{The \deva{व} $\rightarrow$ \deva{ब} Shift (\deva{बसिया}):} The Sanskrit root for dwelling is \deva{वस्} (Vas). True to form, the Awadhi immediately hardens that initial $v$ into a $b$, giving us \deva{बसिया} (Basiyaa).
    
\end{itemize}
\end{tcolorbox}



\newpage


% ----- CHAUPAI 9 -----

\newpage

\subsection*{\deva{नवम चौपाई} (Ninth Chaupai)}
\addcontentsline{toc}{subsection}{\deva{नवम चौपाई} (Ninth Chaupai) — \deva{सूक्ष्म रूप...}}

\begin{minipage}[t]{0.55\textwidth}
\vspace{0pt}

{\large \linespread{1.5}\selectfont
\deva{सूक्ष्म रूप धरि सियहिं दिखावा।}\\
\deva{बिकट रूप धरि लंक जरावा॥}
\par}

\vspace{0.5em}

{\large \linespread{1.5}\selectfont
\telu{సూక్ష్మ రూప ధరి సియహిఁ దిఖావా।}\\
\telu{బికట రూప ధరి లంక జరావా॥}
\par}

\vspace{0.5em}

{\large \linespread{1.3}\selectfont
\textit{sūkṣma rūpa dhari siyahiṃ dikhāvā |}\\
\textit{bikaṭa rūpa dhari laṃka jarāvā ||}
\par}
\end{minipage}%
\hfill
\begin{minipage}[t]{0.42\textwidth}
\vspace{0pt}
\begin{tcolorbox}[anchorbox]
\textbf{Emotional Intelligence and Scale:} Hanuman knew exactly how to scale his approach. He approached allies in a tiny, unthreatening form (\textit{Sukshma Rupa}) to build trust and offer comfort, but assumed a massive, formidable form (\textit{Bikata Rupa}) when he needed to defend himself. Emotional intelligence requires adapting your approach to match the needs of the situation.
\vspace{0.5em}\newline He made himself intentionally unthreatening to sneak into Lanka, showing that true strength knows when to be subtle.\newline\null\hfill\href{https://www.valmiki.iitk.ac.in/sloka?field_kanda_tid=5&language=dv&field_sarga_value=2}{\textit{Sundara Kanda, 2:49}}
\end{tcolorbox}
\end{minipage}

\vspace{0.5em}
\subsubsection*{\textbf{Visual Rhythm Map:}}
\begin{table}[h!]
\centering
\renewcommand{\arraystretch}{1.2}
\setlength{\tabcolsep}{0pt}
\begin{tabularx}{\textwidth}{|*{16}{>{\centering\arraybackslash\hsize=1.0\hsize}X|}}
\hline
\multicolumn{3}{|>{\centering\arraybackslash\hsize=3.0\hsize}X|}{\textbf{\guru{सू}\laghu{क्ष्म}} \par (3)} &
\multicolumn{3}{>{\centering\arraybackslash\hsize=3.0\hsize}X|}{\textbf{\guru{रू}\laghu{प}} \par (3)} &
\multicolumn{2}{>{\centering\arraybackslash\hsize=2.0\hsize}X|}{\textbf{\laghu{ध}\laghu{रि}} \par (2)} &
\multicolumn{3}{>{\centering\arraybackslash\hsize=3.0\hsize}X|}{\textbf{\laghu{सि}\laghu{य}\laghu{हिँ}} \par (3)} &
\multicolumn{5}{>{\centering\arraybackslash\hsize=5.0\hsize}X|}{\textbf{\laghu{दि}\guru{खा}\guru{वा}} \par (5)} \\
\hline
\end{tabularx}
\vspace{-1.5em}
\end{table}

\begin{table}[h!]
\centering
\renewcommand{\arraystretch}{1.2}
\setlength{\tabcolsep}{0pt}
\begin{tabularx}{\textwidth}{|*{16}{>{\centering\arraybackslash\hsize=1.0\hsize}X|}}
\hline
\multicolumn{3}{|>{\centering\arraybackslash\hsize=3.0\hsize}X|}{\textbf{\laghu{बि}\laghu{क}\laghu{ट}} \par (3)} &
\multicolumn{3}{>{\centering\arraybackslash\hsize=3.0\hsize}X|}{\textbf{\guru{रू}\laghu{प}} \par (3)} &
\multicolumn{2}{>{\centering\arraybackslash\hsize=2.0\hsize}X|}{\textbf{\laghu{ध}\laghu{रि}} \par (2)} &
\multicolumn{3}{>{\centering\arraybackslash\hsize=3.0\hsize}X|}{\textbf{\guru{लं}\laghu{क}} \par (3)} &
\multicolumn{5}{>{\centering\arraybackslash\hsize=5.0\hsize}X|}{\textbf{\laghu{ज}\guru{रा}\guru{वा}} \par (5)} \\
\hline
\end{tabularx}
\vspace{-1.5em}
\end{table}

\vspace{1em}
\textbf{ \deva{सरल-संस्कृतम्}:}\\
 \deva{भवान् सूक्ष्मं (लघु) रूपं धृत्वा सीतायै (सियहिं) दर्शनं दत्तवान्।}\\
 \deva{भवान् विकटं (भयानकं/विशालं) रूपं धृत्वा लङ्कां भस्मीकृतवान् (जरावा)।}\\


Assuming a microscopic form, you revealed yourself to Mother Sita.\\
Assuming a massive/terrifying form, you burned down Lanka.


\vspace{1em}
\newpage
\textbf{\deva{प्रतिपदार्थः} (Meaning of each word):}
\begin{table}[h!]
\centering
\renewcommand{\arraystretch}{1.2}
\begin{tabularx}{\textwidth}{@{} l l X X @{}}
\toprule
\textbf{Awadhi} & \textbf{Sanskrit} & \textbf{English} & \textbf{Grammar \& Notes} \\
\midrule
\deva{सूक्ष्म रूप} / \telu{సూక్ష్మ రూప}  &  \deva{सूक्ष्म-रूपम्}  &  Microscopic form  &   \\
\deva{धरि} / \telu{ధరి}  &  \deva{धृत्वा}  &  Having assumed  &   \\
\deva{सियहिं} \gotchaicon / \telu{సియహిఁ}  &  \deva{सीतायै}  &  To Mother Sita  & Dative case suffix -hiṃ \\
\deva{दिखावा} / \telu{దిఖావా}  &  \deva{दर्शितवान्}  &  Showed / Revealed  &   \\
\deva{बिकट रूप} / \telu{బికట రూప}  &  \deva{विकट-रूपम्}  &  Terrifying form  & \\ 
\deva{धरि} / \telu{ధరి}  &  \deva{धृत्वा}  &  Having assumed  &   \\
\deva{लंक} / \telu{లంక}  &  \deva{लङ्काम्}  &  Lanka  &   \\
\deva{जरावा} / \telu{జరావా}  &  \deva{ज्वालितवान्}  &  Burned  & Jval $\rightarrow$ Jaraavaa \\ 
\bottomrule
\end{tabularx}
\end{table}

\begin{tcolorbox}[gotchabox]
\begin{itemize}
    \item \textbf{The Dative Suffix (\deva{सियहिं}):} In Awadhi, the formal name \deva{सीता} softens to \deva{सिय} (Siya). The \deva{-हिं} (-hiṃ) is the objective marker we saw earlier in \deva{मोहिं} (to me). So \deva{सियहिं} means "To Mother Siya".
    \item \textbf{Phonetic Shifts:} Sanskrit \deva{विकट} (Vikata) becomes \deva{बिकट} (Bikata) via the $v \rightarrow b$ shift. Also, the causative verb \deva{जरावा} (Jaraavaa) comes from the root \deva{ज्वल्} (to burn), meaning "caused to burn".
\end{itemize}
\end{tcolorbox}



\newpage


% ----- CHAUPAI 10 -----

\newpage

\subsection*{\deva{दशम चौपाई} (Tenth Chaupai)}
\addcontentsline{toc}{subsection}{\deva{दशम चौपाई} (Tenth Chaupai) — \deva{भीम रूप...}}

\begin{minipage}[t]{0.55\textwidth}
\vspace{0pt}

{\large \linespread{1.5}\selectfont
\deva{भीम रूप धरि असुर सँहारे।}\\
\deva{रामचंद्र के काज सँवारे॥}
\par}

\vspace{0.5em}

{\large \linespread{1.5}\selectfont
\telu{భీమ రూప ధరి అసుర సఁహారే।}\\
\telu{రామచంద్ర కే కాజ సఁవారే॥}
\par}

\vspace{0.5em}

{\large \linespread{1.3}\selectfont
\textit{bhīma rūpa dhari asura sãhāre |}\\
\textit{rāmacandra ke kāja sãvāre ||}
\par}
\end{minipage}%
\hfill
\begin{minipage}[t]{0.42\textwidth}
\vspace{0pt}
\begin{tcolorbox}[anchorbox]
\textbf{Standing Up to Bullies:} Hanuman was inherently peaceful, but when his boundaries were repeatedly broken by toxic forces (\textit{Asura}), he assumed a terrifying, massive form (\textit{Bhima Rupa}) to protect his mission. True strength means being gentle by default, but having the courage and ferocity to stand up decisively to bullies and injustice when necessary.
\vspace{0.5em}\newline When backed into a corner by toxic forces, he unleashed an unapologetic, fierce defense to protect his mission.\newline\null\hfill\href{https://www.valmiki.iitk.ac.in/sloka?field_kanda_tid=5&language=dv&field_sarga_value=47}{\textit{Sundara Kanda, 47:34-37}}
\end{tcolorbox}
\vspace{0.5em}

\end{minipage}

\vspace{0.5em}
\subsubsection*{\textbf{Visual Rhythm Map:}}
\begin{table}[h!]
\centering
\renewcommand{\arraystretch}{1.2}
\noindent\makebox[\textwidth][l]{%
  {%
  \setlength{\tabcolsep}{0pt}%
  \begin{tabularx}{1.0625\textwidth}{|*{17}{>{\centering\arraybackslash\hsize=1.0\hsize}X|}}
  \hline
  \multicolumn{3}{|>{\centering\arraybackslash\hsize=3.0\hsize}X|}{\textbf{\guru{भी}\laghu{म}} \par (3)} &
  \multicolumn{3}{>{\centering\arraybackslash\hsize=3.0\hsize}X|}{\textbf{\guru{रू}\laghu{प}} \par (3)} &
  \multicolumn{2}{>{\centering\arraybackslash\hsize=2.0\hsize}X|}{\textbf{\laghu{ध}\laghu{रि}} \par (2)} &
  \multicolumn{3}{>{\centering\arraybackslash\hsize=3.0\hsize}X|}{\textbf{\laghu{अ}\laghu{सु}\laghu{र}} \par (3)} &
  \multicolumn{6}{>{\centering\arraybackslash\hsize=5.0\hsize}X|}{\textbf{\laghu{सँ}\guru{हा}\guru{रे}} \par (5)} \\
  \hline
  \noalign{\vspace{6pt}}
  \hline
  \multicolumn{6}{|>{\centering\arraybackslash\hsize=6.0\hsize}X|}{\textbf{\guru{रा}\laghu{म}\guru{चं}\laghu{द्र}} \par (6)} &
  \multicolumn{2}{>{\centering\arraybackslash\hsize=2.0\hsize}X|}{\textbf{\guru{के}} \par (2)} &
  \multicolumn{3}{>{\centering\arraybackslash\hsize=3.0\hsize}X|}{\textbf{\guru{का}\laghu{ज}} \par (3)} &
  \multicolumn{6}{>{\centering\arraybackslash\hsize=5.0\hsize}X|}{\textbf{\laghu{सँ}\guru{वा}\guru{रे}} \par (5)} \\
  \hline
  \end{tabularx}%
  }%
}
\vspace{-1.5em}
\end{table}

\vspace{1em}
\textbf{ \deva{सरल-संस्कृतम्}:}\\
 \deva{भवान् भयंकरं (भीमं) रूपं धृत्वा असुरान् संहारितवान्।}\\
 \deva{श्रीरामचन्द्रस्य सर्वाणि कार्याणि भवान् साधितवान् (सफलानि कृतवान्)।}\\


Assuming a massive/terrifying form, you destroyed the demons.\\
And thus, you successfully accomplished the missions of Lord Ramachandra.


\vspace{1em}
\newpage
\textbf{\deva{प्रतिपदार्थः} (Meaning of each word):}
\begin{table}[h!]
\centering
\renewcommand{\arraystretch}{1.2}
\begin{tabularx}{\textwidth}{@{} l l X X @{}}
\toprule
\textbf{Awadhi} & \textbf{Sanskrit} & \textbf{English} & \textbf{Grammar \& Notes} \\
\midrule
\deva{भीम रूप} / \telu{భీమ రూప} / \textit{bhīma rūpa}  &  \deva{भीम-रूपम्}  &  Massive / fearsome form  &   \\
\deva{धरि} / \telu{ధరి} / \textit{dhari}  &  \deva{धृत्वा}  &  Having assumed  &   \\
\deva{असुर} / \telu{అసుర} / \textit{asura}  &  \deva{असुरान्}  &  Demons  &   \\
\deva{सँहारे} / \telu{సఁహారే} / \textit{sãhāre}  &  \deva{संहारितवान्}  &  Destroyed  &   \\
\deva{रामचंद्र के} / \telu{రామచంద్ర కే} / \textit{rāmacandra ke}  &  \deva{श्रीरामचन्द्रस्य}  &  Ramachandra's  & Genitive case marker `ke` \\
\deva{काज} / \telu{కాజ} / \textit{kāja}  &  \deva{कार्याणि}  &  Works/Missions  & Prakrit shift () \\ 
\deva{सँवारे} / \telu{సఁవారే} / \textit{sãvāre}  &  \deva{साधितवान्}  &  Accomplished  &   \\
\bottomrule
\end{tabularx}
\end{table}

\begin{tcolorbox}[poetrybox]
\textbf{The 17-Beat Anomaly \& Bhava-Pradhana:}\\
Mathematically, both lines of this verse sum up to 17 beats due to the heavy weights of words like \deva{संहारे} and \deva{संवारे}. 
But Tulsi das, avoids that problem by using  \deva{सँहारे} and  \deva{सँवारे} ! 
The Hanuman Chalisa is a \textit{\deva{भाव प्रधान}} (emotion foremost) poem.
Singers naturally absorb this 17th beat by slightly eliding or compressing syllables where needed. 
These metrical inconsistencies actually serve as beautiful mnemonic devices that add to the beauty of the singing experience!\\
\textbf{Rhyming Finale:}\\
Notice the brilliant rhyme: by killing the demons (\deva{सँहारे}), he makes Rama's mission a success (\deva{सँवारे}).
\end{tcolorbox}

\newpage



% ==========================================
\section*{Group 4: Praises from Rama and the Heavens}
\addcontentsline{toc}{section}{Group 4: Praises from Rama and the Heavens}
% ==========================================

% ----- CHAUPAI 11 -----
\subsection*{\deva{एकादश चौपाई} (Eleventh Chaupai)}
\addcontentsline{toc}{subsection}{\deva{एकादश चौपाई} (Eleventh Chaupai) — \deva{लाय सजीवन...}}

\begin{minipage}[t]{0.55\textwidth}
\vspace{0pt}

{\large \linespread{1.5}\selectfont
\deva{लाय सजीवन लखन जियाये।}\\
\deva{श्रीरघुबीर हरषि उर लाये॥}
\par}

\vspace{0.5em}

{\large \linespread{1.5}\selectfont
\telu{లాయ సజీవన లఖన జియాయే।}\\
\telu{శ్రీరఘుబీర హరషి ఉర లాయే॥}
\par}

\vspace{0.5em}

{\large \linespread{1.3}\selectfont
\textit{lāya sajīvana lakhana jiyāye |}\\
\textit{śrīraghubīra haraṣi ura lāye ||}
\par}
\end{minipage}%
\hfill
\begin{minipage}[t]{0.42\textwidth}
\vspace{0pt}
\begin{tcolorbox}[anchorbox]
\textbf{The Ultimate Team Player:} When his teammate fell in battle, Hanuman didn't hesitate; he crossed continents overnight and worked tirelessly to bring the life-saving cure (\textit{Sanjeevani}). He teaches us the absolute importance of resourcefulness, reliability, and showing up unconditionally for our friends in their darkest hours.
\vspace{0.5em}\newline Hanuman worked tirelessly through the night to save his fallen commander, defining the essence of absolute reliability.\newline\null\hfill\href{https://www.valmiki.iitk.ac.in/sloka?field_kanda_tid=6&language=dv&field_sarga_value=74}{\textit{Yuddha Kanda, 74:71-74}}
\end{tcolorbox}
\end{minipage}

\vspace{0.5em}
\subsubsection*{\textbf{Visual Rhythm Map:}}
\begin{table}[h!]
\centering
\renewcommand{\arraystretch}{1.2}
\setlength{\tabcolsep}{0pt}
\begin{tabularx}{\textwidth}{|*{16}{>{\centering\arraybackslash\hsize=1.0\hsize}X|}}
\hline
\multicolumn{3}{|>{\centering\arraybackslash\hsize=3.0\hsize}X|}{\textbf{\guru{ला}\laghu{य}} \par (3)} &
\multicolumn{5}{>{\centering\arraybackslash\hsize=5.0\hsize}X|}{\textbf{\laghu{स}\guru{जी}\laghu{व}\laghu{न}} \par (5)} &
\multicolumn{3}{>{\centering\arraybackslash\hsize=3.0\hsize}X|}{\textbf{\laghu{ल}\laghu{ख}\laghu{न}} \par (3)} &
\multicolumn{5}{>{\centering\arraybackslash\hsize=5.0\hsize}X|}{\textbf{\laghu{जि}\guru{या}\guru{ये}} \par (5)} \\
\hline
\end{tabularx}
\vspace{-1.5em}
\end{table}

\begin{table}[h!]
\centering
\renewcommand{\arraystretch}{1.2}
\setlength{\tabcolsep}{0pt}
\begin{tabularx}{\textwidth}{|*{16}{>{\centering\arraybackslash\hsize=1.0\hsize}X|}}
\hline
\multicolumn{7}{|>{\centering\arraybackslash\hsize=7.0\hsize}X|}{\textbf{\guru{श्री}\laghu{र}\laghu{घु}\guru{बी}\laghu{र}} \par (7)} &
\multicolumn{3}{>{\centering\arraybackslash\hsize=3.0\hsize}X|}{\textbf{\laghu{ह}\laghu{र}\laghu{षि}} \par (3)} &
\multicolumn{2}{>{\centering\arraybackslash\hsize=2.0\hsize}X|}{\textbf{\laghu{उ}\laghu{र}} \par (2)} &
\multicolumn{4}{>{\centering\arraybackslash\hsize=4.0\hsize}X|}{\textbf{\guru{ला}\guru{ये}} \par (4)} \\
\hline
\end{tabularx}
\vspace{-1.5em}
\end{table}

\vspace{1em}
\textbf{ \deva{सरल-संस्कृतम्}:}\\
 \deva{भवान् सञ्जीवनी-बूटीम् आनीय लक्ष्मणस्य प्राणान् रक्षितवान् (जियाये)।}\\
 \deva{तत् दृष्ट्वा श्रीरामः (रघुवीरः) हर्षितः सन् भवन्तं वक्षसि (उर) स्थापितवान् (आलिङ्गनं कृतवान्)॥}\\


Bringing the Sanjeevani herb, you resurrected Lakshmana.\\
Seeing this, Lord Rama joyfully embraced you to his chest/heart.


\vspace{1em}
\newpage
\textbf{\deva{प्रतिपदार्थः} (Meaning of each word):}
\begin{table}[h!]
\centering
\renewcommand{\arraystretch}{1.2}
\begin{tabularx}{\textwidth}{@{} l l X X @{}}
\toprule
\textbf{Awadhi} & \textbf{Sanskrit} & \textbf{English} & \textbf{Grammar \& Notes} \\
\midrule
\deva{लाय} / \telu{లాయ} / \textit{lāya}  &  \deva{आनीय}  &  Having brought  &   \\
\deva{सजीवन} \gotchaicon / \telu{సజీవన} / \textit{sajīvana}  &  \deva{सञ्जीवनीम्}  &  Sanjeevani herb  &   \\
\deva{लखन} / \telu{లఖన} / \textit{lakhana}  &  \deva{लक्ष्मणम्}  &  Lakshmana  & \\ 
\deva{जियाये} / \telu{జియాయే} / \textit{jiyāye}  &  \deva{जीवितवान्}  &  Resurrected / Revived  &   \\
\deva{श्रीरघुबीर} / \telu{శ్రీరఘుబీర} / \textit{śrīraghubīra}  &  \deva{श्रीरघुवीरः}  &  Lord Rama  & \\ 
\deva{हरषि} / \telu{హరషి} / \textit{haraṣi}  &  \deva{हर्षितः}  &  Joyfully  & Short ending 'i' \\
\deva{उर} / \telu{ఉర} / \textit{ura}  &  \deva{उरसि / वक्षसि}  &  In the chest/heart  &   \\
\deva{लाये} / \telu{లాయే} / \textit{lāye}  &  \deva{स्थापितवान्}  &  Brought to (embraced)  &   \\
\bottomrule
\end{tabularx}
\end{table}

\begin{tcolorbox}[gotchabox]
\begin{itemize}
    \item \textbf{Phonetic Shifts:} \deva{रघुवीर} (Raghuveera) hardens into \deva{रघुबीर} (Raghubeera). The Sanskrit `jñ/ñj` cluster in \deva{सञ्जीवनी} (Sanjeevani) simplifies into the pure Awadhi \deva{सजीवन} (Sajeevan).
\end{itemize}
\end{tcolorbox}



\newpage


% ----- CHAUPAI 12 -----

\newpage

\subsection*{\deva{द्वादश चौपाई} (Twelfth Chaupai)}
\addcontentsline{toc}{subsection}{\deva{द्वादश चौपाई} (Twelfth Chaupai) — \deva{रघुपति कीन्ही...}}

\begin{minipage}[t]{0.55\textwidth}
\vspace{0pt}

{\large \linespread{1.5}\selectfont
\deva{रघुपति कीन्ही बहुत बड़ाई।}\\
\deva{तुम मम प्रिय भरतहि सम भाई॥}
\par}

\vspace{0.5em}

{\large \linespread{1.5}\selectfont
\telu{రఘుపతి కీన్హీ బహుత బడాయీ।}\\
\telu{తుమ మమ ప్రియ భరతహి సమ భాయీ॥}
\par}

\vspace{0.5em}

{\large \linespread{1.3}\selectfont
\textit{raghupati kīnhī bahuta baṛāī |}\\
\textit{tuma mama priya bharatahi sama bhāī ||}
\par}
\end{minipage}%
\hfill
\begin{minipage}[t]{0.42\textwidth}
\vspace{0pt}
\begin{tcolorbox}[anchorbox]
\textbf{Recognition of Merit:} Lord Rama praised Hanuman by stating he was equal to his own royal brother Bharata. A true leader creates a meritocracy. They judge people by their character, loyalty, and hard work, not by their family background, species, or social status.
\vspace{0.5em}\newline The supreme leader publicly embraced and rewarded him, proving that merit and hard work transcend all social titles.\newline\null\hfill\href{https://www.valmiki.iitk.ac.in/sloka?field_kanda_tid=6&language=dv&field_sarga_value=1}{\textit{Yuddha Kanda, 1:12-14}}
\end{tcolorbox}
\end{minipage}

\vspace{0.5em}
\subsubsection*{\textbf{Visual Rhythm Map:}}
\begin{table}[h!]
\centering
\renewcommand{\arraystretch}{1.2}
\setlength{\tabcolsep}{0pt}
\begin{tabularx}{\textwidth}{|*{16}{>{\centering\arraybackslash\hsize=1.0\hsize}X|}}
\hline
\multicolumn{4}{|>{\centering\arraybackslash\hsize=4.0\hsize}X|}{\textbf{\laghu{र}\laghu{घु}\laghu{प}\laghu{ति}} \par (4)} &
\multicolumn{4}{>{\centering\arraybackslash\hsize=4.0\hsize}X|}{\textbf{\guru{कीं}\guru{न्ही}} \par (4)} &
\multicolumn{3}{>{\centering\arraybackslash\hsize=3.0\hsize}X|}{\textbf{\laghu{ब}\laghu{हु}\laghu{त}} \par (3)} &
\multicolumn{5}{>{\centering\arraybackslash\hsize=5.0\hsize}X|}{\textbf{\laghu{ब}\guru{ड़ा}\guru{ई}} \par (5)} \\
\hline
\end{tabularx}
\vspace{-1.5em}
\end{table}

\begin{table}[h!]
\centering
\renewcommand{\arraystretch}{1.2}
\setlength{\tabcolsep}{0pt}
\begin{tabularx}{\textwidth}{|*{16}{>{\centering\arraybackslash\hsize=1.0\hsize}X|}}
\hline
\multicolumn{2}{|>{\centering\arraybackslash\hsize=2.0\hsize}X|}{\textbf{\laghu{तु}\laghu{म}} \par (2)} &
\multicolumn{2}{>{\centering\arraybackslash\hsize=2.0\hsize}X|}{\textbf{\laghu{म}\laghu{म}} \par (2)} &
\multicolumn{2}{>{\centering\arraybackslash\hsize=2.0\hsize}X|}{\textbf{\laghu{प्रि}\laghu{य}} \par (2)} &
\multicolumn{4}{>{\centering\arraybackslash\hsize=4.0\hsize}X|}{\textbf{\laghu{भ}\laghu{र}\laghu{त}\laghu{हि}} \par (4)} &
\multicolumn{2}{>{\centering\arraybackslash\hsize=2.0\hsize}X|}{\textbf{\laghu{स}\laghu{म}} \par (2)} &
\multicolumn{4}{>{\centering\arraybackslash\hsize=4.0\hsize}X|}{\textbf{\guru{भा}\guru{ई}} \par (4)} \\
\hline
\end{tabularx}
\vspace{-1.5em}
\end{table}

\vspace{1em}
\textbf{ \deva{सरल-संस्कृतम्}:}\\
 \deva{श्रीरामः (रघुपतिः) भवतः अतीव प्रशंसां (बड़ाई) कृतवान्।}\\
 \deva{"त्वं मम (भरतेन समानः) प्रिय-भ्राता असि"॥}\\

The Lord of the Raghus praised you immensely.
He said, "You are as dear a brother to me as Bharata."


\vspace{1em}
\newpage
\textbf{\deva{प्रतिपदार्थः} (Meaning of each word):}
\begin{table}[h!]
\centering
\renewcommand{\arraystretch}{1.2}
\begin{tabularx}{\textwidth}{@{} l l X X @{}}
\toprule
\textbf{Awadhi} & \textbf{Sanskrit} & \textbf{English} & \textbf{Grammar \& Notes} \\
\midrule
\deva{रघुपति} / \telu{రఘుపతి}  &  \deva{रघुपतिः}  &  Lord Rama  &   \\
\deva{कीन्ही} / \telu{కీన్హీ}  &  \deva{कृतवान्}  &  Did  & Awadhi past tense verb \\
\deva{बहुत} / \telu{బహుత}  &  \deva{अतीव / बहु}  &  Much / Lots of  &   \\
\deva{बड़ाई} / \telu{బడాయీ}  &  \deva{प्रशंसा / महिमा}  &  Praise / Greatness  &   \\
\deva{तुम} / \telu{తుమ}  &  \deva{त्वम्}  &  You  &   \\
\deva{मम} / \telu{మమ}  &  \deva{मम}  &  My  & Pure Sanskrit! \\
\deva{प्रिय} / \telu{ప్రియ}  &  \deva{प्रियः}  &  Dear  & Pure Sanskrit! \\
\deva{भरतहि} / \telu{భరతహి}  &  \deva{भरताय / भरतेन}  &  Bharata  & Objective suffix `-hi` \\
\deva{सम} / \telu{సమ}  &  \deva{समानः}  &  Equal / Like  &   \\
\deva{भाई} / \telu{భాయీ}  &  \deva{भ्राता}  &  Brother  &   \\
\bottomrule
\end{tabularx}
\end{table}



\newpage


% ----- CHAUPAI 13 -----

\newpage

\subsection*{\deva{त्रयोदश चौपाई} (Thirteenth Chaupai)}
\addcontentsline{toc}{subsection}{\deva{त्रयोदश चौपाई} (Thirteenth Chaupai) — \deva{सहस बदन...}}

\begin{minipage}[t]{0.55\textwidth}
\vspace{0pt}

{\large \linespread{1.5}\selectfont
\deva{सहस बदन तुम्हरो जस गावैं।}\\
\deva{अस कहि श्रीपति कंठ लगावैं॥}
\par}

\vspace{0.5em}

{\large \linespread{1.5}\selectfont
\telu{సహస బదన తుమ్హరో జస గావైఁ।}\\
\telu{అస కహి శ్రీపతి కంఠ లగావైఁ॥}
\par}

\vspace{0.5em}

{\large \linespread{1.3}\selectfont
\textit{sahasa badana tumharo jasa gāvaĩ |}\\
\textit{asa kahi śrīpati kaṇṭha lagāvaĩ ||}
\par}
\end{minipage}%
\hfill
\begin{minipage}[t]{0.42\textwidth}
\vspace{0pt}
\begin{tcolorbox}[anchorbox]
\textbf{Building a Legacy:} Lord Rama tells him that the entire universe will sing his praises (\textit{Jasa Gaavai}). When you dedicate your life to selfless service, hard work, and excellence, you never need to brag or demand attention. Your reputation and your results will eventually sing for themselves.
\vspace{0.5em}\newline His monumental achievements were so undeniably excellent that his reputation sang for itself globally.\newline\null\hfill\href{https://www.valmiki.iitk.ac.in/sloka?field_kanda_tid=6&language=dv&field_sarga_value=1}{\textit{Yuddha Kanda, 1:11}}
\end{tcolorbox}
\end{minipage}

\vspace{0.5em}
\subsubsection*{\textbf{Visual Rhythm Map:}}
\begin{table}[h!]
\centering
\renewcommand{\arraystretch}{1.2}
\noindent\makebox[\textwidth][l]{%
  {%
  \setlength{\tabcolsep}{0pt}%
  \begin{tabularx}{1.0625\textwidth}{|*{17}{>{\centering\arraybackslash\hsize=1.0\hsize}X|}}
  \hline
  \multicolumn{3}{|>{\centering\arraybackslash\hsize=3.0\hsize}X|}{\textbf{\laghu{स}\laghu{ह}\laghu{स}} \par (3)} &
  \multicolumn{3}{>{\centering\arraybackslash\hsize=3.0\hsize}X|}{\textbf{\laghu{ब}\laghu{द}\laghu{न}} \par (3)} &
  \multicolumn{5}{>{\centering\arraybackslash\hsize=5.0\hsize}X|}{\textbf{\guru{तु}\laghu{म्ह}\guru{रो}} \par (5)} &
  \multicolumn{2}{>{\centering\arraybackslash\hsize=2.0\hsize}X|}{\textbf{\laghu{ज}\laghu{स}} \par (2)} &
  \multicolumn{4}{>{\centering\arraybackslash\hsize=4.0\hsize}X|}{\textbf{\guru{गा}\guru{वैं}} \par (4)} \\
  \hline
  \end{tabularx}%
  }%
  \hspace{0.01\textwidth}%
  \begin{minipage}{0.1\textwidth}
  \vspace{-0.5em}
  \centering
  {\Huge \textbf{17}}
  \end{minipage}%
}
\vspace{-1.5em}
\end{table}

\begin{table}[h!]
\centering
\renewcommand{\arraystretch}{1.2}
\setlength{\tabcolsep}{0pt}
\begin{tabularx}{\textwidth}{|*{16}{>{\centering\arraybackslash\hsize=1.0\hsize}X|}}
\hline
\multicolumn{2}{|>{\centering\arraybackslash\hsize=2.0\hsize}X|}{\textbf{\laghu{अ}\laghu{स}} \par (2)} &
\multicolumn{2}{>{\centering\arraybackslash\hsize=2.0\hsize}X|}{\textbf{\laghu{क}\laghu{हि}} \par (2)} &
\multicolumn{4}{>{\centering\arraybackslash\hsize=4.0\hsize}X|}{\textbf{\guru{श्री}\laghu{प}\laghu{ति}} \par (4)} &
\multicolumn{3}{>{\centering\arraybackslash\hsize=3.0\hsize}X|}{\textbf{\guru{कं}\laghu{ठ}} \par (3)} &
\multicolumn{5}{>{\centering\arraybackslash\hsize=5.0\hsize}X|}{\textbf{\laghu{ल}\guru{गा}\guru{वैं}} \par (5)} \\
\hline
\end{tabularx}
\vspace{-1.5em}
\end{table}

\vspace{1em}
\textbf{ \deva{सरल-संस्कृतम्}:}\\
 \deva{"सहस्रं (1000) मुखानि (शेषानागः) भवतः यशः गायन्ति।"}\\
 \deva{एवं उक्त्वा श्रियः-पतिः (श्रीरामः) भवन्तं कण्ठे आलिङ्गितवान्।}\\


"The thousand-faced serpent (Shesha Naga) sings your glory"\\
Thus saying, Sri Rama embraced you to his neck.


\vspace{1em}
\newpage
\textbf{\deva{प्रतिपदार्थः} (Meaning of each word):}
\begin{table}[h!]
\centering
\renewcommand{\arraystretch}{1.2}
\begin{tabularx}{\textwidth}{@{} l l X X @{}}
\toprule
\textbf{Awadhi} & \textbf{Sanskrit} & \textbf{English} & \textbf{Grammar \& Notes} \\
\midrule
\deva{सहस} / \telu{సహస}  &  \deva{सहस्रम्}  &  Thousand  &   \\
\deva{बदन} / \telu{బదన}  &  \deva{वदनम् / मुखानि}  &  Faces  & \\ 
\deva{तुम्हरो} / \telu{తుమ్హరో}  &  \deva{भवतः}  &  Your  &   \\
\deva{जस} \gotchaicon / \telu{జస}  &  \deva{यशः}  &  Glory / Fame  & \\ 
\deva{गावैं} \gotchaicon / \telu{గావైఁ}  &  \deva{गायन्ति}  &  (They) Sing  & Plural nasal `-ain` \\
\deva{अस} / \telu{అస}  &  \deva{एवम्}  &  Like this / Thus  &   \\
\deva{कहि} / \telu{కహి}  &  \deva{उक्त्वा}  &  Having said  & Absolutive verb `-i` \\
\deva{श्रीपति} / \telu{శ్రీపతి}  &  \deva{श्रियः-पतिः}  &  Lord of Fortune (Rama)  &   \\
\deva{कंठ} / \telu{కంఠ}  &  \deva{कण्ठे}  &  By the neck  &   \\
\deva{लगावैं} \gotchaicon / \telu{లగావైఁ}  &  \deva{आलिङ्गितवान्}  &  Hugged / Applied  &   \\
\bottomrule
\end{tabularx}
\end{table}

\begin{tcolorbox}[gotchabox]
\begin{itemize}
    \item \textbf{The Double Shift (\deva{यश $\rightarrow$ जस}):} Pure Awadhi transformation! The gliding \deva{य} (y) hardens to the plosive \deva{ज} (j), and the palatal \deva{श} (sh) softens to the dental \deva{स} (s). Sanskrit \deva{यश} (Yasha) becomes Awadhi \deva{जस} (Jasa).
    \item \textbf{The Plural Verb Nasal (\deva{गावैं/लगावैं}):} The nasalization \deva{ँ} at the end of verbs denotes plural or respect. Here, it denotes that the 1000 faces of the serpent \textit{sing} (Gaavain), and that Lord Rama (with immense respect) \textit{embraces} (Lagaavain).
\end{itemize}
\end{tcolorbox}



\newpage


% ----- CHAUPAI 14 -----

\newpage

\subsection*{\deva{चतुर्दश चौपाई} (Fourteenth Chaupai)}
\addcontentsline{toc}{subsection}{\deva{चतुर्दश चौपाई} (Fourteenth Chaupai) — \deva{सनकादिक ब्रह्मादि...}}

\begin{minipage}[t]{0.55\textwidth}
\vspace{0pt}

{\large \linespread{1.5}\selectfont
\deva{सनकादिक ब्रह्मादि मुनीसा।}\\
\deva{नारद सारद सहित अहीसा॥}
\par}

\vspace{0.5em}

{\large \linespread{1.5}\selectfont
\telu{సనకాదిక బ్రహ్మాది మునీసా।}\\
\telu{నారద సారద సహిత అహీసా॥}
\par}

\vspace{0.5em}

{\large \linespread{1.3}\selectfont
\textit{sanakādika brahmādi munīsā |}\\
\textit{nārada sārada sahita ahīsā ||}
\par}
\end{minipage}%
\hfill
\begin{minipage}[t]{0.42\textwidth}
\vspace{0pt}
\end{minipage}

\vspace{0.5em}
\subsubsection*{\textbf{Visual Rhythm Map:}}
\begin{table}[h!]
\centering
\renewcommand{\arraystretch}{1.2}
\setlength{\tabcolsep}{0pt}
\begin{tabularx}{\textwidth}{|*{16}{>{\centering\arraybackslash\hsize=1.0\hsize}X|}}
\hline
\multicolumn{6}{|>{\centering\arraybackslash\hsize=6.0\hsize}X|}{\textbf{\laghu{स}\laghu{न}\guru{का}\laghu{दि}\laghu{क}} \par (6)} &
\multicolumn{5}{>{\centering\arraybackslash\hsize=5.0\hsize}X|}{\textbf{\guru{ब्र}\guru{ह्मा}\laghu{दि}} \par (5)} &
\multicolumn{5}{>{\centering\arraybackslash\hsize=5.0\hsize}X|}{\textbf{\laghu{मु}\guru{नी}\guru{सा}} \par (5)} \\
\hline
\end{tabularx}
\vspace{-1.5em}
\end{table}

\begin{table}[h!]
\centering
\renewcommand{\arraystretch}{1.2}
\setlength{\tabcolsep}{0pt}
\begin{tabularx}{\textwidth}{|*{16}{>{\centering\arraybackslash\hsize=1.0\hsize}X|}}
\hline
\multicolumn{4}{|>{\centering\arraybackslash\hsize=4.0\hsize}X|}{\textbf{\guru{ना}\laghu{र}\laghu{द}} \par (4)} &
\multicolumn{4}{>{\centering\arraybackslash\hsize=4.0\hsize}X|}{\textbf{\guru{सा}\laghu{र}\laghu{द}} \par (4)} &
\multicolumn{3}{>{\centering\arraybackslash\hsize=3.0\hsize}X|}{\textbf{\laghu{स}\laghu{हि}\laghu{त}} \par (3)} &
\multicolumn{5}{>{\centering\arraybackslash\hsize=5.0\hsize}X|}{\textbf{\laghu{अ}\guru{ही}\guru{सा}} \par (5)} \\
\hline
\end{tabularx}
\vspace{-1.5em}
\end{table}

\vspace{1em}
\textbf{ \deva{सरल-संस्कृतम्}:}\\
 \deva{सनकादयः (सनक-सनन्दन-सनातन-सनत्कुमाराः), ब्रह्मा-आदयः देवाः, मुनीशाः च}\\
 \deva{नारदः, देवी सरस्वती (शारदा), शेषनाग-सहिताः (अहीशः) अपि यस्य यशः सम्यक् वर्णयितुं न शक्नुवन्ति।}\\


Sages like Sanaka, gods like Brahma, and great ascetics\\
even Narada, Saraswati, and the Lord of Serpents.


\vspace{1em}
\newpage
\textbf{\deva{प्रतिपदार्थः} (Meaning of each word):}
\begin{table}[h!]
\centering
\renewcommand{\arraystretch}{1.2}
\begin{tabularx}{\textwidth}{@{} l l X X @{}}
\toprule
\textbf{Awadhi} & \textbf{Sanskrit} & \textbf{English} & \textbf{Grammar \& Notes} \\
\midrule
\deva{सनकादिक} / \telu{సనకాదిక}  &  \deva{सनकादयः}  &  Sages like Sanaka  &   \\
\deva{ब्रह्मादि} / \telu{బ్రహ్మాది}  &  \deva{ब्रह्मादयः देवाः}  &  Gods like Brahma  &   \\
\deva{मुनीसा} \gotchaicon / \telu{మునీసా}  &  \deva{मुनीशाः (मुनि+ईश)}  &  Chief sages  & \\ 
\deva{नारद} / \telu{నారద}  &  \deva{नारदः}  &  Sage Narada  &   \\
\deva{सारद} \gotchaicon / \telu{సారద}  &  \deva{शारदा / सरस्वती}  &  Goddess Saraswati  & \\ 
\deva{सहित} / \telu{సహిత}  &  \deva{सहिताः / सह}  &  Along with  &   \\
\deva{अहीसा} \gotchaicon / \telu{అహీసా}  &  \deva{अहीशः (अहि+ईश)}  &  Lord of Snakes (Shesha)  & \\ 
\bottomrule
\end{tabularx}
\end{table}

\begin{tcolorbox}[gotchabox]
\begin{itemize}
    \item \textbf{The \deva{श/ष} $\rightarrow$ \deva{स} Shift in Action:} Beautiful symmetry here! Sanskrit \deva{मुनीश} (Muneesha), \deva{शारदा} (Shaarada), and \deva{अहीश} (Aheesha) drop their heavy palatal \deva{श} (sh) and soften into the pure Awadhi rhymes: \deva{मुनीसा} (Muneesaa), \deva{सारद} (Saarad), and \deva{अहीसा} (Aheesaa).
    \item \textbf{The 'Aadi' (\deva{आदि}) suffix:} The suffix \deva{आदि} (ādi) is used to denote "and others like them" or "starting from". 
     This chaupai starts a list of greats, but we will know about why they are mentioned here in the next chaupai.
\end{itemize}
\end{tcolorbox}

\newpage


% ----- CHAUPAI 15 -----

\newpage

\subsection*{\deva{पञ्चदश चौपाई} (Fifteenth Chaupai)}
\addcontentsline{toc}{subsection}{\deva{पञ्चदश चौपाई} (Fifteenth Chaupai) — \deva{जमु कुबेर...}}

\begin{minipage}[t]{0.55\textwidth}
\vspace{0pt}

{\large \linespread{1.5}\selectfont
\deva{जमु कुबेर दिगपाल जहाँ ते।}\\
\deva{कबि कोबिद कहि सके कहाँ ते॥}
\par}

\vspace{0.5em}

{\large \linespread{1.5}\selectfont
\telu{జము కుబేర దిగపాల జహాఁ తే।}\\
\telu{కబి కోబిద కహి సకే కహాఁ తే॥}
\par}

\vspace{0.5em}

{\large \linespread{1.3}\selectfont
\textit{jamu kubera digapāla jahāṃ te |}\\
\textit{kabi kobida kahi sake kahāṃ te ||}
\par}
\end{minipage}%
\hfill
\begin{minipage}[t]{0.42\textwidth}
\vspace{0pt}
\begin{tcolorbox}[anchorbox]
\textbf{Remaining Grounded:} Even when surrounded by celestial beings, absolute authorities, and the masters of wealth and death who cannot measure his greatness, Hanuman remains entirely humble. True confidence doesn't require constant external validation, titles, or showing off.
\vspace{0.5em}\newline He carried immense, cosmic capabilities without ever letting his ego demand external validation from poets or scholars.\newline\null\hfill\href{https://www.valmiki.iitk.ac.in/sloka?field_kanda_tid=4&language=dv&field_sarga_value=66}{\textit{Kishkindha Kanda, 66:26-29}}
\end{tcolorbox}
\end{minipage}

\vspace{0.5em}
\subsubsection*{\textbf{Visual Rhythm Map:}}
\begin{table}[h!]
\centering
\renewcommand{\arraystretch}{1.2}
\setlength{\tabcolsep}{0pt}
\begin{tabularx}{\textwidth}{|*{16}{>{\centering\arraybackslash\hsize=1.0\hsize}X|}}
\hline
\multicolumn{2}{|>{\centering\arraybackslash\hsize=2.0\hsize}X|}{\textbf{\laghu{ज}\laghu{मु}} \par (2)} &
\multicolumn{4}{>{\centering\arraybackslash\hsize=4.0\hsize}X|}{\textbf{\laghu{कु}\guru{बे}\laghu{र}} \par (4)} &
\multicolumn{5}{>{\centering\arraybackslash\hsize=5.0\hsize}X|}{\textbf{\laghu{दि}\laghu{ग}\guru{पा}\laghu{ल}} \par (5)} &
\multicolumn{3}{>{\centering\arraybackslash\hsize=3.0\hsize}X|}{\textbf{\laghu{ज}\guru{हाँ}} \par (3)} &
\multicolumn{2}{>{\centering\arraybackslash\hsize=2.0\hsize}X|}{\textbf{\guru{ते}} \par (2)} \\
\hline
\end{tabularx}
\vspace{-1.5em}
\end{table}

\begin{table}[h!]
\centering
\renewcommand{\arraystretch}{1.2}
\setlength{\tabcolsep}{0pt}
\begin{tabularx}{\textwidth}{|*{16}{>{\centering\arraybackslash\hsize=1.0\hsize}X|}}
\hline
\multicolumn{2}{|>{\centering\arraybackslash\hsize=2.0\hsize}X|}{\textbf{\laghu{क}\laghu{बि}} \par (2)} &
\multicolumn{4}{>{\centering\arraybackslash\hsize=4.0\hsize}X|}{\textbf{\guru{को}\laghu{बि}\laghu{द}} \par (4)} &
\multicolumn{2}{>{\centering\arraybackslash\hsize=2.0\hsize}X|}{\textbf{\laghu{क}\laghu{हि}} \par (2)} &
\multicolumn{3}{>{\centering\arraybackslash\hsize=3.0\hsize}X|}{\textbf{\laghu{स}\guru{के}} \par (3)} &
\multicolumn{3}{>{\centering\arraybackslash\hsize=3.0\hsize}X|}{\textbf{\laghu{क}\guru{हाँ}} \par (3)} &
\multicolumn{2}{>{\centering\arraybackslash\hsize=2.0\hsize}X|}{\textbf{\guru{ते}} \par (2)} \\
\hline
\end{tabularx}
\vspace{-1.5em}
\end{table}

\vspace{1em}
\textbf{ \deva{सरल-संस्कृतम्}:}\\
 \deva{यत्र यमराजः, कुबेरः, दिक्पालाः च}\\
 \deva{(तव यशः वर्णयितुं न शक्नुवन्ति), ततः कथं कवयः कोविदाः (पण्डिताः) च तत् वक्तुं शक्नुवन्ति?}\\


When Yama (God of Death), Kubera (God of Wealth), and the Guardians of the Directions\\
fail to describe your glory, then how can mere poets and scholars ever describe it?


\vspace{1em}
\newpage
\textbf{\deva{प्रतिपदार्थः} (Meaning of each word):}
\begin{table}[h!]
\centering
\renewcommand{\arraystretch}{1.2}
\begin{tabularx}{\textwidth}{@{} l l X X @{}}
\toprule
\textbf{Awadhi} & \textbf{Sanskrit} & \textbf{English} & \textbf{Grammar \& Notes} \\
\midrule
\deva{जमु} / \telu{జము}  &  \deva{यमः}  &  God of Death  & \\ 
\deva{कुबेर} / \telu{కుబేర}  &  \deva{कुबेरः}  &  God of Wealth  &   \\
\deva{दिगपाल} / \telu{దిగపాల}  &  \deva{दिक्पालाः}  &  Direction Guardians  & $k \rightarrow g$ shift \\ 
\deva{जहाँ ते} / \telu{జహాఁ తే}  &  \deva{यत्र / यस्मात्}  &  Where from / Even they  &   \\
\deva{कबि} / \telu{కబి}  &  \deva{कविः}  &  Poet  & \\ 
\deva{कोबिद} / \telu{కోబిద}  &  \deva{कोविदः (पण्डितः)}  &  Scholar  & \\ 
\deva{कहि सके} / \telu{కహి సకే}  &  \deva{वक्तुं शक्नुवन्ति}  &  Can say/describe  &   \\
\deva{कहाँ ते} / \telu{కహాఁ తే}  &  \deva{कुतः / कथम्}  &  How / From where  &   \\
\bottomrule
\end{tabularx}
\end{table}

\newpage



% ==========================================
\section*{Group 5: The King-Maker and Protector}
\addcontentsline{toc}{section}{Group 5: The King-Maker and Protector}
% ==========================================
\newpage

% ----- CHAUPAI 16 -----
\subsection*{\deva{षोडश चौपाई} \(Sixteenth Chaupai\)}
\addcontentsline{toc}{subsection}{\deva{षोडश चौपाई} \(Sixteenth Chaupai\) — \deva{तुम उपकार...}}


\begin{minipage}[t]{0.55\textwidth}
\vspace{0pt}
{\large \linespread{1.5}\selectfont
\deva{तुम उपकार सुग्रीवहिं कीन्हा।}\\
\deva{राम मिलाय राज पद दीन्हा॥}
\par}

\vspace{0.5em}

{\large \linespread{1.5}\selectfont
\telu{తుమ ఉపకార సుగ్రీవహిఁ కీన్హా।}\\
\telu{రామ మిలాయ రాజ పద దీన్హా॥}
\par}

\vspace{0.5em}

{\large \linespread{1.3}\selectfont
\textit{tuma upakāra sugrīvahĩ kīnhā |}\\
\textit{rāma milāya rāja pada dīnhā ||}
\par}
\end{minipage}%
\hfill
\begin{minipage}[t]{0.42\textwidth}
\vspace{0pt}
\begin{tcolorbox}[anchorbox]
\textbf{Lifting Others Up:} Sugriva was a displaced outcast who had lost everything. Hanuman didn't just use his skills to advance his own career; he connected Sugriva to Rama and helped restore him to a position of leadership. A truly great leader uses their network and talents to empower and uplift those who are struggling.
\vspace{0.5em}\newline Hanuman leveraged his own networking and diplomatic skills to uplift a struggling, displaced friend into a position of leadership.\newline\null\hfill\href{https://www.valmiki.iitk.ac.in/sloka?field_kanda_tid=4&language=dv&field_sarga_value=5}{\textit{Kishkindha Kanda, 5:15-18}}
\end{tcolorbox}
\end{minipage}


\vspace{0.5em}
\subsubsection*{\textbf{Visual Rhythm Map:}}
\begin{table}[h!]
\centering
\renewcommand{\arraystretch}{1.2}
\setlength{\tabcolsep}{0pt}
\begin{tabularx}{\textwidth}{|*{16}{>{\centering\arraybackslash\hsize=1.0\hsize}X|}}
\hline
\multicolumn{2}{|>{\centering\arraybackslash\hsize=2.0\hsize}X|}{\textbf{\laghu{तु}\laghu{म}} \par (2)} &
\multicolumn{5}{>{\centering\arraybackslash\hsize=5.0\hsize}X|}{\textbf{\laghu{उ}\laghu{प}\guru{का}\laghu{र}} \par (5)} &
\multicolumn{5}{>{\centering\arraybackslash\hsize=5.0\hsize}X|}{\textbf{\laghu{सु}\guru{ग्री}\laghu{व}\laghu{हिं}} \par (5)} &
\multicolumn{4}{>{\centering\arraybackslash\hsize=4.0\hsize}X|}{\textbf{\guru{की}\guru{न्हा}} \par (4)} \\
\hline
\end{tabularx}
\vspace{-1.5em}
\end{table}

\begin{table}[h!]
\centering
\renewcommand{\arraystretch}{1.2}
\setlength{\tabcolsep}{0pt}
\begin{tabularx}{\textwidth}{|*{16}{>{\centering\arraybackslash\hsize=1.0\hsize}X|}}
\hline
\multicolumn{3}{|>{\centering\arraybackslash\hsize=3.0\hsize}X|}{\textbf{\guru{रा}\laghu{म}} \par (3)} &
\multicolumn{4}{>{\centering\arraybackslash\hsize=4.0\hsize}X|}{\textbf{\laghu{मि}\guru{ला}\laghu{य}} \par (4)} &
\multicolumn{3}{>{\centering\arraybackslash\hsize=3.0\hsize}X|}{\textbf{\guru{रा}\laghu{ज}} \par (3)} &
\multicolumn{2}{>{\centering\arraybackslash\hsize=2.0\hsize}X|}{\textbf{\laghu{प}\laghu{द}} \par (2)} &
\multicolumn{4}{>{\centering\arraybackslash\hsize=4.0\hsize}X|}{\textbf{\guru{दी}\guru{न्हा}} \par (4)} \\
\hline
\end{tabularx}
\vspace{-1.5em}
\end{table}

\vspace{0.5em}
\textbf{ \deva{सरल-संस्कृतम्}:}\\
 \deva{भवान् सुग्रीवाय महान्तम् उपकारं कृतवान्।}\\
 \deva{श्रीरामचन्द्रेण सह तस्य मेलनं कारयित्वा, तस्मै राजपदं दत्तवान्।}\\

You did a great favor for Sugriva.\\
By uniting him with Lord Rama, you bestowed the royal throne upon him.


\vspace{1em}
\newpage
\textbf{\deva{प्रतिपदार्थः} (Meaning of each word):}
\begin{table}[h!]
\centering
\renewcommand{\arraystretch}{1.2}
\begin{tabularx}{\textwidth}{@{} l l X X @{}}
\toprule
\textbf{अवधी} & \textbf{संस्कृतम्} & \textbf{English} & \textbf{Grammar \& Notes} \\
\midrule
\deva{तुम} \gotchaicon / \telu{తుమ} / \textit{tuma}  &  \deva{त्वम्}  &  You  &   \\
\deva{उपकार} / \telu{ఉపకార} / \textit{upakāra}  &  \deva{उपकारम्}  &  Favor / Service  &   \\
\deva{सुग्रीवहिं} \gotchaicon / \telu{సుగ్రీవహిఁ} / \textit{sugrīvahĩ}  &  \deva{सुग्रीवाय}  &  To Sugriva  & Objective suffix `-hiṃ` \\
\deva{कीन्हा} \gotchaicon / \telu{కీన్హా} / \textit{kīnhā}  &  \deva{कृतवान्}  &  Did  & Awadhi past tense \\
\deva{राम मिलाय} / \telu{రామ మిలాయ} / \textit{rāma milāya}  &  \deva{रामेण मेलयित्वा}  &  Having united with Rama  &   \\
\deva{राज पद} / \telu{రాజ పద} / \textit{rāja pada}  &  \deva{राज-पदम्}  &  Royal Throne  &   \\
\deva{दीन्हा} / \telu{దీన్హా} / \textit{dīnhā}  &  \deva{दत्तवान्}  &  Gave  & Rhymes with keenhaa \\
\bottomrule
\end{tabularx}
\end{table}

\begin{tcolorbox}[poetrybox]
\textbf{\deva{त्वम्} $\rightarrow$ \deva{तुम} \gotchaicon\  Pronoun Across Millennia:}\\
Sanskrit uses \deva{त्वम्} for ''you'' (nominative singular) that Awadhi collapses into the simple, \deva{तुम}: the exact same form that survived into modern Hindi, Urdu, and Bengali. 
Note also that \deva{तुम} in Awadhi works for all cases;
 there is no separate accusative or dative form.\\[0.4em]
\textbf{\deva{कृतवान्} $\rightarrow$ \deva{कीन्हा} \gotchaicon\  The Awadhi Past Tense Signature:}\\
\deva{कीन्हा} is the canonical Awadhi perfect past form of \deva{करना} (to do). 
Notice the dramatic reduction: Sanskrit requires an elaborate perfect participle \deva{कृतवान्} (''one who has done''), while Awadhi achieves the same meaning with a single crisp word. 
The rhyming pair \deva{कीन्हा} / \deva{दीन्हा} (did / gave) is a beautiful structural device Tulsidas uses to link this couplet's two actions musically.\\[0.4em]
\textbf{\deva{सुग्रीवहिं} — A Metrical Flex:}\\
\deva{सुग्रीवहिं} carries 5 matras in the rhythm map (\deva{सु}{\scriptsize(1)} + \deva{ग्री}{\scriptsize(2)} + \deva{व}{\scriptsize(1)} + \deva{हिं}{\scriptsize(2)} = 5), yet a phonologically valid 6-matra reading also exists (\deva{सु}{\scriptsize(2)} + \deva{ग्री}{\scriptsize(2)} + \deva{व}{\scriptsize(1)} + \deva{हि}{\scriptsize(1)} = 6). Compare this with \deva{पुत्र} from Chaupai 2 — both words contain a conjunct consonant (\deva{संयुक्ताक्षर}), yet they are weighted differently. The distinction rests on a classical \textit{Apavāda Sūtra} (exception rule) governing stress (\textit{Pratibhā}):

\begin{itemize}
  \item \textbf{\deva{पुत्र} (pu-tra):} The conjunct \deva{त्र} (t+ra) involves a \textit{Sparśa} (stop) consonant. 
  The hard stop forces weight onto the preceding \deva{पु}, which becomes \textit{Guru} by the standard \textit{Saṃyuktapūrva} rule. 
  Result: \textit{Guru}+\textit{Laghu} = 3 matras.

  \item \textbf{\deva{सुग्री} (su-grī):} The conjunct \deva{ग्री} involves a soft consonant followed by \deva{र} (Repha). 
  Sanskrit and Telugu prosody permit a \textit{Śithila} (soft/lax) exception here: when the cluster is pronounced lightly, the preceding short vowel need \emph{not} become \textit{Guru}, keeping \deva{सु} as \textit{Laghu}. 
  Result: \textit{Laghu}+\textit{Guru} = 3 matras, but the opening syllable is light.
\end{itemize}

In the 16-matra Chaupai framework, Tulsidas invokes this \textit{Śithila} license to hold \deva{सुग्रीवहिं} at 5 matras, fitting the meter perfectly. 
The 6-matra reading remains phonologically available for singers who choose to lean into the conjunct. 
This is \textit{Bhāva Pradhāna} prosody at its finest: the emotional weight of the meaning guides the syllabic weight of the delivery.
\end{tcolorbox}

\newpage


% ----- CHAUPAI 17 -----

\newpage

\subsection*{\deva{सप्तदश चौपाई} \(Seventeenth Chaupai\)}
\addcontentsline{toc}{subsection}{\deva{सप्तदश चौपाई} \(Seventeenth Chaupai\) — \deva{तुम्हरो मंत्र...}}

\begin{minipage}[t]{0.55\textwidth}
\vspace{0pt}

{\large \linespread{1.5}\selectfont
\deva{तुम्हरो मंत्र बिभीषन माना।}\\
\deva{लंकेस्वर भए सब जग जाना॥}
\par}

\vspace{0.5em}

{\large \linespread{1.5}\selectfont
\telu{తుమ్హరో మంత్ర బిభీషన మానా।}\\
\telu{లంకేశ్వర భఏ సబ జగ జానా॥}
\par}

\vspace{0.5em}

{\large \linespread{1.3}\selectfont
\textit{tumharo mantra bibhīṣana mānā |}\\
\textit{laṅkesvara bhae saba jaga jānā ||}
\par}
\end{minipage}%
\hfill
\begin{minipage}[t]{0.42\textwidth}
\vspace{0pt}
\begin{tcolorbox}[anchorbox]
\textbf{The Bridge of Advocacy:} In the *Valmiki Ramayana* (*Yuddha Kanda, 17:50-68*), when Vibhishana sought refuge, Sugriva and other leaders suspected him of being a spy. Hanuman was the sole advisor who delivered a brilliant, analytical speech to Lord Rama, advocating for Vibhishana's acceptance. Because Vibhishana was accepted and protected based on Hanuman's fair advocacy (this 'Mantra' or counsel to Rama), he was crowned King of Lanka (*Lankeshvara*). It shows how a true leader's strategic advocacy can lift others to greatness.
\vspace{0.5em}\newline Hanuman's brilliant, logical counsel on behalf of Vibhishana secured his sanctuary and his ultimate elevation to king.\newline\null\hfill\href{https://www.valmiki.iitk.ac.in/sloka?field_kanda_tid=6&language=dv&field_sarga_value=17}{\textit{Yuddha Kanda, 17:50-68}}
\end{tcolorbox}
\end{minipage}

\vspace{0.5em}
\subsubsection*{\textbf{Visual Rhythm Map:}}
\begin{table}[h!]
\centering
\renewcommand{\arraystretch}{1.2}
\setlength{\tabcolsep}{0pt}
\begin{tabularx}{\textwidth}{|*{16}{>{\centering\arraybackslash\hsize=1.0\hsize}X|}}
\hline
\multicolumn{4}{|>{\centering\arraybackslash\hsize=4.0\hsize}X|}{\textbf{\laghu{तु}\laghu{म्ह}\guru{रो}} \par (4)} &
\multicolumn{3}{>{\centering\arraybackslash\hsize=3.0\hsize}X|}{\textbf{\guru{मं}\laghu{त्र}} \par (3)} &
\multicolumn{5}{>{\centering\arraybackslash\hsize=5.0\hsize}X|}{\textbf{\laghu{बि}\guru{भी}\laghu{ष}\laghu{न}} \par (5)} &
\multicolumn{4}{>{\centering\arraybackslash\hsize=4.0\hsize}X|}{\textbf{\guru{मा}\guru{ना}} \par (4)} \\
\hline
\end{tabularx}
\vspace{-1.5em}
\end{table}

\begin{table}[h!]
\centering
\renewcommand{\arraystretch}{1.2}
\noindent\makebox[\textwidth][l]{%
  {%
  \setlength{\tabcolsep}{0pt}%
  \begin{tabularx}{1.0625\textwidth}{|*{17}{>{\centering\arraybackslash\hsize=1.0\hsize}X|}}
  \hline
  \multicolumn{6}{|>{\centering\arraybackslash\hsize=6.0\hsize}X|}{\textbf{\guru{लं}\guru{के}\laghu{स्व}\laghu{र}} \par (6)} &
  \multicolumn{3}{>{\centering\arraybackslash\hsize=3.0\hsize}X|}{\textbf{\laghu{भ}\guru{ए}} \par (3)} &
  \multicolumn{2}{>{\centering\arraybackslash\hsize=2.0\hsize}X|}{\textbf{\laghu{स}\laghu{ब}} \par (2)} &
  \multicolumn{2}{>{\centering\arraybackslash\hsize=2.0\hsize}X|}{\textbf{\laghu{ज}\laghu{ग}} \par (2)} &
  \multicolumn{4}{>{\centering\arraybackslash\hsize=4.0\hsize}X|}{\textbf{\guru{जा}\guru{ना}} \par (4)} \\
  \hline
  \end{tabularx}%
  }%
  \hspace{0.01\textwidth}%
  \begin{minipage}{0.1\textwidth}
  \vspace{-0.5em}
  \centering
  {\Huge \textbf{17}}
  \end{minipage}%
}
\vspace{-1.5em}
\end{table}

\vspace{0.5em}
\textbf{ \deva{सरल-संस्कृतम्}:}\\
 \deva{भवतः मन्त्रं (परामर्शं / उपदेशं) विभीषणः अङ्गीकृतवान्।}\\
 \deva{तेन सः लङ्केश्वरः (लङ्कायाः राजा) अभवत् इति सर्वं जगत् जानाति।} \\

Vibhishana heeded your counsel. 
By doing so, he became the King of Lanka, a fact the entire world knows.

\newpage

\vspace{1em}
\textbf{\deva{प्रतिपदार्थः} (Meaning of each word):}
\begin{table}[h!]
\centering
\renewcommand{\arraystretch}{1.2}
\begin{tabularx}{\textwidth}{@{} l l X X @{}}
\toprule
\textbf{अवधी} & \textbf{संस्कृतम्} & \textbf{English} & \textbf{Grammar \& Notes} \\
\midrule
\deva{तुम्हरो} / \telu{తుమ్హరో} / \textit{tumharo}  &  \deva{भवतः}  &  Your  &  Awadhi possessive; syncopated from \deva{तुम्हारो} \\
\deva{मंत्र} / \telu{మంత్ర} / \textit{mantra}  &  \deva{मन्त्रम् (उपदेशम्)}  &  Word / Counsel  &   \\
\deva{बिभीषन} / \telu{బిభీషన} / \textit{bibhīṣana}  &  \deva{विभीषणः}  &  Vibhishana  &  $v \rightarrow b$  \\
\deva{माना} / \telu{మానా} / \textit{mānā}  &  \deva{अङ्गीकृतवान्}  &  Accepted / Heeded  &   \\
\deva{लंकेस्वर} / \telu{లంకేశ్వర} / \textit{laṅkesvara}  &  \deva{लङ्केश्वरः}  &  King of Lanka  &  Compound: \deva{लंका} + \deva{ईश्वर} \\
\deva{भए} / \telu{భఏ} / \textit{bhae}  &  \deva{अभवत्}  &  Became  &   \\
\deva{सब जग} / \telu{సబ జగ} / \textit{saba jaga}  &  \deva{सर्वं जगत्}  &  The whole world  &   \\
\deva{जाना} / \telu{జానా} / \textit{jānā}  &  \deva{जानाति}  &  Knows  &   \\
\bottomrule
\end{tabularx}
\end{table}



\newpage


% ----- CHAUPAI 18 -----

\newpage

\subsection*{\deva{अष्टादश चौपाई} \(Eighteenth Chaupai\)}
\addcontentsline{toc}{subsection}{\deva{अष्टादश चौपाई} \(Eighteenth Chaupai\) — \deva{जुग सहस्र...}}

\begin{minipage}[t]{0.55\textwidth}
\vspace{0pt}

{\large \linespread{1.5}\selectfont
\deva{जुग सहस्र जोजन पर भानू।}\\
\deva{लील्यो ताहि मधुर फल जानू॥}
\par}

\vspace{0.5em}

{\large \linespread{1.5}\selectfont
\telu{జుగ సహస్ర జోజన పర భానూ।}\\
\telu{లీల్యో తాహి మధుర ఫల జానూ॥}
\par}

\vspace{0.5em}

{\large \linespread{1.3}\selectfont
\textit{juga sahasra jojana para bhānū |}\\
\textit{līlyo tāhi madhura phala jānū ||}
\par}
\end{minipage}%
\hfill
\begin{minipage}[t]{0.42\textwidth}
\vspace{0pt}
\begin{tcolorbox}[anchorbox]
\textbf{Fearless Curiosity:} As a child, he literally leaped for the sun because he wanted to understand it. Fearless curiosity and ambition are the primary engines of learning. Students should never be afraid to reach high or ask bold questions, even if they occasionally make mistakes along the way.
\vspace{0.5em}\newline Driven by fearless childhood curiosity, he leaped into the unknown, teaching us to never fear making bold mistakes while learning.\newline\null\hfill\href{https://www.valmiki.iitk.ac.in/sloka?field_kanda_tid=4&language=dv&field_sarga_value=66}{\textit{Kishkindha Kanda, 66:21-23}}
\end{tcolorbox}
\end{minipage}

\vspace{0.5em}
\subsubsection*{\textbf{Visual Rhythm Map:}}
\begin{table}[h!]
\centering
\renewcommand{\arraystretch}{1.2}
\setlength{\tabcolsep}{0pt}
\begin{tabularx}{\textwidth}{|*{16}{>{\centering\arraybackslash\hsize=1.0\hsize}X|}}
\hline
\multicolumn{2}{|>{\centering\arraybackslash\hsize=2.0\hsize}X|}{\textbf{\laghu{जु}\laghu{ग}} \par (2)} &
\multicolumn{4}{>{\centering\arraybackslash\hsize=4.0\hsize}X|}{\textbf{\laghu{स}\guru{ह}\laghu{स्र}} \par (4)} &
\multicolumn{4}{>{\centering\arraybackslash\hsize=4.0\hsize}X|}{\textbf{\guru{जो}\laghu{ज}\laghu{न}} \par (4)} &
\multicolumn{2}{>{\centering\arraybackslash\hsize=2.0\hsize}X|}{\textbf{\laghu{प}\laghu{र}} \par (2)} &
\multicolumn{4}{>{\centering\arraybackslash\hsize=4.0\hsize}X|}{\textbf{\guru{भा}\guru{नू}} \par (4)} \\
\hline
\end{tabularx}
\vspace{-1.5em}
\end{table}

\begin{table}[h!]
\centering
\renewcommand{\arraystretch}{1.2}
\setlength{\tabcolsep}{0pt}
\begin{tabularx}{\textwidth}{|*{16}{>{\centering\arraybackslash\hsize=1.0\hsize}X|}}
\hline
\multicolumn{4}{|>{\centering\arraybackslash\hsize=4.0\hsize}X|}{\textbf{\guru{ली}\guru{ल्यो}} \par (4)} &
\multicolumn{3}{>{\centering\arraybackslash\hsize=3.0\hsize}X|}{\textbf{\guru{ता}\laghu{हि}} \par (3)} &
\multicolumn{3}{>{\centering\arraybackslash\hsize=3.0\hsize}X|}{\textbf{\laghu{म}\laghu{धु}\laghu{र}} \par (3)} &
\multicolumn{2}{>{\centering\arraybackslash\hsize=2.0\hsize}X|}{\textbf{\laghu{फ}\laghu{ल}} \par (2)} &
\multicolumn{4}{>{\centering\arraybackslash\hsize=4.0\hsize}X|}{\textbf{\guru{जा}\guru{नू}} \par (4)} \\
\hline
\end{tabularx}
\vspace{-1.5em}
\end{table}

\vspace{0.5em}
\textbf{ \deva{सरल-संस्कृतम्}:}\\
 \deva{युग-सहस्र-योजन-दूरे स्थितं भानुं (सूर्यम्)}\\
 \deva{मधुरं फलम् इति मत्वा भवान् लीलया निगीर्णवान् (भक्षितवान्)।}\\

The sun, situated thousands of Yojanas away \\
thinking it to be a sweet fruit, you playfully swallowed it.


\vspace{1em}
\newpage
\textbf{\deva{प्रतिपदार्थः} (Meaning of each word):}
\begin{table}[h!]
\centering
\renewcommand{\arraystretch}{1.2}
\begin{tabularx}{\textwidth}{@{} l l X X @{}}
\toprule
\textbf{अवधी} & \textbf{संस्कृतम्} & \textbf{English} & \textbf{Grammar \& Notes} \\
\midrule
\deva{जुग} \glossaryicon / \telu{జుగ} / \textit{juga}  &  \deva{युग}  &  Yuga (Era/Distance)  & \\ 
\deva{सहस्र} / \telu{సహస్ర} / \textit{sahasra}  &  \deva{सहस्र}  &  Thousand  &   \\
\deva{जोजन} / \telu{జోజన} / \textit{jojana}  &  \deva{योजन}  &  Yojana (approx 8 miles)  & \\ 
\deva{पर} / \telu{పర} / \textit{para}  &  \deva{दूरे / परत्र}  &  Away / At a distance  &   \\
\deva{भानू} / \telu{భానూ} / \textit{bhānū}  &  \deva{भानुम्}  &  The Sun  & Long `ū` for rhyme \\
\deva{लील्यो} / \telu{లీల్యో} / \textit{līlyo}  &  \deva{लीलया निगीर्णवान्}  &  Swallowed playfully  &   \\
\deva{ताहि} / \telu{తాహి} / \textit{tāhi}  &  \deva{तम्}  &  It (the sun)  & Awadhi objective \\
\deva{मधुर फल} / \telu{మధుర ఫల} / \textit{madhura phala}  &  \deva{मधुर-फलम्}  &  Sweet fruit  &   \\
\deva{जानू} / \telu{జానూ} / \textit{jānū}  &  \deva{ज्ञात्वा / मत्वा}  &  Knowing / Thinking  & Long `ū` for rhyme \\
\bottomrule
\end{tabularx}
\end{table}

\begin{tcolorbox}[poetrybox]
\textbf{The Verse the Internet Loves — The 96-Million-Mile Question:}\\
This is the most-quoted verse of the Hanuman Chalisa in modern popular science discussions. The calculation goes:
\deva{जुग} (\textit{Yuga}) = 12,000 $\times$ \deva{सहस्र} = 1,000 $\times$ \deva{जोजन} $\approx$ 8 miles:
\[
12{,}000 \times 1{,}000 \times 8~\text{miles} ~=~ 96{,}000{,}000~\text{miles}
\]
The modern scientific distance from the Earth to the Sun is approximately 93,000,000 miles ($\approx$ 150 million km). This similarity has made the verse enormously popular.

A few points of honest scholarly context:
\begin{itemize}
  \item The Hanuman Chalisa is a devotional poem, not an astronomical treatise. Neither it nor Tulsidas's Ramcharitmanas makes an explicit claim of astronomical accuracy.
  \item \deva{जुग} most commonly refers to a unit of time, not a distance scalar.  
    It can also mean \textit{joining}, allowing us to interpret the verse as a praise of Hanuman for joining with the sun which is thousands of \textit{yojanas} away.
  \item We cannot determine whether Tulsidas intentionally encoded the distance between the Earth and the Sun, or whether this is a beautiful numerical coincidence.  
  The poet could have used common knowledge fo the time, or he might have used it to refere to a large number, while fitting the metrical requirements of the Chalisa. 
\end{itemize}
 Neither interpretation diminishes the devotional or literary power of the verse.
\end{tcolorbox}

\newpage


% ----- CHAUPAI 19 -----

\newpage

\subsection*{\deva{एकोनविंशतितम चौपाई} \(Nineteenth Chaupai\)}
\addcontentsline{toc}{subsection}{\deva{एकोनविंशतितम चौपाई} \(Nineteenth Chaupai\) — \deva{प्रभु मुद्रिका...}}

\begin{minipage}[t]{0.55\textwidth}
\vspace{0pt}

{\large \linespread{1.5}\selectfont
\deva{प्रभु मुद्रिका मेलि मुख माहीं।}\\
\deva{जलधि लांघि गये अचरज नाहीं॥}
\par}

\vspace{0.5em}

{\large \linespread{1.5}\selectfont
\telu{ప్రభు ముద్రికా మేలి ముఖ మాహీఁ।}\\
\telu{జలధి లాంఘి గయే అచరజ నాహీఁ॥}
\par}

\vspace{0.5em}

{\large \linespread{1.3}\selectfont
\textit{prabhu mudrikā meli mukha māhīṃ |}\\
\textit{jaladhi lāṃghi gaye acaraja nāhīṃ ||}
\par}
\end{minipage}%
\hfill
\begin{minipage}[t]{0.42\textwidth}
\vspace{0pt}
\begin{tcolorbox}[anchorbox]
\textbf{Focus Under Pressure:} He carried the absolute weight of a massive responsibility (Rama's signet ring) across a terrifying, treacherous ocean without dropping it. Taking on great responsibility in education or life requires unshakeable focus and the determination not to let distractions sink your goals.
\vspace{0.5em}\newline He carried an enormous responsibility across a treacherous, demon-patrolled ocean without ever losing focus.\newline\null\hfill\href{https://www.valmiki.iitk.ac.in/sloka?field_kanda_tid=5&language=dv&field_sarga_value=36}{\textit{Sundara Kanda, 36:2}}
\end{tcolorbox}
\end{minipage}

\vspace{0.5em}
\subsubsection*{\textbf{Visual Rhythm Map:}}
\begin{table}[h!]
\centering
\renewcommand{\arraystretch}{1.2}
\setlength{\tabcolsep}{0pt}
\begin{tabularx}{\textwidth}{|*{16}{>{\centering\arraybackslash\hsize=1.0\hsize}X|}}
\hline
\multicolumn{2}{|>{\centering\arraybackslash\hsize=2.0\hsize}X|}{\textbf{\laghu{प्र}\laghu{भु}} \par (2)} &
\multicolumn{5}{>{\centering\arraybackslash\hsize=5.0\hsize}X|}{\textbf{\guru{मु}\laghu{द्रि}\guru{का}} \par (5)} &
\multicolumn{3}{>{\centering\arraybackslash\hsize=3.0\hsize}X|}{\textbf{\guru{मे}\laghu{लि}} \par (3)} &
\multicolumn{2}{>{\centering\arraybackslash\hsize=2.0\hsize}X|}{\textbf{\laghu{मु}\laghu{ख}} \par (2)} &
\multicolumn{4}{>{\centering\arraybackslash\hsize=4.0\hsize}X|}{\textbf{\guru{मा}\guru{हीं}} \par (4)} \\
\hline
\end{tabularx}
\vspace{-1.5em}
\end{table}

\begin{table}[h!]
\centering
\renewcommand{\arraystretch}{1.2}
\noindent\makebox[\textwidth][l]{%
  {%
  \setlength{\tabcolsep}{0pt}%
  \begin{tabularx}{1.0625\textwidth}{|*{17}{>{\centering\arraybackslash\hsize=1.0\hsize}X|}}
  \hline
  \multicolumn{3}{|>{\centering\arraybackslash\hsize=3.0\hsize}X|}{\textbf{\laghu{ज}\laghu{ल}\laghu{धि}} \par (3)} &
  \multicolumn{3}{>{\centering\arraybackslash\hsize=3.0\hsize}X|}{\textbf{\guru{लां}\laghu{घि}} \par (3)} &
  \multicolumn{3}{>{\centering\arraybackslash\hsize=3.0\hsize}X|}{\textbf{\laghu{ग}\guru{ये}} \par (3)} &
  \multicolumn{4}{>{\centering\arraybackslash\hsize=4.0\hsize}X|}{\textbf{\laghu{अ}\laghu{च}\laghu{र}\laghu{ज}} \par (4)} &
  \multicolumn{4}{>{\centering\arraybackslash\hsize=4.0\hsize}X|}{\textbf{\guru{ना}\guru{हीं}} \par (4)} \\
  \hline
  \end{tabularx}%
  }%
  \hspace{0.01\textwidth}%
  \begin{minipage}{0.1\textwidth}
  \vspace{-0.5em}
  \centering
  {\Huge \textbf{17}}
  \end{minipage}%
}
\vspace{-1.5em}
\end{table}

\vspace{0.5em}
\textbf{ \deva{सरल-संस्कृतम्}:}\\
 \deva{प्रभोः अङ्गुलीयकं (मुद्रिकाम्) मुखे निधाय }\\
 \deva{भवान् समुद्रं (जलधिं) उल्लङ्घितवान्, अत्र किञ्चित् आश्चर्यं नास्ति।}

Keeping the Lord's ring in your mouth you leaped across the vast ocean, 
and there is no surprise in this!

\vspace{1em}
\begin{tcolorbox}[poetrybox]
\textbf{The 17-Beat Squeeze \& Singing Flow:}\\
Line 2 mathematically counts to 17 beats because \deva{गये} is fully 3 beats native. But because this poem is \textit{Bhava Pradhan} (driven by devotion), strict meter bends to the singer's emotion. Singers effortlessly compress \deva{गये} into 2 quick beats (\textit{ga-e}). The 16-beat rule is a flexible guide, and these anomalies are part of the lively, organic beauty of the Awadhi oral tradition!
\end{tcolorbox}
\newpage
\textbf{\deva{प्रतिपदार्थः} (Meaning of each word):}
\begin{table}[h!]
\centering
\renewcommand{\arraystretch}{1.2}
\begin{tabularx}{\textwidth}{@{} l l X X @{}}
\toprule
\textbf{अवधी} & \textbf{संस्कृतम्} & \textbf{English} & \textbf{Grammar \& Notes} \\
\midrule
\deva{प्रभु} / \telu{ప్రభు}  &  \deva{प्रभोः / रामस्य}  &  Lord's  &   \\
\deva{मुद्रिका} / \telu{ముద్రికా}  &  \deva{मुद्रिका / अङ्गुलीयकम्}  &  Ring  &   \\
\deva{मेलि} / \telu{మేలి}  &  \deva{निधाय}  &  Having kept  & Absolutive `-i` ending \\
\deva{मुख माहीं} / \telu{ముఖ మాహీఁ}  &  \deva{मुख-मध्ये}  &  Inside the mouth  &   \\
\deva{जलधि} / \telu{జలధి}  &  \deva{जलधिः (समुद्रः)}  &  Ocean  &   \\
\deva{लांघि} / \telu{లాంఘి}  &  \deva{उल्लङ्घ्य}  &  Having crossed/leapt  &   \\
\deva{गये} / \telu{గయే}  &  \deva{गतवान्}  &  Went  &   \\
\deva{अचरज} / \telu{అచరజ}  &  \deva{आश्चर्यम्}  &  Surprise  & Swara-bhakti \\ 
\deva{नाहीं} / \telu{నాహీఁ}  &  \deva{नास्ति}  &  Is not  &   \\
\bottomrule
\end{tabularx}
\end{table}



\newpage


% ----- CHAUPAI 20 -----

\newpage

\subsection*{\deva{विंशतितम चौपाई} \(Twentieth Chaupai\)}
\addcontentsline{toc}{subsection}{\deva{विंशतितम चौपाई} \(Twentieth Chaupai\) — \deva{दुर्गम काज...}}

\begin{minipage}[t]{0.55\textwidth}
\vspace{0pt}

{\large \linespread{1.5}\selectfont
\deva{दुर्गम काज जगत के जेते।}\\
\deva{सुगम अनुग्रह तुम्हरे तेते॥}
\par}

\vspace{0.5em}

{\large \linespread{1.5}\selectfont
\telu{దుర్గమ కాజ జగత కే జేతే।}\\
\telu{సుగమ అనుగ్రహ తుమ్హరే తేతే॥}
\par}

\vspace{0.5em}

{\large \linespread{1.3}\selectfont
\textit{durgama kāja jagata ke jete |}\\
\textit{sugama anugraha tumhare tete ||}
\par}
\end{minipage}%
\hfill
\begin{minipage}[t]{0.42\textwidth}
\vspace{0pt}
\begin{tcolorbox}[anchorbox]
\textbf{The Ultimate Problem Solver:} He takes 'impossible' tasks (\textit{Durgam}) and makes them manageable (\textit{Sugam}). When students face overwhelming exams, massive projects, or complex challenges, remembering his laser-like focus and unmatched work ethic provides a blueprint for making the impossible solvable.
\vspace{0.5em}\newline He specialized in taking wildly impossible tasks (Durgama) and breaking them down into solvable, manageable steps.\newline\null\hfill\href{https://www.valmiki.iitk.ac.in/sloka?field_kanda_tid=5&language=dv&field_sarga_value=39}{\textit{Sundara Kanda, 39:53}}
\end{tcolorbox}
\end{minipage}

\vspace{0.5em}
\subsubsection*{\textbf{Visual Rhythm Map:}}
\begin{table}[h!]
\centering
\renewcommand{\arraystretch}{1.2}
\setlength{\tabcolsep}{0pt}
\begin{tabularx}{\textwidth}{|*{16}{>{\centering\arraybackslash\hsize=1.0\hsize}X|}}
\hline
\multicolumn{4}{|>{\centering\arraybackslash\hsize=4.0\hsize}X|}{\textbf{\guru{दु}\laghu{र्ग}\laghu{म}} \par (4)} &
\multicolumn{3}{>{\centering\arraybackslash\hsize=3.0\hsize}X|}{\textbf{\guru{का}\laghu{ज}} \par (3)} &
\multicolumn{3}{>{\centering\arraybackslash\hsize=3.0\hsize}X|}{\textbf{\laghu{ज}\laghu{ग}\laghu{त}} \par (3)} &
\multicolumn{2}{>{\centering\arraybackslash\hsize=2.0\hsize}X|}{\textbf{\guru{के}} \par (2)} &
\multicolumn{4}{>{\centering\arraybackslash\hsize=4.0\hsize}X|}{\textbf{\guru{जे}\guru{ते}} \par (4)} \\
\hline
\end{tabularx}
\vspace{-1.5em}
\end{table}

\begin{table}[h!]
\centering
\renewcommand{\arraystretch}{1.2}
\setlength{\tabcolsep}{0pt}
\begin{tabularx}{\textwidth}{|*{16}{>{\centering\arraybackslash\hsize=1.0\hsize}X|}}
\hline
\multicolumn{3}{|>{\centering\arraybackslash\hsize=3.0\hsize}X|}{\textbf{\laghu{सु}\laghu{ग}\laghu{म}} \par (3)} &
\multicolumn{5}{>{\centering\arraybackslash\hsize=5.0\hsize}X|}{\textbf{\laghu{अ}\guru{नु}\laghu{ग्र}\laghu{ह}} \par (5)} &
\multicolumn{4}{>{\centering\arraybackslash\hsize=4.0\hsize}X|}{\textbf{\laghu{तु}\laghu{म्ह}\guru{रे}} \par (4)} &
\multicolumn{4}{>{\centering\arraybackslash\hsize=4.0\hsize}X|}{\textbf{\guru{ते}\guru{ते}} \par (4)} \\
\hline
\end{tabularx}
\vspace{-1.5em}
\end{table}

\vspace{0.5em}
\textbf{ \deva{सरल-संस्कृतम्}:}\\
 \deva{जगतः यावन्ति अपि दुर्गमाणि (कठिनानि) कार्याणि सन्ति}\\
 \deva{तानि सर्वाणि भवतः अनुग्रहेण (कृपया) सुगमानि भवन्ति॥}\\

Whatever difficult tasks exist in the world\\
they all become easily attainable by your grace.


\vspace{1em}
\newpage
\textbf{\deva{प्रतिपदार्थः} (Meaning of each word):}
\begin{table}[h!]
\centering
\renewcommand{\arraystretch}{1.2}
\begin{tabularx}{\textwidth}{@{} l l X X @{}}
\toprule
\textbf{अवधी} & \textbf{संस्कृतम्} & \textbf{English} & \textbf{Grammar \& Notes} \\
\midrule
\deva{दुर्गम} / \telu{దుర్గమ} / \textit{durgama}  &  \deva{दुर्गमानि}  &  Difficult / Hard to cross  &   \\
\deva{काज} / \telu{కాజ} / \textit{kāja}  &  \deva{कार्याणि}  &  Tasks / Work  & \\ 
\deva{जगत के} / \telu{జగత కే} / \textit{jagata ke}  &  \deva{जगतः}  &  Of the world  & Genitive Postposition `ke` \\
\deva{जेते} / \telu{జేతే} / \textit{jete}  &  \deva{यावन्ति}  &  However many  & Rhymes with `tete` \\
\deva{सुगम} / \telu{సుగమ} / \textit{sugama}  &  \deva{सुగమాని}  &  Accessible / Easy  &   \\
\deva{अनुग्रह} / \telu{అనుగ్రహ} / \textit{anugraha}  &  \deva{अनुग्रहेण}  &  By grace / blessing  &   \\
\deva{तुम्हरे} \gotchaicon / \telu{తుమ్హరే} / \textit{tumhare}  &  \deva{భవతః}  &  Your  &   \\
\deva{तेते} / \telu{తేతే} / \textit{tete}  &  \deva{तावन्ति}  &  That many  &   \\
\bottomrule
\end{tabularx}
\end{table}

\begin{tcolorbox}[gotchabox]
\begin{itemize}
    \item \textbf{Aspirated Consonants (\deva{म्ह} and \deva{न्ह}):} In words like \deva{तुम्हरे} (Tumhare) and \deva{कीन्हा} (Keenhaa), the `h` merges with `m` and `n` to form a single aspirated sound. In classical Awadhi meter, this does \textit{not} make the preceding short vowel heavy. Therefore, \deva{तु} (tu) in Tumhare remains exactly 1 beat!
\end{itemize}
\end{tcolorbox}

\newpage


% ----- CHAUPAI 21 -----

\newpage

\subsection*{\deva{एकविंशतितम चौपाई} \(Twenty-First Chaupai\)}
\addcontentsline{toc}{subsection}{\deva{एकविंशतितम चौपाई} \(Twenty-First Chaupai\) — \deva{राम दुआरे...}}

\begin{minipage}[t]{0.55\textwidth}
\vspace{0pt}

{\large \linespread{1.5}\selectfont
\deva{राम दुआरे तुम रखवारे।}\\
\deva{होत न आज्ञा बिनु पैसारे॥}
\par}

\vspace{0.5em}

{\large \linespread{1.5}\selectfont
\telu{రామ దుఆరే తుమ రఖవారే।}\\
\telu{హోత న ఆజ్ఞా బిను పైసారే॥}
\par}

\vspace{0.5em}

{\large \linespread{1.3}\selectfont
\textit{rāma duāre tuma rakhavāre |}\\
\textit{hota na ājñā binu paisāre ||}
\par}
\end{minipage}%
\hfill
\begin{minipage}[t]{0.42\textwidth}
\vspace{0pt}
\begin{tcolorbox}[anchorbox]
\textbf{The Gatekeeper of Discipline:} He guards the door to success (\textit{Ram Duaare Tum Rakhvare}). Nothing worthwhile in life is achieved without passing through the gates of discipline. Hanuman represents the hard work, routine, and discipline required before anyone can attain mastery or peace.
\vspace{0.5em}\newline Positioned at the very gates of the King, he became the ultimate symbol that success requires passing through strict discipline.\newline\null\hfill\href{https://www.valmiki.iitk.ac.in/sloka?field_kanda_tid=6&language=dv&field_sarga_value=128}{\textit{Yuddha Kanda, 128:88}}
\end{tcolorbox}
\end{minipage}

\vspace{0.5em}
\subsubsection*{\textbf{Visual Rhythm Map:}}
\begin{table}[h!]
\centering
\renewcommand{\arraystretch}{1.2}
\setlength{\tabcolsep}{0pt}
\begin{tabularx}{\textwidth}{|*{16}{>{\centering\arraybackslash\hsize=1.0\hsize}X|}}
\hline
\multicolumn{3}{|>{\centering\arraybackslash\hsize=3.0\hsize}X|}{\textbf{\guru{रा}\laghu{म}} \par (3)} &
\multicolumn{5}{>{\centering\arraybackslash\hsize=5.0\hsize}X|}{\textbf{\laghu{दु}\guru{आ}\guru{रे}} \par (5)} &
\multicolumn{2}{>{\centering\arraybackslash\hsize=2.0\hsize}X|}{\textbf{\laghu{तु}\laghu{म}} \par (2)} &
\multicolumn{6}{>{\centering\arraybackslash\hsize=6.0\hsize}X|}{\textbf{\laghu{र}\laghu{ख}\guru{वा}\guru{रे}} \par (6)} \\
\hline
\end{tabularx}
\vspace{-1.5em}
\end{table}

\begin{table}[h!]
\centering
\renewcommand{\arraystretch}{1.2}
\setlength{\tabcolsep}{0pt}
\begin{tabularx}{\textwidth}{|*{16}{>{\centering\arraybackslash\hsize=1.0\hsize}X|}}
\hline
\multicolumn{3}{|>{\centering\arraybackslash\hsize=3.0\hsize}X|}{\textbf{\guru{हो}\laghu{त}} \par (3)} &
\multicolumn{1}{>{\centering\arraybackslash\hsize=1.0\hsize}X|}{\textbf{\laghu{न}} \par (1)} &
\multicolumn{4}{>{\centering\arraybackslash\hsize=4.0\hsize}X|}{\textbf{\guru{आ}\guru{ज्ञा}} \par (4)} &
\multicolumn{2}{>{\centering\arraybackslash\hsize=2.0\hsize}X|}{\textbf{\laghu{बि}\laghu{नु}} \par (2)} &
\multicolumn{6}{>{\centering\arraybackslash\hsize=6.0\hsize}X|}{\textbf{\guru{पै}\guru{सा}\guru{रे}} \par (6)} \\
\hline
\end{tabularx}
\vspace{-1.5em}
\end{table}

\vspace{0.5em}
\textbf{ \deva{सरल-संस्कृतम्}:}\\
 \deva{श्रीरामस्य द्वारे भवान् एव रक्षकः (द्वारपालः) अस्ति।}\\
 \deva{भवतः आज्ञां विना तत्र प्रवेशः न भवति॥}\\

At the door of Lord Rama, you are the protector/guard.\\
Without your permission, entry is not possible.


\vspace{1em}
\newpage
\textbf{\deva{प्रतिपदार्थः} (Meaning of each word):}
\begin{table}[h!]
\centering
\renewcommand{\arraystretch}{1.2}
\begin{tabularx}{\textwidth}{@{} l l X X @{}}
\toprule
\textbf{अवधी} & \textbf{संस्कृतम्} & \textbf{English} & \textbf{Grammar \& Notes} \\
\midrule
\deva{राम दुआरे} / \telu{రామ దుఆరే}  &  \deva{रामस्य द्वारे}  &  At Rama's door  & `duaare` is locative \\
\deva{तुम} \gotchaicon / \telu{తుమ}  &  \deva{त्वम्}  &  You  &   \\
\deva{रखवारे} / \telu{రఖవారే}  &  \deva{रक्षकः}  &  Protector / Guard  & \\ 
\deva{होत न} / \telu{హోత న}  &  \deva{न भवति}  &  Does not happen  &   \\
\deva{आज्ञा} / \telu{ఆజ్ఞా}  &  \deva{आज्ञाम्}  &  Command / Permission  &   \\
\deva{बिनु} / \telu{బిను}  &  \deva{विना}  &  Without  &   \\
\deva{पैसारे} / \telu{పైసారే}  &  \deva{प्रवेशः}  &  Entry / Passage  & Awadhi noun \\
\bottomrule
\end{tabularx}
\end{table}

\begin{tcolorbox}[poetrybox]
\textbf{\deva{त्वम्} $\rightarrow$ \deva{तुम} \gotchaicon\ — The Eternal ''You'':}\\
Sanskrit \deva{त्वम्} (nominative ''you'') traces back to Proto-Indo-European \textit{*túh₂} — the same root that gives English ''thou,'' German \textit{du}, and Persian \textit{tu}. Awadhi condenses this to the unflinching \deva{तुम}, stripping away all case distinctions. Here the word carries exceptional weight: Tulsidas addresses Hanuman directly and personally, making this verse a powerful statement that \textit{only you} — not any other deity — stand guard at Rama's door. It is an intimate, almost devotional ''you.''
\end{tcolorbox}

\newpage


% ----- CHAUPAI 22 -----

\newpage

\subsection*{\deva{द्वाविंशतितम चौपाई} \(Twenty-Second Chaupai\)}
\addcontentsline{toc}{subsection}{\deva{द्वाविंशतितम चौपाई} \(Twenty-Second Chaupai\) — \deva{सब सुख...}}

\begin{minipage}[t]{0.55\textwidth}
\vspace{0pt}

{\large \linespread{1.5}\selectfont
\deva{सब सुख लहै तुम्हारी सरना।}\\
\deva{तुम रच्छक काहू को डर ना॥}
\par}

\vspace{0.5em}

{\large \linespread{1.5}\selectfont
\telu{సబ సుఖ లహై తుమ్హారీ సరనా।}\\
\telu{తుమ రచ్ఛక కాహూ కో డర నా॥}
\par}

\vspace{0.5em}

{\large \linespread{1.3}\selectfont
\textit{saba sukha lahai tumhārī saranā |}\\
\textit{tuma racchaka kāhū ko ḍara nā ||}
\par}
\end{minipage}%
\hfill
\begin{minipage}[t]{0.42\textwidth}
\vspace{0pt}
\begin{tcolorbox}[anchorbox]
\textbf{Creating Psychological Safety:} A good team captain or leader provides an umbrella of protection so their team can operate without fear (\textit{Dara na}). When you take shelter in good habits and associate with strong, supportive friends, daily anxieties naturally fade away.
\vspace{0.5em}\newline By speaking gently and confidently, he provided immediate psychological refuge to someone trapped in deep anxiety.\newline\null\hfill\href{https://www.valmiki.iitk.ac.in/sloka?field_kanda_tid=5&language=dv&field_sarga_value=34}{\textit{Sundara Kanda, 34:38-40}}
\end{tcolorbox}
\end{minipage}

\vspace{0.5em}
\subsubsection*{\textbf{Visual Rhythm Map:}}
\begin{table}[h!]
\centering
\renewcommand{\arraystretch}{1.2}
\setlength{\tabcolsep}{0pt}
\begin{tabularx}{\textwidth}{|*{16}{>{\centering\arraybackslash\hsize=1.0\hsize}X|}}
\hline
\multicolumn{2}{|>{\centering\arraybackslash\hsize=2.0\hsize}X|}{\textbf{\laghu{स}\laghu{ब}} \par (2)} &
\multicolumn{2}{>{\centering\arraybackslash\hsize=2.0\hsize}X|}{\textbf{\laghu{सु}\laghu{ख}} \par (2)} &
\multicolumn{3}{>{\centering\arraybackslash\hsize=3.0\hsize}X|}{\textbf{\laghu{ल}\guru{है}} \par (3)} &
\multicolumn{5}{>{\centering\arraybackslash\hsize=5.0\hsize}X|}{\textbf{\laghu{तु}\guru{म्हा}\guru{री}} \par (5)} &
\multicolumn{4}{>{\centering\arraybackslash\hsize=4.0\hsize}X|}{\textbf{\laghu{स}\laghu{र}\guru{ना}} \par (4)} \\
\hline
\end{tabularx}
\vspace{-1.5em}
\end{table}

\begin{table}[h!]
\centering
\renewcommand{\arraystretch}{1.2}
\setlength{\tabcolsep}{0pt}
\begin{tabularx}{\textwidth}{|*{16}{>{\centering\arraybackslash\hsize=1.0\hsize}X|}}
\hline
\multicolumn{2}{|>{\centering\arraybackslash\hsize=2.0\hsize}X|}{\textbf{\laghu{तु}\laghu{म}} \par (2)} &
\multicolumn{4}{>{\centering\arraybackslash\hsize=4.0\hsize}X|}{\textbf{\guru{र}\laghu{च्छ}\laghu{क}} \par (4)} &
\multicolumn{4}{>{\centering\arraybackslash\hsize=4.0\hsize}X|}{\textbf{\guru{का}\guru{हू}} \par (4)} &
\multicolumn{2}{>{\centering\arraybackslash\hsize=2.0\hsize}X|}{\textbf{\guru{को}} \par (2)} &
\multicolumn{2}{>{\centering\arraybackslash\hsize=2.0\hsize}X|}{\textbf{\laghu{ड}\laghu{र}} \par (2)} &
\multicolumn{2}{>{\centering\arraybackslash\hsize=2.0\hsize}X|}{\textbf{\guru{ना}} \par (2)} \\
\hline
\end{tabularx}
\vspace{-1.5em}
\end{table}

\vspace{0.5em}
\textbf{ \deva{सरल-संस्कृतम्}:}\\
 \deva{यः भवतः शरणम् आगच्छति, सः सर्वाणि सुखानि लभते।}\\
 \deva{यदा भवान् रक्षकः अस्ति, तदा कस्यापि भयं नास्ति॥}\\

Whoever takes your shelter, obtains all comforts and joys.\\
When you are the protector, there is nothing to fear.


\vspace{1em}
\newpage
\textbf{\deva{प्रतिपदार्थः} (Meaning of each word):}
\begin{table}[h!]
\centering
\renewcommand{\arraystretch}{1.2}
\begin{tabularx}{\textwidth}{@{} l l X X @{}}
\toprule
\textbf{अवधी} & \textbf{संस्कृतम्} & \textbf{English} & \textbf{Grammar \& Notes} \\
\midrule
\deva{सब सुख} / \telu{సబ సుఖ} / \textit{saba sukha}  &  \deva{सर्वाणि सुखानि}  &  All happiness  &   \\
\deva{लहै} / \telu{లహై} / \textit{lahai}  &  \deva{लभते}  &  Obtains / gets  & Verb ending in `-ai` \\
\deva{तुम्हारी} / \telu{తుమ్హారీ} / \textit{tumhārī}  &  \deva{भवतः}  &  Your  &   \\
\deva{सरना} / \telu{సరనా} / \textit{saranā}  &  \deva{शरणम्}  &  Shelter / Refuge  & \\ 
\deva{तुम} \gotchaicon / \telu{తుమ} / \textit{tuma}  &  \deva{त्वम्}  &  You  &   \\
\deva{रच्छक} \gotchaicon / \telu{రచ్ఛక} / \textit{racchaka}  &  \deva{रक्षकः}  &  Protector  & \\ 
\deva{काहू को} / \telu{కాహూ కో} / \textit{kāhū ko}  &  \deva{कस्यापि}  &  For anyone  &   \\
\deva{डर ना} / \telu{డర నా} / \textit{ḍara nā}  &  \deva{भयं नास्ति}  &  No fear  &   \\
\bottomrule
\end{tabularx}
\end{table}

\begin{tcolorbox}[gotchabox]
\begin{itemize}
    \item \textbf{The \deva{क्ष} $\rightarrow$ \deva{च्छ} Consonant Shift:} Look closely at \deva{रच्छक} (Racchaka). The mighty Sanskrit conjunct \deva{क्ष} (kṣa) is very hard on the tongue. In Awadhi, it either softens to \deva{ख} (kha) as in \deva{लखन} (Lakhan for Lakshman), or it softens into the palatal \deva{च्छ} (ccha) as we see here (\deva{रक्षक} $\rightarrow$ \deva{रच्छक}).
\end{itemize}
\end{tcolorbox}

\newpage


% ----- CHAUPAI 23 -----

\newpage

\subsection*{\deva{त्रयोविंशतितम चौपाई} \(Twenty-Third Chaupai\)}
\addcontentsline{toc}{subsection}{\deva{त्रयोविंशतितम चौपाई} \(Twenty-Third Chaupai\) — \deva{आपन तेज...}}

\begin{minipage}[t]{0.55\textwidth}
\vspace{0pt}

{\large \linespread{1.5}\selectfont
\deva{आपन तेज सम्हारो आपै।}\\
\deva{तीनों लोक हाँक तें काँपै॥}
\par}

\vspace{0.5em}

{\large \linespread{1.5}\selectfont
\telu{ఆపన తేజ సమ్హారో ఆపై।}\\
\telu{తీనోం లోక హాఁక తేఁ కాఁపై॥}
\par}

\vspace{0.5em}

{\large \linespread{1.3}\selectfont
\textit{āpana teja samhāro āpai |}\\
\textit{tīnõ loka hā̃ka tẽ kā̃pai ||}
\par}
\end{minipage}%
\hfill
\begin{minipage}[t]{0.42\textwidth}
\vspace{0pt}
\begin{tcolorbox}[anchorbox]
\textbf{Self-Regulation:} He possessed immense, earth-shaking energy and talent, but he kept it under perfect control (\textit{Samharo Aapai}). Raw intelligence, talent, and energy are completely useless without the self-discipline to regulate and direct them properly.
\vspace{0.5em}\newline He completely regulated his blinding, destructive energy, using it only precisely when the situation demanded it.\newline\null\hfill\href{https://www.valmiki.iitk.ac.in/sloka?field_kanda_tid=5&language=dv&field_sarga_value=42}{\textit{Sundara Kanda, 42:34}}
\end{tcolorbox}
\end{minipage}

\vspace{0.5em}
\subsubsection*{\textbf{Visual Rhythm Map:}}
\begin{table}[h!]
\centering
\renewcommand{\arraystretch}{1.2}
\setlength{\tabcolsep}{0pt}
\begin{tabularx}{\textwidth}{|*{16}{>{\centering\arraybackslash\hsize=1.0\hsize}X|}}
\hline
\multicolumn{4}{|>{\centering\arraybackslash\hsize=4.0\hsize}X|}{\textbf{\guru{आ}\laghu{प}\laghu{न}} \par (4)} &
\multicolumn{3}{>{\centering\arraybackslash\hsize=3.0\hsize}X|}{\textbf{\guru{ते}\laghu{ज}} \par (3)} &
\multicolumn{5}{>{\centering\arraybackslash\hsize=5.0\hsize}X|}{\textbf{\laghu{स}\guru{म्हा}\guru{रो}} \par (5)} &
\multicolumn{4}{>{\centering\arraybackslash\hsize=4.0\hsize}X|}{\textbf{\guru{आ}\guru{पै}} \par (4)} \\
\hline
\end{tabularx}
\vspace{-1.5em}
\end{table}

\begin{table}[h!]
\centering
\renewcommand{\arraystretch}{1.2}
\setlength{\tabcolsep}{0pt}
\begin{tabularx}{\textwidth}{|*{16}{>{\centering\arraybackslash\hsize=1.0\hsize}X|}}
\hline
\multicolumn{4}{|>{\centering\arraybackslash\hsize=4.0\hsize}X|}{\textbf{\guru{ती}\guru{नों}} \par (4)} &
\multicolumn{3}{>{\centering\arraybackslash\hsize=3.0\hsize}X|}{\textbf{\guru{लो}\laghu{क}} \par (3)} &
\multicolumn{3}{>{\centering\arraybackslash\hsize=3.0\hsize}X|}{\textbf{\guru{हाँ}\laghu{क}} \par (3)} &
\multicolumn{2}{>{\centering\arraybackslash\hsize=2.0\hsize}X|}{\textbf{\guru{तें}} \par (2)} &
\multicolumn{4}{>{\centering\arraybackslash\hsize=4.0\hsize}X|}{\textbf{\guru{काँ}\guru{पै}} \par (4)} \\
\hline
\end{tabularx}
\vspace{-1.5em}
\end{table}

\vspace{0.5em}
\textbf{ \deva{सरल-संस्कृतम्}:}\\
 \deva{भवतः प्रचण्डं तेजः भवान् स्वयम् एव सम्भालितुं शक्नोति।}\\
 \deva{भवतः एकया गर्जनया (हुङ्कारेण) त्रयः लोकाः कम्पन्ते॥}\\

Only you yourself can control your own immense splendor and power.\\
At your single roar, all three worlds tremble.


\vspace{1em}
\newpage
\textbf{\deva{प्रतिपदार्थः} (Meaning of each word):}
\begin{table}[h!]
\centering
\renewcommand{\arraystretch}{1.2}
\begin{tabularx}{\textwidth}{@{} l l X X @{}}
\toprule
\textbf{अवधी} & \textbf{संस्कृतम्} & \textbf{English} & \textbf{Grammar \& Notes} \\
\midrule
\deva{आपन} / \telu{ఆపన} / \textit{āpana}  &  \deva{आत्मीयम्}  &  Own  &   \\
\deva{तेज} / \telu{తేజ} / \textit{teja}  &  \deva{तेजः}  &  Power / Splendor  &   \\
\deva{सम्हारो} / \telu{సమ్హారో} / \textit{samhāro}  &  \deva{सम्भालनम् / नियन्त्रणम्}  &  Control / Handle  &   \\
\deva{आपै} / \telu{ఆపై} / \textit{āpai}  &  \deva{स्वयम् एव}  &  You alone / By yourself  &   \\
\deva{तीनों लोक} / \telu{తీనోం లోక} / \textit{tīnõ loka}  &  \deva{त्रयः लोकाः}  &  Three worlds  &   \\
\deva{हाँक तें} / \telu{హాఁక తేఁ} / \textit{hā̃ka tẽ}  &  \deva{हुङ्कारेण / गर्जनया}  &  By your roar  &   \\
\deva{काँपै} / \telu{కాఁపై} / \textit{kā̃pai}  &  \deva{कम्पन्ते}  &  Tremble / Shiver  &   \\
\bottomrule
\end{tabularx}
\end{table}



\newpage


% ----- CHAUPAI 24 -----

\newpage

\subsection*{\deva{चतुर्विंशतितम चौपाई} \(Twenty-Fourth Chaupai\)}
\addcontentsline{toc}{subsection}{\deva{चतुर्विंशतितम चौपाई} \(Twenty-Fourth Chaupai\) — \deva{भूत पिसाच...}}

\begin{minipage}[t]{0.55\textwidth}
\vspace{0pt}

{\large \linespread{1.5}\selectfont
\deva{भूत पिसाच निकट नहिं आवै।}\\
\deva{महाबीर जब नाम सुनावै॥}
\par}

\vspace{0.5em}

{\large \linespread{1.5}\selectfont
\telu{భూత పిసాచ నికట నహిఁ ఆవై।}\\
\telu{మహాబీర జబ నామ సునావై॥}
\par}

\vspace{0.5em}

{\large \linespread{1.3}\selectfont
\textit{bhūta pisāca nikaṭa nahĩ āvai |}\\
\textit{mahābīra jaba nāma sunāvai ||}
\par}
\end{minipage}%
\hfill
\begin{minipage}[t]{0.42\textwidth}
\vspace{0pt}
\begin{tcolorbox}[anchorbox]
\textbf{Mental Armor:} The "ghosts" (\textit{Bhuta Pisacha}) represent performance anxiety and self-doubt. In the *Valmiki Ramayana* (*Sundara Kanda, 50:14-16*), when standing captive alone in Ravana's hostile court, Hanuman maintains complete composure. He objectively observes the grandeur of his enemies without being intimidated or losing focus. Keeping strong role models in mind acts as mental armor against external pressures.
\vspace{0.5em}\newline Standing alone in a hostile court, Hanuman's objective composure acted as mental armor, keeping him focused and unintimidated.\newline\null\hfill\href{https://www.valmiki.iitk.ac.in/sloka?field_kanda_tid=5&language=dv&field_sarga_value=50}{\textit{Sundara Kanda, 50:14-16}}
\end{tcolorbox}
\end{minipage}

\vspace{0.5em}
\subsubsection*{\textbf{Visual Rhythm Map:}}
\begin{table}[h!]
\centering
\renewcommand{\arraystretch}{1.2}
\setlength{\tabcolsep}{0pt}
\begin{tabularx}{\textwidth}{|*{16}{>{\centering\arraybackslash\hsize=1.0\hsize}X|}}
\hline
\multicolumn{3}{|>{\centering\arraybackslash\hsize=3.0\hsize}X|}{\textbf{\guru{भू}\laghu{त}} \par (3)} &
\multicolumn{4}{>{\centering\arraybackslash\hsize=4.0\hsize}X|}{\textbf{\laghu{पि}\guru{सा}\laghu{च}} \par (4)} &
\multicolumn{3}{>{\centering\arraybackslash\hsize=3.0\hsize}X|}{\textbf{\laghu{नि}\laghu{क}\laghu{ट}} \par (3)} &
\multicolumn{2}{>{\centering\arraybackslash\hsize=2.0\hsize}X|}{\textbf{\laghu{न}\laghu{हिं}} \par (2)} &
\multicolumn{4}{>{\centering\arraybackslash\hsize=4.0\hsize}X|}{\textbf{\guru{आ}\guru{वै}} \par (4)} \\
\hline
\end{tabularx}
\vspace{-1.5em}
\end{table}

\begin{table}[h!]
\centering
\renewcommand{\arraystretch}{1.2}
\setlength{\tabcolsep}{0pt}
\begin{tabularx}{\textwidth}{|*{16}{>{\centering\arraybackslash\hsize=1.0\hsize}X|}}
\hline
\multicolumn{6}{|>{\centering\arraybackslash\hsize=6.0\hsize}X|}{\textbf{\laghu{म}\guru{हा}\guru{बी}\laghu{र}} \par (6)} &
\multicolumn{2}{>{\centering\arraybackslash\hsize=2.0\hsize}X|}{\textbf{\laghu{ज}\laghu{ब}} \par (2)} &
\multicolumn{3}{>{\centering\arraybackslash\hsize=3.0\hsize}X|}{\textbf{\guru{ना}\laghu{म}} \par (3)} &
\multicolumn{5}{>{\centering\arraybackslash\hsize=5.0\hsize}X|}{\textbf{\laghu{सु}\guru{ना}\guru{वै}} \par (5)} \\
\hline
\end{tabularx}
\vspace{-1.5em}
\end{table}

\vspace{0.5em}
\textbf{ \deva{सरल-संस्कृतम्}:}\\
 \deva{भूताः पिशाचाः च समीपम् अपि न आगच्छन्ति}\\
 \deva{यदा कोऽपि "महावीर" इति भवतः नाम श्रावयति (जपति)॥}\\

Ghosts and evil spirits do not even come near\\
whenever your name 'Mahaveer' is chanted aloud.


\vspace{1em}
\newpage
\textbf{\deva{प्रतिपदार्थः} (Meaning of each word):}
\begin{table}[h!]
\centering
\renewcommand{\arraystretch}{1.2}
\begin{tabularx}{\textwidth}{@{} l l X X @{}}
\toprule
\textbf{अवधी} & \textbf{संस्कृतम्} & \textbf{English} & \textbf{Grammar \& Notes} \\
\midrule
\deva{भूत} / \telu{భూత} / \textit{bhūta}  &  \deva{भूताः}  &  Ghosts  &   \\
\deva{पिसाच} / \telu{పిసాచ} / \textit{pisāca}  &  \deva{पिशाचाः}  &  Evil spirits / Ghouls  & \\ 
\deva{निकट} / \telu{నికట} / \textit{nikaṭa}  &  \deva{निकटम् / समीपम्}  &  Near / Close  &   \\
\deva{नहिं आवै} / \telu{నహిఁ ఆవై} / \textit{nahĩ āvai}  &  \deva{न आगच्छन्ति}  &  Do not come  &   \\
\deva{महाबीर} / \telu{మహాబీర} / \textit{mahābīra}  &  \deva{महावीरः}  &  Mahaveer (Great Hero)  & \\ 
\deva{जब} / \telu{జబ} / \textit{jaba}  &  \deva{यदा}  &  When  & \\ 
\deva{नाम} / \telu{నామ} / \textit{nāma}  &  \deva{नाम}  &  Name  &   \\
\deva{सुनावै} / \telu{సునావై} / \textit{sunāvai}  &  \deva{श्रावयति}  &  Is heard / recited  & Causative verb \\
\bottomrule
\end{tabularx}
\end{table}

\newpage



% ==========================================
\section*{Group 6: The Remover of Obstacles \& Suffering}
\addcontentsline{toc}{section}{Group 6: The Remover of Obstacles \& Suffering}
% ==========================================

% ----- CHAUPAI 25 -----
\subsection*{\deva{पञ्चविंशतितम चौपाई} (Twenty-Fifth Chaupai)}
\addcontentsline{toc}{subsection}{\deva{पञ्चविंशतितम चौपाई} (Twenty-Fifth Chaupai) — \deva{नासै रोग...}}

\begin{minipage}[t]{0.55\textwidth}
\vspace{0pt}

{\large \linespread{1.5}\selectfont
\deva{नासै रोग हरै सब पीरा।}\\
\deva{जपत निरन्तर हनुमत बीरा॥}
\par}

\vspace{0.5em}

{\large \linespread{1.5}\selectfont
\telu{నాసై రోగ హరై సబ పీరా।}\\
\telu{జపత నిరంతర హనుమత బీరా॥}
\par}

\vspace{0.5em}

{\large \linespread{1.3}\selectfont
\textit{nāsai roga harai saba pīrā |}\\
\textit{japata nirantara hanumata bīrā ||}
\par}
\end{minipage}%
\hfill
\begin{minipage}[t]{0.42\textwidth}
\vspace{0pt}
\begin{tcolorbox}[anchorbox]
\textbf{Consistency is the Cure:} Just as studying requires daily effort, maintaining mental health and overcoming life's pains (\textit{Peera}) requires constant (\textit{Nirantara}), disciplined practice. You cannot cram for character; it must be built consistently over time.
\vspace{0.5em}\newline He proved that consistent, relentless action acts as the ultimate medicine for healing deep physical and mental exhaustion.\newline\null\hfill\href{https://www.valmiki.iitk.ac.in/sloka?field_kanda_tid=6&language=dv&field_sarga_value=74}{\textit{Yuddha Kanda, 74:73-75}}
\end{tcolorbox}
\end{minipage}

\vspace{0.5em}
\subsubsection*{\textbf{Visual Rhythm Map:}}
\begin{table}[h!]
\centering
\renewcommand{\arraystretch}{1.2}
\setlength{\tabcolsep}{0pt}
\begin{tabularx}{\textwidth}{|*{16}{>{\centering\arraybackslash\hsize=1.0\hsize}X|}}
\hline
\multicolumn{4}{|>{\centering\arraybackslash\hsize=4.0\hsize}X|}{\textbf{\guru{ना}\guru{सै}} \par (4)} &
\multicolumn{3}{>{\centering\arraybackslash\hsize=3.0\hsize}X|}{\textbf{\guru{रो}\laghu{ग}} \par (3)} &
\multicolumn{3}{>{\centering\arraybackslash\hsize=3.0\hsize}X|}{\textbf{\laghu{ह}\guru{रै}} \par (3)} &
\multicolumn{2}{>{\centering\arraybackslash\hsize=2.0\hsize}X|}{\textbf{\laghu{स}\laghu{ब}} \par (2)} &
\multicolumn{4}{>{\centering\arraybackslash\hsize=4.0\hsize}X|}{\textbf{\guru{पी}\guru{रा}} \par (4)} \\
\hline
\end{tabularx}
\vspace{-1.5em}
\end{table}

\begin{table}[h!]
\centering
\renewcommand{\arraystretch}{1.2}
\setlength{\tabcolsep}{0pt}
\begin{tabularx}{\textwidth}{|*{16}{>{\centering\arraybackslash\hsize=1.0\hsize}X|}}
\hline
\multicolumn{3}{|>{\centering\arraybackslash\hsize=3.0\hsize}X|}{\textbf{\laghu{ज}\laghu{प}\laghu{त}} \par (3)} &
\multicolumn{5}{>{\centering\arraybackslash\hsize=5.0\hsize}X|}{\textbf{\laghu{नि}\guru{रं}\laghu{त}\laghu{र}} \par (5)} &
\multicolumn{4}{>{\centering\arraybackslash\hsize=4.0\hsize}X|}{\textbf{\laghu{ह}\laghu{नु}\laghu{म}\laghu{त}} \par (4)} &
\multicolumn{4}{>{\centering\arraybackslash\hsize=4.0\hsize}X|}{\textbf{\guru{बी}\guru{रा}} \par (4)} \\
\hline
\end{tabularx}
\vspace{-1.5em}
\end{table}

\vspace{1em}
\textbf{ \deva{सरल-संस्कृतम्}:}\\
 \deva{सर्वे रोगाः नश्यन्ति, सर्वाः पीडाः च हरिताः भवन्ति}\\
 \deva{यः निरन्तरं हे वीर हनुमन्! भवतः नाम जपति॥}\\


All diseases are destroyed, and all suffering/pain is taken away\\
for one who continuously chants your name, O brave Hanuman!


\vspace{1em}
\newpage
\textbf{\deva{प्रतिपदार्थः} (Meaning of each word):}
\begin{table}[h!]
\centering
\renewcommand{\arraystretch}{1.2}
\begin{tabularx}{\textwidth}{@{} l l X X @{}}
\toprule
\textbf{Awadhi} & \textbf{Sanskrit} & \textbf{English} & \textbf{Grammar \& Notes} \\
\midrule
\deva{नासै} \gotchaicon / \telu{నాసై} / \textit{nāsai}  &  \deva{नश्यति}  &  Is destroyed  & Verb ending in `-ai` \\
\deva{रोग} / \telu{రోగ} / \textit{roga}  &  \deva{रोगाः}  &  Diseases  &   \\
\deva{हरै} \gotchaicon / \telu{హరై} / \textit{harai}  &  \deva{हरति}  &  Takes away / destroys  & Verb ending in `-ai` \\
\deva{सब पीरा} / \telu{సబ పీరా} / \textit{saba pīrā}  &  \deva{सर्वाः पीडाः}  &  All pain  & \deva{पीडा $\rightarrow$ पीरा} \\ 
\deva{जपत} / \telu{జపత} / \textit{japata}  &  \deva{जपन्}  &  Chanting  & Continuous action `-ta` \\
\deva{निरन्तर} / \telu{నిరంతర} / \textit{nirantara}  &  \deva{निरन्तरम्}  &  Continuously  &   \\
\deva{हनुमत} / \telu{హనుమత} / \textit{hanumata}  &  \deva{हनुमन्}  &  Hanuman  &   \\
\deva{बीरा} / \telu{బీరా} / \textit{bīrā}  &  \deva{वीरः}  &  Brave/Hero  & \par  \deva{} \\ 
\bottomrule
\end{tabularx}
\end{table}

\begin{tcolorbox}[gotchabox]
\begin{itemize}
    \item \textbf{The \deva{ड} $\rightarrow$ \deva{र} Shift:} The Sanskrit word for pain is \deva{पीडा} (Peedaa). Over time, the retroflex 'D' softens into a rolling 'R' sound, giving the beautiful Awadhi word \deva{पीरा} (Peeraa).
    \item \textbf{Verbs ending in -ai:} Just like \deva{गावैं} and \deva{लहै}, the verbs \deva{नासै} (Destroy) and \deva{हरै} (Remove) are standard present-tense conjugations in Awadhi poetry.
\end{itemize}
\end{tcolorbox}

\newpage


% ----- CHAUPAI 26 -----

\newpage

\subsection*{\deva{षड्विंशतितम चौपाई} (Twenty-Sixth Chaupai)}
\addcontentsline{toc}{subsection}{\deva{षड्विंशतितम चौपाई} (Twenty-Sixth Chaupai) — \deva{संकट तें...}}

\begin{minipage}[t]{0.55\textwidth}
\vspace{0pt}

{\large \linespread{1.5}\selectfont
\deva{संकट तें हनुमान छुड़ावै।}\\
\deva{मन क्रम बचन ध्यान जो लावै॥}
\par}

\vspace{0.5em}

{\large \linespread{1.5}\selectfont
\telu{సంకట తేఁ హనుమాన ఛుడావై।}\\
\telu{మన క్రమ బచన ధ్యాన జో లావై॥}
\par}

\vspace{0.5em}

{\large \linespread{1.3}\selectfont
\textit{saṅkaṭa tẽ hanumāna chuṛāvai |}\\
\textit{mana krama bacana dhyāna jo lāvai ||}
\par}
\end{minipage}%
\hfill
\begin{minipage}[t]{0.42\textwidth}
\vspace{0pt}
\begin{tcolorbox}[anchorbox]
\textbf{The Perfect Alignment:} When your mind (\textit{Mana}), words (\textit{Bachan}), and actions (\textit{Krama}) are perfectly aligned toward a single goal, you can navigate your way out of any crisis (\textit{Sankat}). Hypocrisy and distraction are the true enemies of success.
\vspace{0.5em}\newline He turned a terrifying crisis around by perfectly aligning his sharp mind with his decisive, practical actions.\newline\null\hfill\href{https://www.valmiki.iitk.ac.in/sloka?field_kanda_tid=5&language=dv&field_sarga_value=53}{\textit{Sundara Kanda, 53:25-27}}
\end{tcolorbox}
\end{minipage}

\vspace{0.5em}
\subsubsection*{\textbf{Visual Rhythm Map:}}
\begin{table}[h!]
\centering
\renewcommand{\arraystretch}{1.2}
\setlength{\tabcolsep}{0pt}
\begin{tabularx}{\textwidth}{|*{16}{>{\centering\arraybackslash\hsize=1.0\hsize}X|}}
\hline
\multicolumn{4}{|>{\centering\arraybackslash\hsize=4.0\hsize}X|}{\textbf{\guru{सं}\laghu{क}\laghu{ट}} \par (4)} &
\multicolumn{2}{>{\centering\arraybackslash\hsize=2.0\hsize}X|}{\textbf{\guru{तें}} \par (2)} &
\multicolumn{5}{>{\centering\arraybackslash\hsize=5.0\hsize}X|}{\textbf{\laghu{ह}\laghu{नु}\guru{मा}\laghu{न}} \par (5)} &
\multicolumn{5}{>{\centering\arraybackslash\hsize=5.0\hsize}X|}{\textbf{\laghu{छु}\guru{ड़ा}\guru{वै}} \par (5)} \\
\hline
\end{tabularx}
\vspace{-1.5em}
\end{table}

\begin{table}[h!]
\centering
\renewcommand{\arraystretch}{1.2}
\setlength{\tabcolsep}{0pt}
\begin{tabularx}{\textwidth}{|*{16}{>{\centering\arraybackslash\hsize=1.0\hsize}X|}}
\hline
\multicolumn{2}{|>{\centering\arraybackslash\hsize=2.0\hsize}X|}{\textbf{\laghu{म}\laghu{न}} \par (2)} &
\multicolumn{2}{>{\centering\arraybackslash\hsize=2.0\hsize}X|}{\textbf{\laghu{क्र}\laghu{म}} \par (2)} &
\multicolumn{3}{>{\centering\arraybackslash\hsize=3.0\hsize}X|}{\textbf{\laghu{ब}\laghu{च}\laghu{न}} \par (3)} &
\multicolumn{3}{>{\centering\arraybackslash\hsize=3.0\hsize}X|}{\textbf{\guru{ध्या}\laghu{न}} \par (3)} &
\multicolumn{2}{>{\centering\arraybackslash\hsize=2.0\hsize}X|}{\textbf{\guru{जो}} \par (2)} &
\multicolumn{4}{>{\centering\arraybackslash\hsize=4.0\hsize}X|}{\textbf{\guru{ला}\guru{वै}} \par (4)} \\
\hline
\end{tabularx}
\vspace{-1.5em}
\end{table}

\vspace{1em}
\textbf{ \deva{सरल-संस्कृतम्}:}\\
 \deva{सङ्कटात् तमेव हनुमान् मोचयति}\\
 \deva{यः मनसा, कर्मणा, वचनेन च तस्य ध्यानं करोति॥}\\


Hanuman rescues that person from crisis and danger\\
who meditates on him with mind, action, and words/speech.


\vspace{1em}
\newpage
\textbf{\deva{प्रतिपदार्थः} (Meaning of each word):}
\begin{table}[h!]
\centering
\renewcommand{\arraystretch}{1.2}
\begin{tabularx}{\textwidth}{@{} l l X X @{}}
\toprule
\textbf{Awadhi} & \textbf{Sanskrit} & \textbf{English} & \textbf{Grammar \& Notes} \\
\midrule
\deva{संकट तें} / \telu{సంకట తేఁ} / \textit{saṅkaṭa tẽ}  &  \deva{सङ्कटात्}  &  From crisis  & Ablative postposition `-tẽ` \\
\deva{हनुमान} / \telu{హనుమాన} / \textit{hanumāna}  &  \deva{हनुमान्}  &  Hanuman  &   \\
\deva{छुड़ावै} / \telu{ఛుడావై} / \textit{chuṛāvai}  &  \deva{मोचयति}  &  Rescues/Frees  &   \\
\deva{मन} / \telu{మన} / \textit{mana}  &  \deva{मनसा}  &  With mind  &   \\
\deva{क्रम} / \telu{క్రమ} / \textit{krama}  &  \deva{कर्मणा}  &  With actions  & Metathesis of Karma \\
\deva{बचन} / \telu{బచన} / \textit{bacana}  &  \deva{वचनेन}  &  With words  & \\ 
\deva{ध्यान} / \telu{ధ్యాన} / \textit{dhyāna}  &  \deva{ध्यानम्}  &  Meditation  &   \\
\deva{जो लावै} / \telu{జో లావై} / \textit{jo lāvai}  &  \deva{यः करोति (आनयति)}  &  Whoever applies  &   \\
\bottomrule
\end{tabularx}
\end{table}



\newpage


% ----- CHAUPAI 27 -----

\newpage

\subsection*{\deva{सप्तविंशतितम चौपाई} (Twenty-Seventh Chaupai)}
\addcontentsline{toc}{subsection}{\deva{सप्तविंशतितम चौपाई} (Twenty-Seventh Chaupai) — \deva{सब पर...}}

\begin{minipage}[t]{0.55\textwidth}
\vspace{0pt}

{\large \linespread{1.5}\selectfont
\deva{सब पर राम तपस्वी राजा।}\\
\deva{तिन के काज सकल तुम साजा॥}
\par}

\vspace{0.5em}

{\large \linespread{1.5}\selectfont
\telu{సబ పర రామ తపస్వీ రాజా।}\\
\telu{తిన కే కాజ సకల తుమ సాజా॥}
\par}

\vspace{0.5em}

{\large \linespread{1.3}\selectfont
\textit{saba para rāma tapasvī rājā |}\\
\textit{tina ke kāja sakala tuma sājā ||}
\par}
\end{minipage}%
\hfill
\begin{minipage}[t]{0.42\textwidth}
\vspace{0pt}
\begin{tcolorbox}[anchorbox]
\textbf{Executing the Vision:} Even the greatest kings and visionaries rely on dedicated individuals to execute the hardest tasks on the ground. A grand vision or a brilliant idea is meaningless without a hard-working team willing to put in the effort to make it a reality.
\vspace{0.5em}\newline Even the greatest visionary leaders required his relentless, on-the-ground execution to turn their goals into reality.\newline\null\hfill\href{https://www.valmiki.iitk.ac.in/sloka?field_kanda_tid=5&language=dv&field_sarga_value=65}{\textit{Sundara Kanda, 65:26}}
\end{tcolorbox}
\end{minipage}

\vspace{0.5em}
\subsubsection*{\textbf{Visual Rhythm Map:}}
\begin{table}[h!]
\centering
\renewcommand{\arraystretch}{1.2}
\setlength{\tabcolsep}{0pt}
\begin{tabularx}{\textwidth}{|*{16}{>{\centering\arraybackslash\hsize=1.0\hsize}X|}}
\hline
\multicolumn{2}{|>{\centering\arraybackslash\hsize=2.0\hsize}X|}{\textbf{\laghu{स}\laghu{ब}} \par (2)} &
\multicolumn{2}{>{\centering\arraybackslash\hsize=2.0\hsize}X|}{\textbf{\laghu{प}\laghu{र}} \par (2)} &
\multicolumn{3}{>{\centering\arraybackslash\hsize=3.0\hsize}X|}{\textbf{\guru{रा}\laghu{म}} \par (3)} &
\multicolumn{5}{>{\centering\arraybackslash\hsize=5.0\hsize}X|}{\textbf{\laghu{त}\guru{प}\guru{स्वी}} \par (5)} &
\multicolumn{4}{>{\centering\arraybackslash\hsize=4.0\hsize}X|}{\textbf{\guru{रा}\guru{जा}} \par (4)} \\
\hline
\end{tabularx}
\vspace{-1.5em}
\end{table}

\begin{table}[h!]
\centering
\renewcommand{\arraystretch}{1.2}
\setlength{\tabcolsep}{0pt}
\begin{tabularx}{\textwidth}{|*{16}{>{\centering\arraybackslash\hsize=1.0\hsize}X|}}
\hline
\multicolumn{2}{|>{\centering\arraybackslash\hsize=2.0\hsize}X|}{\textbf{\laghu{ति}\laghu{न}} \par (2)} &
\multicolumn{2}{>{\centering\arraybackslash\hsize=2.0\hsize}X|}{\textbf{\guru{के}} \par (2)} &
\multicolumn{3}{>{\centering\arraybackslash\hsize=3.0\hsize}X|}{\textbf{\guru{का}\laghu{ज}} \par (3)} &
\multicolumn{3}{>{\centering\arraybackslash\hsize=3.0\hsize}X|}{\textbf{\laghu{स}\laghu{क}\laghu{ल}} \par (3)} &
\multicolumn{2}{>{\centering\arraybackslash\hsize=2.0\hsize}X|}{\textbf{\laghu{तु}\laghu{म}} \par (2)} &
\multicolumn{4}{>{\centering\arraybackslash\hsize=4.0\hsize}X|}{\textbf{\guru{सा}\guru{जा}} \par (4)} \\
\hline
\end{tabularx}
\vspace{-1.5em}
\end{table}

\vspace{1em}
\textbf{ \deva{सरल-संस्कृतम्}:}\\
 \deva{सर्वेषामुपरि तपस्वी राजा श्रीरामः अस्ति।}\\
 \deva{तस्य अपि सर्वाणि कार्याणि भवान् एव सिद्धानि कृतवान् (साधितवान्)॥}\\


Above all beings is Lord Rama, the ascetic King.\\
Yet, you are the one who accomplished all of even His tasks!


\vspace{1em}
\newpage
\textbf{\deva{प्रतिपदार्थः} (Meaning of each word):}
\begin{table}[h!]
\centering
\renewcommand{\arraystretch}{1.2}
\begin{tabularx}{\textwidth}{@{} l l X X @{}}
\toprule
\textbf{Awadhi} & \textbf{Sanskrit} & \textbf{English} & \textbf{Grammar \& Notes} \\
\midrule
\deva{सब पर} / \telu{సబ పర} / \textit{saba para}  &  \deva{सर्वेषाम् उपरि}  &  Above everyone  &   \\
\deva{राम} / \telu{రామ} / \textit{rāma}  &  \deva{रामः}  &  Rama  &   \\
\deva{तपस्वी राजा} / \telu{తపస్వీ రాజా} / \textit{tapasvī rājā}  &  \deva{तपस्वी राजा}  &  Ascetic King  & Rama ruled as an ascetic \\
\deva{तिन के} / \telu{తిన కే} / \textit{tina ke}  &  \deva{तस्य अपि}  &  Even His  & Honorific plural genitive \\
\deva{काज} / \telu{కాజ} / \textit{kāja}  &  \deva{कार्याणि}  &  Tasks / Work  & (from Kaarya) \\ 
\deva{सकल} / \telu{సకల} / \textit{sakala}  &  \deva{सकलानि / सर्वाणि}  &  All / Entire  &   \\
\deva{तुम} / \telu{తుమ} / \textit{tuma}  &  \deva{त्वम्}  &  You  &   \\
\deva{साजा} / \telu{సాజా} / \textit{sājā}  &  \deva{साधितवान्}  &  Accomplished / Arranged  &   \\
\bottomrule
\end{tabularx}
\end{table}



\newpage


% ----- CHAUPAI 28 -----

\newpage

\subsection*{\deva{अष्टाविंशतितम चौपाई} (Twenty-Eighth Chaupai)}
\addcontentsline{toc}{subsection}{\deva{अष्टाविंशतितम चौपाई} (Twenty-Eighth Chaupai) — \deva{और मनोरथ...}}

\begin{minipage}[t]{0.55\textwidth}
\vspace{0pt}

{\large \linespread{1.5}\selectfont
\deva{और मनोरथ जो कोइ लावै।}\\
\deva{सोइ अमित जीवन फल पावै॥}
\par}

\vspace{0.5em}

{\large \linespread{1.5}\selectfont
\telu{ఔర మనోరథ జో కోఇ లావై।}\\
\telu{సోఇ అమిత జీవన ఫల పావై॥}
\par}

\vspace{0.5em}

{\large \linespread{1.3}\selectfont
\textit{aura manoratha jo koi lāvai |}\\
\textit{soi amita jīvana phala pāvai ||}
\par}
\end{minipage}%
\hfill
\begin{minipage}[t]{0.42\textwidth}
\vspace{0pt}
\begin{tcolorbox}[anchorbox]
\textbf{Meaningful Ambitions:} True success isn't just about wishing for petty material things; it is about developing the clarity to seek the \textit{Amit Jeevan Phal}—the limitless fruit of a deeply purposeful, well-lived, and meaningful life.
\vspace{0.5em}\newline He accepted the ultimate reward only because his truest ambition was meaningful dedication, not material wealth.\newline\null\hfill\href{https://www.valmiki.iitk.ac.in/sloka?field_kanda_tid=6&language=dv&field_sarga_value=128}{\textit{Yuddha Kanda, 128:82-83}}
\end{tcolorbox}
\end{minipage}

\vspace{0.5em}
\subsubsection*{\textbf{Visual Rhythm Map:}}
\begin{table}[h!]
\centering
\renewcommand{\arraystretch}{1.2}
\noindent\makebox[\textwidth][l]{%
  {%
  \setlength{\tabcolsep}{0pt}%
  \begin{tabularx}{1.125\textwidth}{|*{18}{>{\centering\arraybackslash\hsize=1.0\hsize}X|}}
  \hline
  \multicolumn{3}{|>{\centering\arraybackslash\hsize=3.0\hsize}X|}{\textbf{\guru{औ}\laghu{र}} \par (3)} &
  \multicolumn{5}{>{\centering\arraybackslash\hsize=5.0\hsize}X|}{\textbf{\laghu{म}\guru{नो}\laghu{र}\laghu{थ}} \par (5)} &
  \multicolumn{2}{>{\centering\arraybackslash\hsize=2.0\hsize}X|}{\textbf{\guru{जो}} \par (2)} &
  \multicolumn{3}{>{\centering\arraybackslash\hsize=3.0\hsize}X|}{\textbf{\guru{को}\laghu{इ}} \par (3)} &
  \multicolumn{4}{>{\centering\arraybackslash\hsize=4.0\hsize}X|}{\textbf{\guru{ला}\guru{वै}} \par (4)} \\
  \hline
  \end{tabularx}%
  }%
  \hspace{0.01\textwidth}%
  \begin{minipage}{0.1\textwidth}
  \vspace{-0.5em}
  \centering
  {\Huge \textbf{17}}
  \end{minipage}%
}
\vspace{-1.5em}
\end{table}

\begin{table}[h!]
\centering
\renewcommand{\arraystretch}{1.2}
\setlength{\tabcolsep}{0pt}
\begin{tabularx}{\textwidth}{|*{16}{>{\centering\arraybackslash\hsize=1.0\hsize}X|}}
\hline
\multicolumn{3}{|>{\centering\arraybackslash\hsize=3.0\hsize}X|}{\textbf{\guru{सो}\laghu{इ}} \par (3)} &
\multicolumn{3}{>{\centering\arraybackslash\hsize=3.0\hsize}X|}{\textbf{\laghu{अ}\laghu{मि}\laghu{त}} \par (3)} &
\multicolumn{4}{>{\centering\arraybackslash\hsize=4.0\hsize}X|}{\textbf{\guru{जी}\laghu{व}\laghu{न}} \par (4)} &
\multicolumn{2}{>{\centering\arraybackslash\hsize=2.0\hsize}X|}{\textbf{\laghu{फ}\laghu{ल}} \par (2)} &
\multicolumn{4}{>{\centering\arraybackslash\hsize=4.0\hsize}X|}{\textbf{\guru{पा}\guru{वै}} \par (4)} \\
\hline
\end{tabularx}
\vspace{-1.5em}
\end{table}

\vspace{1em}
\textbf{ \deva{सरल-संस्कृतम्}:}\\
 \deva{यः कश्चित् अन्यम् अपि मनोरथम् (इच्छां) आनयति}\\
 \deva{सः जीवने अमितम् (असीमं) फलं प्राप्नोति॥}\\


Whoever brings any other desire or wish to you\\
that person obtains limitless fruits in their life.


\vspace{1em}
\newpage
\textbf{\deva{प्रतिपदार्थः} (Meaning of each word):}
\begin{table}[h!]
\centering
\renewcommand{\arraystretch}{1.2}
\begin{tabularx}{\textwidth}{@{} l l X X @{}}
\toprule
\textbf{Awadhi} & \textbf{Sanskrit} & \textbf{English} & \textbf{Grammar \& Notes} \\
\midrule
\deva{और} / \telu{ఔర}  &  \deva{अन्य / अपरः}  &  And / Other  &   \\
\deva{मनोरथ} / \telu{మనోరథ}  &  \deva{मनोरथम्}  &  Desire / Wish  &   \\
\deva{जो कोइ} / \telu{జో కోఇ}  &  \deva{यः कश्चित्}  &  Whoever  &   \\
\deva{लावै} / \telu{లావై}  &  \deva{आनयति}  &  Brings  &   \\
\deva{सोइ} / \telu{సోఇ}  &  \deva{सः एव}  &  That very person  & Awadhi pronoun \\
\deva{अमित} / \telu{అమిత}  &  \deva{अमितम् / असीमम्}  &  Limitless / Infinite  &   \\
\deva{जीवन फल} / \telu{జీవన ఫల}  &  \deva{जीवन-फलम्}  &  Fruit of life  &   \\
\deva{पावै} / \telu{పావై}  &  \deva{प्राप्नोति}  &  Obtains  &   \\
\bottomrule
\end{tabularx}
\end{table}

\begin{tcolorbox}[poetrybox]
\textbf{The 18-Beat Flow \& Bhava-Pradhan:}\\
Line 1 mathematically adds up to 17 beats. However, keeping in mind that this is a \textit{Bhava Pradhan} poem, the 16-matra rule is a convenience, not a strict law like in Sanskrit prosody. 
Singers joyfully and instinctively compress \deva{कोई} into a quick \textit{koy} without interrupting the singing flow, 
making this inconsistency a beautiful mnemonic quirk of the oral tradition!
\end{tcolorbox}

\newpage


% ----- CHAUPAI 29 -----

\newpage

\subsection*{\deva{एकोनत्रिंशत्तम चौपाई} (Twenty-Ninth Chaupai)}
\addcontentsline{toc}{subsection}{\deva{एकोनत्रिंशत्तम चौपाई} (Twenty-Ninth Chaupai) — \deva{चारों जुग...}}

\begin{minipage}[t]{0.55\textwidth}
\vspace{0pt}

{\large \linespread{1.5}\selectfont
\deva{चारों जुग परताप तुम्हारा।}\\
\deva{है परसिद्ध जगत उजियारा॥}
\par}

\vspace{0.5em}

{\large \linespread{1.5}\selectfont
\telu{చారోఁ జుగ పరతాప తుమ్హారా।}\\
\telu{హై పరసిద్ధ జగత ఉజియారా॥}
\par}

\vspace{0.5em}

{\large \linespread{1.3}\selectfont
\textit{cārõ juga paratāpa tumhārā |}\\
\textit{hai parasiddha jagata ujiyārā ||}
\par}
\end{minipage}%
\hfill
\begin{minipage}[t]{0.42\textwidth}
\vspace{0pt}
\begin{tcolorbox}[anchorbox]
\textbf{Timeless Principles:} His glory shines in all four epochs (\textit{Charon Yug}). In the *Valmiki Ramayana* (*Sundara Kanda, 1:198*), the celestials praise Hanuman, declaring that anyone who possesses his four timeless qualities—courage, vision, intellect, and skill—will never fail in their endeavors. These core principles remain globally relevant across all generations.
\vspace{0.5em}\newline Hanuman's four foundational leadership qualities guarantee success that remains relevant across all space and time.\newline\null\hfill\href{https://www.valmiki.iitk.ac.in/sloka?field_kanda_tid=5&language=dv&field_sarga_value=1}{\textit{Sundara Kanda, 1:198}}
\end{tcolorbox}
\end{minipage}

\vspace{0.5em}
\subsubsection*{\textbf{Visual Rhythm Map:}}
\begin{table}[h!]
\centering
\renewcommand{\arraystretch}{1.2}
\setlength{\tabcolsep}{0pt}
\begin{tabularx}{\textwidth}{|*{16}{>{\centering\arraybackslash\hsize=1.0\hsize}X|}}
\hline
\multicolumn{4}{|>{\centering\arraybackslash\hsize=4.0\hsize}X|}{\textbf{\guru{चा}\guru{रों}} \par (4)} &
\multicolumn{2}{>{\centering\arraybackslash\hsize=2.0\hsize}X|}{\textbf{\laghu{जु}\laghu{ग}} \par (2)} &
\multicolumn{5}{>{\centering\arraybackslash\hsize=5.0\hsize}X|}{\textbf{\laghu{प}\laghu{र}\guru{ता}\laghu{प}} \par (5)} &
\multicolumn{5}{>{\centering\arraybackslash\hsize=5.0\hsize}X|}{\textbf{\laghu{तु}\guru{म्हा}\guru{रा}} \par (5)} \\
\hline
\end{tabularx}
\vspace{-1.5em}
\end{table}

\begin{table}[h!]
\centering
\renewcommand{\arraystretch}{1.2}
\setlength{\tabcolsep}{0pt}
\begin{tabularx}{\textwidth}{|*{16}{>{\centering\arraybackslash\hsize=1.0\hsize}X|}}
\hline
\multicolumn{2}{|>{\centering\arraybackslash\hsize=2.0\hsize}X|}{\textbf{\guru{है}} \par (2)} &
\multicolumn{5}{>{\centering\arraybackslash\hsize=5.0\hsize}X|}{\textbf{\laghu{प}\laghu{र}\guru{सि}\laghu{द्ध}} \par (5)} &
\multicolumn{3}{>{\centering\arraybackslash\hsize=3.0\hsize}X|}{\textbf{\laghu{ज}\laghu{ग}\laghu{त}} \par (3)} &
\multicolumn{6}{>{\centering\arraybackslash\hsize=6.0\hsize}X|}{\textbf{\laghu{उ}\laghu{जि}\guru{या}\guru{रा}} \par (6)} \\
\hline
\end{tabularx}
\vspace{-1.5em}
\end{table}

\vspace{1em}
\textbf{ \deva{सरल-संस्कृतम्}:}\\
 \deva{चतुर्षु युगेषु भवतः प्रतापः (महिमा) अस्ति।}\\
 \deva{संपूर्ण-जगति भवतः प्रकाशः प्रसिद्धः अस्ति॥}\\


Your glory and majesty exist throughout the four Yugas / ages.\\
Your illuminating light is famous across the universe.


\vspace{1em}
\newpage
\textbf{\deva{प्रतिपदार्थः} (Meaning of each word):}
\begin{table}[h!]
\centering
\renewcommand{\arraystretch}{1.2}
\begin{tabularx}{\textwidth}{@{} l l X X @{}}
\toprule
\textbf{Awadhi} & \textbf{Sanskrit} & \textbf{English} & \textbf{Grammar \& Notes} \\
\midrule
\deva{चारों जुग} \glossaryicon / \telu{చారోం జుగ} / \textit{cārõ juga}  &  \deva{चतुर्षु युगेषु}  &  In the four Yugas  & \\ 
\deva{परताप} \gotchaicon / \telu{పరతాప} / \textit{paratāpa}  &  \deva{प्रतापः}  &  Glory / Majesty  & Swara-Bhakti \\
\deva{तुम्हारा} / \telu{తుమ్హారా} / \textit{tumhārā}  &  \deva{भवतः}  &  Your  &   \\
\deva{है} \gotchaicon / \telu{హై} / \textit{hai}  &  \deva{अस्ति}  &  Is  &   \\
\deva{परसिद्ध} \gotchaicon / \telu{పరసిద్ధ} / \textit{parasiddha}  &  \deva{प्रसिद्धः}  &  Famous / Known  & Swara-Bhakti \\
\deva{जगत} / \telu{జగత} / \textit{jagata}  &  \deva{जगति}  &  In the world  &   \\
\deva{उजियारा} / \telu{ఉజియారా} / \textit{ujiyārā}  &  \deva{प्रकाशः / उज्ज्वलः}  &  Light / Radiance  &   \\
\bottomrule
\end{tabularx}
\end{table}

\begin{tcolorbox}[gotchabox]
\begin{itemize}
    \item \textbf{Swara-Bhakti in Action:} Both \deva{परताप} (Parataap) and \deva{परसिद्ध} (Parasiddha) exhibit pure Awadhi Swara-Bhakti (anaptyxis). The strict Sanskrit conjuncts \deva{प्रताप} (Prataapa) and \deva{प्रसिद्ध} (Prasiddha) are broken open by inserting a short `a` vowel between $P$ and $R$. This makes the words musically softer and easier to sing while perfectly occupying the required musical beats.
\end{itemize}
\end{tcolorbox}

\newpage


% ----- CHAUPAI 30 -----

\newpage

\subsection*{\deva{त्रिंशत्तम चौपाई} (Thirtieth Chaupai)}
\addcontentsline{toc}{subsection}{\deva{त्रिंशत्तम चौपाई} (Thirtieth Chaupai) — \deva{साधु सन्त...}}

\begin{minipage}[t]{0.55\textwidth}
\vspace{0pt}

{\large \linespread{1.5}\selectfont
\deva{साधु सन्त के तुम रखवारे।}\\
\deva{असुर निकंदन राम दुलारे॥}
\par}

\vspace{0.5em}

{\large \linespread{1.5}\selectfont
\telu{సాధు సన్త కే తుమ రఖవారే।}\\
\telu{అసుర నికందన రామ దులారే॥}
\par}

\vspace{0.5em}

{\large \linespread{1.3}\selectfont
\textit{sādhu santa ke tuma rakhavāre |}\\
\textit{asura nikandana rāma dulāre ||}
\par}
\end{minipage}%
\hfill
\begin{minipage}[t]{0.42\textwidth}
\vspace{0pt}
\begin{tcolorbox}[anchorbox]
\textbf{The Righteous Protector:} He is the gentle guardian of the innocent (\textit{Sadhu}) and the fierce uprooter of toxic elements (\textit{Asura Nikandana}). Students must learn to cultivate both immense compassion for others and the firm boundaries needed to stand up against negative influences.
\vspace{0.5em}\newline He struck the perfect balance: fiercely eliminating toxic threats while gently protecting the innocent in his care.\newline\null\hfill\href{https://www.valmiki.iitk.ac.in/sloka?field_kanda_tid=5&language=dv&field_sarga_value=46}{\textit{Sundara Kanda, 46:25-27}}
\end{tcolorbox}
\end{minipage}

\vspace{0.5em}
\subsubsection*{\textbf{Visual Rhythm Map:}}
\begin{table}[h!]
\centering
\renewcommand{\arraystretch}{1.2}
\setlength{\tabcolsep}{0pt}
\begin{tabularx}{\textwidth}{|*{16}{>{\centering\arraybackslash\hsize=1.0\hsize}X|}}
\hline
\multicolumn{3}{|>{\centering\arraybackslash\hsize=3.0\hsize}X|}{\textbf{\guru{सा}\laghu{धु}} \par (3)} &
\multicolumn{3}{>{\centering\arraybackslash\hsize=3.0\hsize}X|}{\textbf{\guru{स}\laghu{न्त}} \par (3)} &
\multicolumn{2}{>{\centering\arraybackslash\hsize=2.0\hsize}X|}{\textbf{\guru{के}} \par (2)} &
\multicolumn{2}{>{\centering\arraybackslash\hsize=2.0\hsize}X|}{\textbf{\laghu{तु}\laghu{म}} \par (2)} &
\multicolumn{6}{>{\centering\arraybackslash\hsize=6.0\hsize}X|}{\textbf{\laghu{र}\laghu{ख}\guru{वा}\guru{रे}} \par (6)} \\
\hline
\end{tabularx}
\vspace{-1.5em}
\end{table}

\begin{table}[h!]
\centering
\renewcommand{\arraystretch}{1.2}
\setlength{\tabcolsep}{0pt}
\begin{tabularx}{\textwidth}{|*{16}{>{\centering\arraybackslash\hsize=1.0\hsize}X|}}
\hline
\multicolumn{3}{|>{\centering\arraybackslash\hsize=3.0\hsize}X|}{\textbf{\laghu{अ}\laghu{सु}\laghu{र}} \par (3)} &
\multicolumn{5}{>{\centering\arraybackslash\hsize=5.0\hsize}X|}{\textbf{\laghu{नि}\guru{कं}\laghu{द}\laghu{न}} \par (5)} &
\multicolumn{3}{>{\centering\arraybackslash\hsize=3.0\hsize}X|}{\textbf{\guru{रा}\laghu{म}} \par (3)} &
\multicolumn{5}{>{\centering\arraybackslash\hsize=5.0\hsize}X|}{\textbf{\laghu{दु}\guru{ला}\guru{रे}} \par (5)} \\
\hline
\end{tabularx}
\vspace{-1.5em}
\end{table}

\vspace{1em}
\textbf{ \deva{सरल-संस्कृतम्}:}\\
 \deva{साधूनां सन्तानां च भवान् रक्षकः अस्ति।}\\
 \deva{भवान् असुराणां विनाशकः, श्रीरामस्य प्रियः च अस्ति॥}\\


You are the protector of sages and saints.\\
You are the destroyer of demons, and the darling of Lord Rama.


\vspace{1em}
\newpage
\textbf{\deva{प्रतिपदार्थः} (Meaning of each word):}
\begin{table}[h!]
\centering
\renewcommand{\arraystretch}{1.2}
\begin{tabularx}{\textwidth}{@{} l l X X @{}}
\toprule
\textbf{Awadhi} & \textbf{Sanskrit} & \textbf{English} & \textbf{Grammar \& Notes} \\
\midrule
\deva{साधु सन्त के} / \telu{సాధు సన్త కే}  &  \deva{साधु-सन्तानाम्}  &  Of sages and saints  &   \\
\deva{तुम} / \telu{తుమ}  &  \deva{त्वम्}  &  You  &   \\
\deva{रखवारे} / \telu{రఖవారే}  &  \deva{रक्षकः}  &  Protector  & \\ 
\deva{असुर} / \telu{అసుర}  &  \deva{असुराणाम्}  &  Of demons  &   \\
\deva{निकंदन} / \telu{నికందన}  &  \deva{विनाशकः / निकृन्दन}  &  Destroyer / Uprooter  &   \\
\deva{राम दुलारे} / \telu{రామ దులారే}  &  \deva{राम-प्रियः}  &  Dear to Rama / Darling  &   \\
\bottomrule
\end{tabularx}
\end{table}

\newpage


% ----- CHAUPAI 31 -----

\newpage

\subsection*{\deva{एकत्रिंशत्तम चौपाई} (Thirty-First Chaupai)}
\addcontentsline{toc}{subsection}{\deva{एकत्रिंशत्तम चौपाई} (Thirty-First Chaupai) — \deva{अष्ट सिद्धि...}}

\begin{minipage}[t]{0.55\textwidth}
\vspace{0pt}

{\large \linespread{1.5}\selectfont
\deva{अष्ट सिद्धि नौ निधि के दाता।}\\
\deva{अस बर दीन जानकी माता॥}
\par}

\vspace{0.5em}

{\large \linespread{1.5}\selectfont
\telu{అష్ట సిద్ధి నౌ నిధి కే దాతా।}\\
\telu{అస బర దీన జానకీ మాతా॥}
\par}

\vspace{0.5em}

{\large \linespread{1.3}\selectfont
\textit{aṣṭa siddhi nau nidhi ke dātā |}\\
\textit{asa bara dīna jānakī mātā ||}
\par}
\end{minipage}%
\hfill
\begin{minipage}[t]{0.42\textwidth}
\vspace{0pt}
\begin{tcolorbox}[anchorbox]
\textbf{The Truest Wealth:} The greatest treasure (\textit{Nidhi}) or power (\textit{Siddhi}) in life is not magical abilities or massive bank accounts. True wealth is finding your unique purpose, achieving self-mastery, and building a community of friends and mentors who value you.
\vspace{0.5em}\newline The greatest power or treasure anyone can receive is discovering their true purpose and a community that values them.\newline\null\hfill\href{https://www.valmiki.iitk.ac.in/sloka?field_kanda_tid=5&language=dv&field_sarga_value=68}{\textit{Sundara Kanda, 68:26-29}}
\end{tcolorbox}
\end{minipage}

\vspace{0.5em}
\subsubsection*{\textbf{Visual Rhythm Map:}}
\begin{table}[h!]
\centering
\renewcommand{\arraystretch}{1.2}
\setlength{\tabcolsep}{0pt}
\begin{tabularx}{\textwidth}{|*{16}{>{\centering\arraybackslash\hsize=1.0\hsize}X|}}
\hline
\multicolumn{3}{|>{\centering\arraybackslash\hsize=3.0\hsize}X|}{\textbf{\guru{अ}\laghu{ष्ट}} \par (3)} &
\multicolumn{3}{>{\centering\arraybackslash\hsize=3.0\hsize}X|}{\textbf{\guru{सि}\laghu{द्धि}} \par (3)} &
\multicolumn{2}{>{\centering\arraybackslash\hsize=2.0\hsize}X|}{\textbf{\guru{नौ}} \par (2)} &
\multicolumn{2}{>{\centering\arraybackslash\hsize=2.0\hsize}X|}{\textbf{\laghu{नि}\laghu{धि}} \par (2)} &
\multicolumn{2}{>{\centering\arraybackslash\hsize=2.0\hsize}X|}{\textbf{\guru{के}} \par (2)} &
\multicolumn{4}{>{\centering\arraybackslash\hsize=4.0\hsize}X|}{\textbf{\guru{दा}\guru{ता}} \par (4)} \\
\hline
\end{tabularx}
\vspace{-1.5em}
\end{table}

\begin{table}[h!]
\centering
\renewcommand{\arraystretch}{1.2}
\setlength{\tabcolsep}{0pt}
\begin{tabularx}{\textwidth}{|*{16}{>{\centering\arraybackslash\hsize=1.0\hsize}X|}}
\hline
\multicolumn{2}{|>{\centering\arraybackslash\hsize=2.0\hsize}X|}{\textbf{\laghu{अ}\laghu{स}} \par (2)} &
\multicolumn{2}{>{\centering\arraybackslash\hsize=2.0\hsize}X|}{\textbf{\laghu{ब}\laghu{र}} \par (2)} &
\multicolumn{3}{>{\centering\arraybackslash\hsize=3.0\hsize}X|}{\textbf{\guru{दी}\laghu{न}} \par (3)} &
\multicolumn{5}{>{\centering\arraybackslash\hsize=5.0\hsize}X|}{\textbf{\guru{जा}\laghu{न}\guru{की}} \par (5)} &
\multicolumn{4}{>{\centering\arraybackslash\hsize=4.0\hsize}X|}{\textbf{\guru{मा}\guru{ता}} \par (4)} \\
\hline
\end{tabularx}
\vspace{-1.5em}
\end{table}

\vspace{1em}
\textbf{ \deva{सरल-संस्कृतम्}:}\\
 \deva{भवान् अष्ट-सिद्धीनां नव-निधीनां च दाता अस्ति।}\\
 \deva{एतादृशः वरः माता-जानक्या भवतः कृते दत्तः॥}\\


You are the bestower of the eight supernatural powers and nine divine treasures\\
Such was the boon given to you by Mother Janaki (Sita).


\vspace{1em}
\newpage
\textbf{\deva{प्रतिपदार्थः} (Meaning of each word):}
\begin{table}[h!]
\centering
\renewcommand{\arraystretch}{1.2}
\begin{tabularx}{\textwidth}{@{} l l X X @{}}
\toprule
\textbf{Awadhi} & \textbf{Sanskrit} & \textbf{English} & \textbf{Grammar \& Notes} \\
\midrule
\deva{अष्ट सिद्धि} \glossaryicon / \telu{అష్ట సిద్ధి} / \textit{aṣṭa siddhi}  &  \deva{अष्ट-सिद्धिः}  &  8 Supernatural Powers  & e.g., shrinking/expanding \\
\deva{नौ निधि} \glossaryicon / \telu{నౌ నిధి} / \textit{nau nidhi}  &  \deva{नव-निधिः}  &  9 Divine Treasures  & e.g., endless wealth \\
\deva{के दाता} / \telu{కే దాతా} / \textit{ke dātā}  &  \deva{प्रदाता}  &  Giver / Bestower of  &   \\
\deva{अस बर} / \telu{అస బర} / \textit{asa bara}  &  \deva{एतादृशं वरम्}  &  Such a boon  &   \\
\deva{दीन} / \telu{దీన} / \textit{dīna}  &  \deva{दत्तवती}  &  Gave  & Past tense verb \\
\deva{जानकी माता} / \telu{జానకీ మాతా} / \textit{jānakī mātā}  &  \deva{माता जानकी}  &  Mother Janaki (Sita)  &   \\
\bottomrule
\end{tabularx}
\end{table}

\newpage



% ==========================================
\section*{Group 7: The Bestower of Grace}
\addcontentsline{toc}{section}{Group 7: The Bestower of Grace}
% ==========================================

% ----- CHAUPAI 32 -----
\subsection*{\deva{द्वात्रिंशत्तम चौपाई} (Thirty-Second Chaupai)}
\addcontentsline{toc}{subsection}{\deva{द्वात्रिंशत्तम चौपाई} (Thirty-Second Chaupai) :  \deva{राम रसायन...}}

\begin{minipage}[t]{0.55\textwidth}
\vspace{0pt}

{\large \linespread{1.5}\selectfont
\deva{राम रसायन तुम्हरे पासा।}\\
\deva{सदा रहो रघुपति के दासा॥}
\par}

\vspace{0.5em}

{\large \linespread{1.5}\selectfont
\telu{రామ రసాయన తుమ్హరే పాసా।}\\
\telu{సదా రహో రఘుపతి కే దాసా॥}
\par}

\vspace{0.5em}

{\large \linespread{1.3}\selectfont
\textit{rāma rasāyana tumhare pāsā |}\\
\textit{sadā raho raghupati ke dāsā ||}
\par}
\end{minipage}%
\hfill
\begin{minipage}[t]{0.42\textwidth}
\vspace{0pt}
\begin{tcolorbox}[anchorbox]
\textbf{The Servant's Heart:} By consciously choosing to remain a humble servant and lifelong learner (\textit{Daasa}) rather than demanding to be the master, he found the ultimate cure (\textit{Rasayana}) for ego and unrest. Humility is the greatest medicine for a peaceful mind.
\vspace{0.5em}\newline He discovered that the ultimate cure for ego and unrest is choosing to remain a humble, lifelong learner.\newline\null\hfill\href{https://www.valmiki.iitk.ac.in/sloka?field_kanda_tid=6&language=dv&field_sarga_value=128}{\textit{Yuddha Kanda, 128:104}}
\end{tcolorbox}
\end{minipage}

\vspace{0.5em}
\subsubsection*{\textbf{Visual Rhythm Map:}}
\begin{table}[h!]
\centering
\renewcommand{\arraystretch}{1.2}
\setlength{\tabcolsep}{0pt}
\begin{tabularx}{\textwidth}{|*{16}{>{\centering\arraybackslash\hsize=1.0\hsize}X|}}
\hline
\multicolumn{3}{|>{\centering\arraybackslash\hsize=3.0\hsize}X|}{\textbf{\guru{रा}\laghu{म}} \par (3)} &
\multicolumn{5}{>{\centering\arraybackslash\hsize=5.0\hsize}X|}{\textbf{\laghu{र}\guru{सा}\laghu{य}\laghu{न}} \par (5)} &
\multicolumn{4}{>{\centering\arraybackslash\hsize=4.0\hsize}X|}{\textbf{\laghu{तु}\laghu{म्ह}\guru{रे}} \par (4)} &
\multicolumn{4}{>{\centering\arraybackslash\hsize=4.0\hsize}X|}{\textbf{\guru{पा}\guru{सा}} \par (4)} \\
\hline
\end{tabularx}
\vspace{-1.5em}
\end{table}

\begin{table}[h!]
\centering
\renewcommand{\arraystretch}{1.2}
\setlength{\tabcolsep}{0pt}
\begin{tabularx}{\textwidth}{|*{16}{>{\centering\arraybackslash\hsize=1.0\hsize}X|}}
\hline
\multicolumn{3}{|>{\centering\arraybackslash\hsize=3.0\hsize}X|}{\textbf{\laghu{स}\guru{दा}} \par (3)} &
\multicolumn{3}{>{\centering\arraybackslash\hsize=3.0\hsize}X|}{\textbf{\laghu{र}\guru{हो}} \par (3)} &
\multicolumn{4}{>{\centering\arraybackslash\hsize=4.0\hsize}X|}{\textbf{\laghu{र}\laghu{घु}\laghu{प}\laghu{ति}} \par (4)} &
\multicolumn{2}{>{\centering\arraybackslash\hsize=2.0\hsize}X|}{\textbf{\guru{के}} \par (2)} &
\multicolumn{4}{>{\centering\arraybackslash\hsize=4.0\hsize}X|}{\textbf{\guru{दा}\guru{सा}} \par (4)} \\
\hline
\end{tabularx}
\vspace{-1.5em}
\end{table}

\vspace{1em}
\textbf{ \deva{सरल-संस्कृतम्}:}\\
 \deva{भवतः पार्श्वे "राम-नाम" रूपी रसायनम् (औषधम्) अस्ति।}\\
 \deva{भवान् सर्वदा रघुपतेः (श्रीरामस्य) दासः भूत्वा एव तिष्ठति॥}\\


You possess the ultimate elixir, the essence of Lord Rama.\\
May you always remain the devoted servant of Lord Rama.


\vspace{1em}
\newpage
\textbf{\deva{प्रतिपदार्थः} (Meaning of each word):}
\begin{table}[h!]
\centering
\renewcommand{\arraystretch}{1.2}
\begin{tabularx}{\textwidth}{@{} l l X X @{}}
\toprule
\textbf{Awadhi} & \textbf{Sanskrit} & \textbf{English} & \textbf{Grammar \& Notes} \\
\midrule
\deva{राम रसायन} / \telu{రామ రసాయన}  &  \deva{राम-रसायनम्}  &  Elixir of Rama  & Rasa (Essence) + Ayana (Path) \\
\deva{तुम्हरे पासा} / \telu{తుమ్హరే పాసా}  &  \deva{भवतः पार्श्वे}  &  In your possession  &   \\
\deva{सदा रहो} / \telu{సదా రహో}  &  \deva{सर्वदा भव / तिष्ठ}  &  Always remain  &   \\
\deva{रघुपति के} / \telu{రఘుపతి కే}  &  \deva{रघुपतेः}  &  Of Lord Rama  &   \\
\deva{दासा} / \telu{దాసా}  &  \deva{दासः / सेवकः}  &  Servant  &   \\
\bottomrule
\end{tabularx}
\end{table}



\newpage


% ----- CHAUPAI 33 -----

\newpage

\subsection*{\deva{त्रयस्त्रिंशत्तम चौपाई} (Thirty-Third Chaupai)}
\addcontentsline{toc}{subsection}{\deva{त्रयस्त्रिंशत्तम चौपाई} (Thirty-Third Chaupai) :   \deva{तुम्हरे भजन...}}

\begin{minipage}[t]{0.55\textwidth}
\vspace{0pt}

{\large \linespread{1.5}\selectfont
\deva{तुम्हरे भजन राम को पावै।}\\
\deva{जनम जनम के दुख बिसरावै॥}
\par}

\vspace{0.5em}

{\large \linespread{1.5}\selectfont
\telu{తుమ్హరే భజన రామ కో పావై।}\\
\telu{జనమ జనమ కే దుఖ బిసరావై॥}
\par}

\vspace{0.5em}

{\large \linespread{1.3}\selectfont
\textit{tumhare bhajana rāma ko pāvai |}\\
\textit{janama janama ke dukha bisarāvai ||}
\par}
\end{minipage}%
\hfill
\begin{minipage}[t]{0.42\textwidth}
\vspace{0pt}
\begin{tcolorbox}[anchorbox]
\textbf{The Clean Slate:} Re-centering our minds on positive role models and healthy habits helps wash away the baggage of our past mistakes and traumas (\textit{Janama Janama ke dukha}), allowing us to move forward with a clean slate and renewed focus.
\vspace{0.5em}\newline The sound of his positive, encouraging voice abruptly snapped someone out of deep depression and past trauma.\newline\null\hfill\href{https://www.valmiki.iitk.ac.in/sloka?field_kanda_tid=5&language=dv&field_sarga_value=31}{\textit{Sundara Kanda, 31:14-18}}
\end{tcolorbox}
\end{minipage}

\vspace{0.5em}
\subsubsection*{\textbf{Visual Rhythm Map:}}
\begin{table}[h!]
\centering
\renewcommand{\arraystretch}{1.2}
\setlength{\tabcolsep}{0pt}
\begin{tabularx}{\textwidth}{|*{16}{>{\centering\arraybackslash\hsize=1.0\hsize}X|}}
\hline
\multicolumn{4}{|>{\centering\arraybackslash\hsize=4.0\hsize}X|}{\textbf{\laghu{तु}\laghu{म्ह}\guru{रे}} \par (4)} &
\multicolumn{3}{>{\centering\arraybackslash\hsize=3.0\hsize}X|}{\textbf{\laghu{भ}\laghu{ज}\laghu{न}} \par (3)} &
\multicolumn{3}{>{\centering\arraybackslash\hsize=3.0\hsize}X|}{\textbf{\guru{रा}\laghu{म}} \par (3)} &
\multicolumn{2}{>{\centering\arraybackslash\hsize=2.0\hsize}X|}{\textbf{\guru{को}} \par (2)} &
\multicolumn{4}{>{\centering\arraybackslash\hsize=4.0\hsize}X|}{\textbf{\guru{पा}\guru{वै}} \par (4)} \\
\hline
\end{tabularx}
\vspace{-1.5em}
\end{table}

\begin{table}[h!]
\centering
\renewcommand{\arraystretch}{1.2}
\setlength{\tabcolsep}{0pt}
\begin{tabularx}{\textwidth}{|*{16}{>{\centering\arraybackslash\hsize=1.0\hsize}X|}}
\hline
\multicolumn{3}{|>{\centering\arraybackslash\hsize=3.0\hsize}X|}{\textbf{\laghu{ज}\laghu{न}\laghu{म}} \par (3)} &
\multicolumn{3}{>{\centering\arraybackslash\hsize=3.0\hsize}X|}{\textbf{\laghu{ज}\laghu{न}\laghu{म}} \par (3)} &
\multicolumn{2}{>{\centering\arraybackslash\hsize=2.0\hsize}X|}{\textbf{\guru{के}} \par (2)} &
\multicolumn{2}{>{\centering\arraybackslash\hsize=2.0\hsize}X|}{\textbf{\laghu{दु}\laghu{ख}} \par (2)} &
\multicolumn{6}{>{\centering\arraybackslash\hsize=6.0\hsize}X|}{\textbf{\laghu{बि}\laghu{स}\guru{रा}\guru{वै}} \par (6)} \\
\hline
\end{tabularx}
\vspace{-1.5em}
\end{table}

\vspace{1em}
\textbf{ \deva{सरल-संस्कृतम्}:}\\
 \deva{भवतः भजनेन भक्तः श्रीरामं प्राप्नोति।}\\
 \deva{तथा जन्म-जन्मान्तराणां दुःखानि विस्मरति॥}\\


By singing your hymns, one attains Lord Rama.\\
And forgets the sorrows of many lifetimes.


\vspace{1em}
\newpage
\textbf{\deva{प्रतिपदार्थः} (Meaning of each word):}
\begin{table}[h!]
\centering
\renewcommand{\arraystretch}{1.2}
\begin{tabularx}{\textwidth}{@{} l l X X @{}}
\toprule
\textbf{Awadhi} & \textbf{Sanskrit} & \textbf{English} & \textbf{Grammar \& Notes} \\
\midrule
\deva{तुम्हरे भजन} / \telu{తుమ్హరే భజన} / \textit{tumhare bhajana}  &  \deva{भवतः भजनेन}  &  By singing your hymns  &   \\
\deva{राम को पावै} / \telu{రామ కో పావై} / \textit{rāma ko pāvai}  &  \deva{रामं प्राप्नोति}  &  Attains Rama  &   \\
\deva{जनम जनम के} / \telu{జనమ జనమ కే} / \textit{janama janama ke}  &  \deva{जन्म-जन्मान्तरस्य}  &  Of lifetimes  & Swara-Bhakti (Janma $\rightarrow$ Janama) \\ 
\deva{दुख} / \telu{దుఖ} / \textit{dukha}  &  \deva{दुःखानि}  &  Sorrows / Pain  &   \\
\deva{बिसरावै} / \telu{బిసరావై} / \textit{bisarāvai}  &  \deva{विस्मरति}  &  Forgets / Washes away  & Causative ending \\
\bottomrule
\end{tabularx}
\end{table}

\newpage


% ----- CHAUPAI 34 -----

\newpage

\subsection*{\deva{चतुस्त्रिंशत्तम चौपाई} (Thirty-Fourth Chaupai)}
\addcontentsline{toc}{subsection}{\deva{चतुस्त्रिंशत्तम चौपाई} (Thirty-Fourth Chaupai) :   \deva{अन्त काल...}}

\begin{minipage}[t]{0.55\textwidth}
\vspace{0pt}

{\large \linespread{1.5}\selectfont
\deva{अंत काल रघुबर पुर जाई।}\\
\deva{जहाँ जन्म हरि-भक्त कहाई॥}
\par}

\vspace{0.5em}

{\large \linespread{1.5}\selectfont
\telu{అంత కాల రఘుబర పుర జాయీ।}\\
\telu{జహాఁ జన్మ హరి-భక్త కహాయీ॥}
\par}

\vspace{0.5em}

{\large \linespread{1.3}\selectfont
\textit{anta kāla raghubara pura jāī |}\\
\textit{jahā̃ janma hari-bhakta kahāī ||}
\par}
\end{minipage}%
\hfill
\begin{minipage}[t]{0.42\textwidth}
\vspace{0pt}
\begin{tcolorbox}[anchorbox]
\textbf{The Ultimate Destination:} By dedicating our lives to service and excellence, our legacy (\textit{Anta Kaal}) becomes immortal. In the *Valmiki Ramayana* (*Yuddha Kanda, 1:13*), Lord Rama laments that Hanuman's monumental, selfless service is so vast that no material reward could ever match it. Rama offers his own embrace as his entire wealth, establishing Hanuman's eternal spiritual legacy.
\vspace{0.5em}\newline Rama's ultimate embrace and praise in Yuddha Kanda established Hanuman's immortal, spiritual legacy above all material rewards.\newline\null\hfill\href{https://www.valmiki.iitk.ac.in/sloka?field_kanda_tid=6&language=dv&field_sarga_value=1}{\textit{Yuddha Kanda, 1:13}}
\end{tcolorbox}
\end{minipage}

\vspace{0.5em}
\subsubsection*{\textbf{Visual Rhythm Map:}}
\begin{table}[h!]
\centering
\renewcommand{\arraystretch}{1.2}
\setlength{\tabcolsep}{0pt}
\begin{tabularx}{\textwidth}{|*{16}{>{\centering\arraybackslash\hsize=1.0\hsize}X|}}
\hline
\multicolumn{3}{|>{\centering\arraybackslash\hsize=3.0\hsize}X|}{\textbf{\guru{अं}\laghu{त}} \par (3)} &
\multicolumn{3}{>{\centering\arraybackslash\hsize=3.0\hsize}X|}{\textbf{\guru{का}\laghu{ल}} \par (3)} &
\multicolumn{4}{>{\centering\arraybackslash\hsize=4.0\hsize}X|}{\textbf{\laghu{र}\laghu{घु}\laghu{ब}\laghu{र}} \par (4)} &
\multicolumn{2}{>{\centering\arraybackslash\hsize=2.0\hsize}X|}{\textbf{\laghu{पु}\laghu{र}} \par (2)} &
\multicolumn{4}{>{\centering\arraybackslash\hsize=4.0\hsize}X|}{\textbf{\guru{जा}\guru{ई}} \par (4)} \\
\hline
\end{tabularx}
\vspace{-1.5em}
\end{table}

\begin{table}[h!]
\centering
\renewcommand{\arraystretch}{1.2}
\setlength{\tabcolsep}{0pt}
\begin{tabularx}{\textwidth}{|*{16}{>{\centering\arraybackslash\hsize=1.0\hsize}X|}}
\hline
\multicolumn{3}{|>{\centering\arraybackslash\hsize=3.0\hsize}X|}{\textbf{\laghu{ज}\guru{हाँ}} \par (3)} &
\multicolumn{3}{>{\centering\arraybackslash\hsize=3.0\hsize}X|}{\textbf{\guru{ज}\laghu{न्म}} \par (3)} &
\multicolumn{5}{>{\centering\arraybackslash\hsize=5.0\hsize}X|}{\textbf{\laghu{ह}\laghu{रि}\guru{भ}\laghu{क्त}} \par (5)} &
\multicolumn{5}{>{\centering\arraybackslash\hsize=5.0\hsize}X|}{\textbf{\laghu{क}\guru{हा}\guru{ई}} \par (5)} \\
\hline
\end{tabularx}
\vspace{-1.5em}
\end{table}

\vspace{1em}
\textbf{ \deva{सरल-संस्कृतम्}:}\\
 \deva{सः अन्त्यकाले श्रीरामस्य धाम (वैकुण्ठं) गच्छति।}\\
 \deva{यत्र तस्य जन्म भवति चेत्, सः भगवतः भक्तः इति उच्यते॥}\\


At the end of time/life, one goes to the celestial abode of Lord Rama.\\
And wherever they might be born anew, they are renowned as a devotee of Hari/God.


\vspace{1em}
\newpage
\textbf{\deva{प्रतिपदार्थः} (Meaning of each word):}
\begin{table}[h!]
\centering
\renewcommand{\arraystretch}{1.2}
\begin{tabularx}{\textwidth}{@{} l l X X @{}}
\toprule
\textbf{Awadhi} & \textbf{Sanskrit} & \textbf{English} & \textbf{Grammar \& Notes} \\
\midrule
\deva{अन्त काल} / \telu{అన్త కాల} / \textit{anta kāla}  &  \deva{अन्त्य-काले}  &  At the end of life  &   \\
\deva{रघुबर पुर} / \telu{రఘుబర పుర} / \textit{raghubara pura}  &  \deva{श्रीराम-धाम}  &  Rama's celestial city  & \\ 
\deva{जाई} / \telu{జాయీ} / \textit{jāī}  &  \deva{गच्छति}  &  Goes / Reaches  &   \\
\deva{जहाँ} / \telu{జహాఁ} / \textit{jahā̃}  &  \deva{यत्र}  &  Where / Wherever  &   \\
\deva{जन्म} / \telu{జన్మ} / \textit{janma}  &  \deva{जन्म}  &  Birth  & Pure Sanskrit \\
\deva{हरिभक्त} / \telu{హరిభక్త} / \textit{hari-bhakta}  &  \deva{हरि-भक्तः}  &  Devotee of God  &   \\
\deva{कहाई} / \telu{కహాయీ} / \textit{kahāī}  &  \deva{कथ्यते / उच्यते}  &  Is called / renowned  & Passive voice \\
\bottomrule
\end{tabularx}
\end{table}

\begin{tcolorbox}[poetrybox]
\textbf{Meter Alert (Heavy Conjuncts):}\\
Pay attention to \deva{जन्म} (Janma) and \deva{हरिभक्त} (Haribhakta) in Line 2. Because of the `nm` and `kt` conjuncts, the preceding vowels (\deva{ज} and \deva{भ}) absorb an extra beat, making them heavy (Guru) to maintain the exact 16-beat rhythm.
\end{tcolorbox}

\newpage


% ----- CHAUPAI 35 -----

\newpage

\subsection*{\deva{पञ्चत्रिंशत्तम चौपाई} (Thirty-Fifth Chaupai)}
\addcontentsline{toc}{subsection}{\deva{पञ्चत्रिंशत्तम चौपाई} (Thirty-Fifth Chaupai) :   \deva{और देवता...}}

\begin{minipage}[t]{0.55\textwidth}
\vspace{0pt}

{\large \linespread{1.5}\selectfont
\deva{और देवता चित्त न धरई।}\\
\deva{हनुमत सेइ सर्ब सुख करई॥}
\par}

\vspace{0.5em}

{\large \linespread{1.5}\selectfont
\telu{ఔర దేవతా చిత్త న ధరఈ।}\\
\telu{హనుమత సేఇ సర్బ సుఖ కరఈ॥}
\par}

\vspace{0.5em}

{\large \linespread{1.3}\selectfont
\textit{aura devatā citta na dharaī |}\\
\textit{hanumata sei sarba sukha karaī ||}
\par}
\end{minipage}%
\hfill
\begin{minipage}[t]{0.42\textwidth}
\vspace{0pt}
\begin{tcolorbox}[anchorbox]
\textbf{One-Pointed Focus:} A magnifying glass only starts a fire if held completely still. True academic or career success demands absolute, laser-focused dedication to one primary goal (\textit{Ananya Bhakti}) rather than constantly bouncing between different distractions.
\vspace{0.5em}\newline He ruthlessly blocked out all other distractions, organizing massive efforts with absolute, single-minded focus on his goal.\newline\null\hfill\href{https://www.valmiki.iitk.ac.in/sloka?field_kanda_tid=4&language=dv&field_sarga_value=39}{\textit{Kishkindha Kanda, 39:10}}
\end{tcolorbox}
\end{minipage}

\vspace{0.5em}
\subsubsection*{\textbf{Visual Rhythm Map:}}
\begin{table}[h!]
\centering
\renewcommand{\arraystretch}{1.2}
\setlength{\tabcolsep}{0pt}
\begin{tabularx}{\textwidth}{|*{16}{>{\centering\arraybackslash\hsize=1.0\hsize}X|}}
\hline
\multicolumn{3}{|>{\centering\arraybackslash\hsize=3.0\hsize}X|}{\textbf{\guru{औ}\laghu{र}} \par (3)} &
\multicolumn{5}{>{\centering\arraybackslash\hsize=5.0\hsize}X|}{\textbf{\guru{दे}\laghu{व}\guru{ता}} \par (5)} &
\multicolumn{3}{>{\centering\arraybackslash\hsize=3.0\hsize}X|}{\textbf{\guru{चि}\laghu{त्त}} \par (3)} &
\multicolumn{1}{>{\centering\arraybackslash\hsize=1.0\hsize}X|}{\textbf{\laghu{न}} \par (1)} &
\multicolumn{4}{>{\centering\arraybackslash\hsize=4.0\hsize}X|}{\textbf{\laghu{ध}\laghu{र}\guru{ई}} \par (4)} \\
\hline
\end{tabularx}
\vspace{-1.5em}
\end{table}

\begin{table}[h!]
\centering
\renewcommand{\arraystretch}{1.2}
\setlength{\tabcolsep}{0pt}
\begin{tabularx}{\textwidth}{|*{16}{>{\centering\arraybackslash\hsize=1.0\hsize}X|}}
\hline
\multicolumn{4}{|>{\centering\arraybackslash\hsize=4.0\hsize}X|}{\textbf{\laghu{ह}\laghu{नु}\laghu{म}\laghu{त}} \par (4)} &
\multicolumn{3}{>{\centering\arraybackslash\hsize=3.0\hsize}X|}{\textbf{\guru{से}\laghu{इ}} \par (3)} &
\multicolumn{3}{>{\centering\arraybackslash\hsize=3.0\hsize}X|}{\textbf{\guru{स}\laghu{र्ब}} \par (3)} &
\multicolumn{2}{>{\centering\arraybackslash\hsize=2.0\hsize}X|}{\textbf{\laghu{सु}\laghu{ख}} \par (2)} &
\multicolumn{4}{>{\centering\arraybackslash\hsize=4.0\hsize}X|}{\textbf{\laghu{क}\laghu{र}\guru{ई}} \par (4)} \\
\hline
\end{tabularx}
\vspace{-1.5em}
\end{table}

\vspace{1em}
\textbf{ \deva{सरल-संस्कृतम्}:}\\
 \deva{मनसि कस्यापि अन्यस्य देवतायाः ध्यानं न धर्तव्यम्।}\\
 \deva{केवलं हनुमतः सेवया (भक्त्या) एव सर्वाणि सुखानि प्राप्यन्ते॥}\\


Let not the mind hold or rely on any other deity.\\
By serving Hanuman alone, one attains all possible happiness.


\vspace{1em}
\newpage
\textbf{\deva{प्रतिपदार्थः} (Meaning of each word):}
\begin{table}[h!]
\centering
\renewcommand{\arraystretch}{1.2}
\begin{tabularx}{\textwidth}{@{} l l X X @{}}
\toprule
\textbf{Awadhi} & \textbf{Sanskrit} & \textbf{English} & \textbf{Grammar \& Notes} \\
\midrule
\deva{और} / \telu{ఔర}  &  \deva{अन्य / परः}  &  Other  &   \\
\deva{देवता} / \telu{దేవతా}  &  \deva{देवता}  &  Deity / Gods  &   \\
\deva{चित्त} / \telu{చిత్త}  &  \deva{चित्ते / मनसि}  &  In the mind  &   \\
\deva{न धरई} / \telu{న ధరఈ}  &  \deva{न धरति}  &  Does not hold/keep  & `y` $\rightarrow$ `ee` verb \\ 
\deva{हनुमंत} / \telu{హనుమంత}  &  \deva{हनुमन्तम्}  &  Hanuman  &   \\
\deva{सेइ} / \telu{సేఇ}  &  \deva{सेवया / आराध्य}  &  By serving  & Absolutive `-i` ending \\
\deva{सर्ब सुख} \gotchaicon / \telu{సర్బ సుఖ}  &  \deva{सर्व-सुखानि}  &  All happiness  &  \\ 
\deva{करई} / \telu{కరఈ}  &  \deva{करोति}  &  Achieves / Does  & `y` $\rightarrow$ `ee` verb \\ 
\bottomrule
\end{tabularx}
\end{table}

\begin{tcolorbox}[poetrybox]
\textbf{\deva{सर्व} $\rightarrow$ \deva{सर्ब} \gotchaicon\ — The Invisible Case Agreement:}\\
Two shifts happen simultaneously in \deva{सर्ब सुख}. First, the familiar \textbf{\deva{व} $\rightarrow$ \deva{ब} shift} turns \deva{सर्व} into \deva{सर्ब}. 
Second, and this is the real gotcha, in classical Sanskrit, \deva{सर्व} must agree in gender, number, and case with the noun it modifies, if they appear as separate words instead of in a compund word like \deva{सर्वसुखसंपन्न}. 
Since \deva{सुख} (happiness) is a \textit{neuter plural}, Sanskrit demands \deva{सर्वाणि सुखानि}.
Awadhi simply ignores this entirely! \deva{सर्ब} is used as a frozen, indeclinable adjective, it never changes its form regardless of the noun's gender or case. 
This ''case-dropping'' is one of the most significant structural differences between Sanskrit and Awadhi.
\end{tcolorbox}

\newpage


% ----- CHAUPAI 36 -----

\newpage

\subsection*{\deva{षट्त्रिंशत्तम चौपाई} (Thirty-Sixth Chaupai)}
\addcontentsline{toc}{subsection}{\deva{षट्त्रिंशत्तम चौपाई} (Thirty-Sixth Chaupai) :  \deva{संकट कटै...}}

\begin{minipage}[t]{0.55\textwidth}
\vspace{0pt}

{\large \linespread{1.5}\selectfont
\deva{संकट कटै मिटै सब पीरा।}\\
\deva{जो सुमिरै हनुमत बलबीरा॥}
\par}

\vspace{0.5em}

{\large \linespread{1.5}\selectfont
\telu{సంకట కటై మిటై సబ పీరా।}\\
\telu{జో సుమిరై హనుమత బలబీరా॥}
\par}

\vspace{0.5em}

{\large \linespread{1.3}\selectfont
\textit{saṅkaṭa kaṭai miṭai saba pīrā |}\\
\textit{jo sumirai hanumata balabīrā ||}
\par}
\end{minipage}%
\hfill
\begin{minipage}[t]{0.42\textwidth}
\vspace{0pt}
\begin{tcolorbox}[anchorbox]
\textbf{Converting Constraints to Opportunities:} In the *Valmiki Ramayana* (*Sundara Kanda, 48:45-50*), when bound by Indrajit's weapon, Hanuman intentionally chose not to break free immediately. He realized that being bound allowed him to be taken directly to Ravana's court, giving him the perfect opportunity to meet the enemy king face-to-face and deliver a message of peace and warning. It teaches us to convert constraints into opportunities.
\vspace{0.5em}\newline Hanuman strategically accepted his temporary captivity, converting a severe constraint into an opportunity to face the enemy king directly.\newline\null\hfill\href{https://www.valmiki.iitk.ac.in/sloka?field_kanda_tid=5&language=dv&field_sarga_value=48}{\textit{Sundara Kanda, 48:45-50}}
\end{tcolorbox}
\end{minipage}

\vspace{0.5em}
\subsubsection*{\textbf{Visual Rhythm Map:}}

\begin{table}[h!]
\centering
\renewcommand{\arraystretch}{1.2}
{%
\setlength{\tabcolsep}{0pt}%
\begin{tabularx}{\textwidth}{|*{16}{>{\centering\arraybackslash\hsize=1.0\hsize}X|}}
\hline
\multicolumn{4}{|>{\centering\arraybackslash\hsize=4.0\hsize}X|}{\textbf{\guru{सं}\laghu{क}\laghu{ट}} \par (4)} & \multicolumn{3}{>{\centering\arraybackslash\hsize=3.0\hsize}X|}{\textbf{\laghu{क}\guru{टै}} \par (3)} & \multicolumn{3}{>{\centering\arraybackslash\hsize=3.0\hsize}X|}{\textbf{\laghu{मि}\guru{टै}} \par (3)} & \multicolumn{2}{>{\centering\arraybackslash\hsize=2.0\hsize}X|}{\textbf{\laghu{स}\laghu{ब}} \par (2)} & \multicolumn{4}{>{\centering\arraybackslash\hsize=4.0\hsize}X|}{\textbf{\guru{पी}\guru{रा}} \par (4)} \\
\hline
\noalign{\vspace{6pt}}
\hline
\multicolumn{2}{|>{\centering\arraybackslash\hsize=2.0\hsize}X|}{\textbf{\guru{जो}} \par (2)} & \multicolumn{4}{>{\centering\arraybackslash\hsize=4.0\hsize}X|}{\textbf{\laghu{सु}\laghu{मि}\guru{रै}} \par (4)} & \multicolumn{4}{>{\centering\arraybackslash\hsize=4.0\hsize}X|}{\textbf{\laghu{ह}\laghu{नु}\laghu{म}\laghu{त}} \par (4)} & \multicolumn{6}{>{\centering\arraybackslash\hsize=6.0\hsize}X|}{\textbf{\laghu{ब}\laghu{ल}\guru{बी}\guru{रा}} \par (6)} \\
\hline
\end{tabularx}%
}%
\end{table}




\vspace{1em}
\textbf{ \deva{सरल-संस्कृतम्}:}\\
 \deva{तस्य सर्वाणि सङ्कटानि छिद्यन्ते, सर्वाः पीडाः च नश्यन्ति}\\
 \deva{यः बलवीरं (महाबल-शालिनं) हनूमन्तं स्मरति॥}\\


All crises are cut away, and all pain is erased for whoever remembers the mighty and brave Hanuman.


\vspace{1em}
\newpage
\textbf{\deva{प्रतिपदार्थः} (Meaning of each word):}
\begin{table}[h!]
\centering
\renewcommand{\arraystretch}{1.2}
\begin{tabularx}{\textwidth}{@{} l l X X @{}}
\toprule
\textbf{Awadhi} & \textbf{Sanskrit} & \textbf{English} & \textbf{Grammar \& Notes} \\
\midrule
\deva{संकट कटै} / \telu{సంకట కటై} / \textit{saṅkaṭa kaṭai}  &  \deva{सङ्कटं छिद्यते}  &  Crisis is cut  & Present tense `-ai` \\
\deva{मिटै} / \telu{మిటై} / \textit{miṭai}  &  \deva{मृज्यते (नश्यति)}  &  Is erased/mitigated  &   \\
\deva{सब पीरा} / \telu{సబ పీరా} / \textit{saba pīrā}  &  \deva{सर्वाः पीडाः}  &  All pain  & \deva{पीडा $\rightarrow$ पीरा} \\ 
\deva{जो सुमिरै} / \telu{జో సుమిరై} / \textit{jo sumirai}  &  \deva{यः स्मरति}  &  Whoever remembers  & Swara-Bhakti (Smara $\rightarrow$ Sumira) \\ 
\deva{हनुमत} / \telu{హనుమత} / \textit{hanumata}  &  \deva{हनुमन्तम्}  &  Hanuman  &   \\
\deva{बलबीरा} \gotchaicon / \telu{బలబీరా} / \textit{balabīrā}  &  \deva{बलवीरम्}  &  Mighty and Brave  &  \\ 
\bottomrule
\end{tabularx}
\end{table}

\begin{tcolorbox}[poetrybox]
\textbf{The Compound \deva{बलवीर} \gotchaicon:}\\
The Sanskrit original is \deva{बलवीरम्} — a \textit{Dvandva} compound meaning ''one who possesses both \deva{बल} (raw physical strength) and \deva{वीर} (heroic courage).'' 
Tulsidas uses it to make a point: Hanuman is not merely a brute-force warrior (\textit{Bala}), 
nor merely a courageous one (\textit{Vira}), 
but a rare soul in whom both qualities are perfectly fused.\\[0.4em]
The familiar \textbf{\deva{व} $\rightarrow$ \deva{ब} shift} then converts \deva{वीर} into \deva{बीर} in Awadhi, giving us \deva{बलबीरा}. 
The final \deva{-आ} is a standard Awadhi masculine rhyming inflection to complete the 6-matra cadence.
\end{tcolorbox}

\newpage


% ----- CHAUPAI 37 -----

\newpage

\subsection*{\deva{सप्तत्रिंशत्तम चौपाई} (Thirty-Seventh Chaupai)}
\addcontentsline{toc}{subsection}{\deva{सप्तत्रिंशत्तम चौपाई} (Thirty-Seventh Chaupai) :  \deva{जै जै...}}

\begin{minipage}[t]{0.55\textwidth}
\vspace{0pt}

{\large \linespread{1.5}\selectfont
\deva{जै जै जै हनुमान गोसाईं।}\\
\deva{कृपा करहु गुरु देव की नाईं॥}
\par}

\vspace{0.5em}

{\large \linespread{1.5}\selectfont
\telu{జై జై జై హనుమాన గోసాయీఁ।}\\
\telu{కృపా కరహు గురు దేవ కీ నాయీఁ॥}
\par}

\vspace{0.5em}

{\large \linespread{1.3}\selectfont
\textit{jai jai jai hanumāna gosāī̃ |}\\
\textit{kṛpā karahu guru deva kī nāī̃ ||}
\par}
\end{minipage}%
\hfill
\begin{minipage}[t]{0.42\textwidth}
\vspace{0pt}
\begin{tcolorbox}[anchorbox]
\textbf{The Ultimate Mentor:} By accepting an ideal role model as our guide (\textit{Gurudeva}), we are asking for the unvarnished truth and the strict discipline required to elevate our lives to their highest potential. A true mentor challenges us to grow.
\vspace{0.5em}\newline By accepting an ideal role model, we invite the strict discipline needed to universally elevate our own potential.\newline\null\hfill\href{https://www.valmiki.iitk.ac.in/sloka?field_kanda_tid=5&language=dv&field_sarga_value=1}{\textit{Sundara Kanda, 1:205}}
\end{tcolorbox}
\end{minipage}

\vspace{0.5em}
\subsubsection*{\textbf{Visual Rhythm Map:}}
\begin{table}[h!]
\centering
\renewcommand{\arraystretch}{1.2}
\noindent\makebox[\textwidth][l]{%
  {%
  \setlength{\tabcolsep}{0pt}%
  \begin{tabularx}{1.0625\textwidth}{|*{17}{>{\centering\arraybackslash\hsize=1.0\hsize}X|}}
  \hline
  \multicolumn{2}{|>{\centering\arraybackslash\hsize=2.0\hsize}X|}{\textbf{\guru{जै}} \par (2)} &
  \multicolumn{2}{>{\centering\arraybackslash\hsize=2.0\hsize}X|}{\textbf{\guru{जै}} \par (2)} &
  \multicolumn{2}{>{\centering\arraybackslash\hsize=2.0\hsize}X|}{\textbf{\guru{जै}} \par (2)} &
  \multicolumn{5}{>{\centering\arraybackslash\hsize=5.0\hsize}X|}{\textbf{\laghu{ह}\laghu{नु}\guru{मा}\laghu{न}} \par (5)} &
  \multicolumn{6}{>{\centering\arraybackslash\hsize=6.0\hsize}X|}{\textbf{\guru{गो}\guru{सा}\guru{ईं}} \par (6)} \\
  \hline
  \end{tabularx}%
  }%
  \hspace{0.01\textwidth}%
  \begin{minipage}{0.1\textwidth}
  \vspace{-0.5em}
  \centering
  {\Huge \textbf{17}}
  \end{minipage}%
}
\vspace{-1.5em}
\end{table}

\begin{table}[h!]
\centering
\renewcommand{\arraystretch}{1.2}
\noindent\makebox[\textwidth][l]{%
  {%
  \setlength{\tabcolsep}{0pt}%
  \begin{tabularx}{1.0625\textwidth}{|*{17}{>{\centering\arraybackslash\hsize=1.0\hsize}X|}}
  \hline
  \multicolumn{3}{|>{\centering\arraybackslash\hsize=3.0\hsize}X|}{\textbf{\laghu{कृ}\guru{पा}} \par (3)} &
  \multicolumn{3}{>{\centering\arraybackslash\hsize=3.0\hsize}X|}{\textbf{\laghu{क}\laghu{र}\laghu{हु}} \par (3)} &
  \multicolumn{2}{>{\centering\arraybackslash\hsize=2.0\hsize}X|}{\textbf{\laghu{गु}\laghu{रु}} \par (2)} &
  \multicolumn{3}{>{\centering\arraybackslash\hsize=3.0\hsize}X|}{\textbf{\guru{दे}\laghu{व}} \par (3)} &
  \multicolumn{2}{>{\centering\arraybackslash\hsize=2.0\hsize}X|}{\textbf{\guru{की}} \par (2)} &
  \multicolumn{4}{>{\centering\arraybackslash\hsize=4.0\hsize}X|}{\textbf{\guru{ना}\guru{ईं}} \par (4)} \\
  \hline
  \end{tabularx}%
  }%
  \hspace{0.01\textwidth}%
  \begin{minipage}{0.1\textwidth}
  \vspace{-0.5em}
  \centering
  {\Huge \textbf{17}}
  \end{minipage}%
}
\vspace{-1.5em}
\end{table}


\vspace{1em}
\textbf{ \deva{सरल-संस्कृतम्}:}\\
 \deva{हे गोस्वामिन् (इन्द्रियाणां स्वामिन्) हनुमन्! भवतः जयः, जयः, जयः अस्तु।}\\
 \deva{भवान् मम उपरि इष्ट-गुरुदेवः इव कृपां करोतु॥}\\


Hail, Hail, All Hail to you, Lord Hanuman, Master of the Senses!\\
Please bestow your grace upon me, just like a Supreme Guru.


\vspace{1em}
\newpage
\textbf{\deva{प्रतिपदार्थः} (Meaning of each word):}
\begin{table}[h!]
\centering
\renewcommand{\arraystretch}{1.2}
\begin{tabularx}{\textwidth}{@{} l l X X @{}}
\toprule
\textbf{Awadhi} & \textbf{Sanskrit} & \textbf{English} & \textbf{Grammar \& Notes} \\
\midrule
\deva{जै जै जै} / \telu{జై జై జై}  &  \deva{जय, जय, जय}  &  Victory / Hail (x3)  &   \\
\deva{हनुमान} / \telu{హనుమాన}  &  \deva{हनुमन्}  &  Hanuman  &   \\
\deva{गोसाईं} / \telu{గోసాయీఁ}  &  \deva{गोस्वामिन्}  &  Lord of Senses  & Go (Senses) + Swami (Lord) \\
\deva{कृपा करहु} / \telu{కృపా కరహు}  &  \deva{कृपां करोतु}  &  Have mercy / grace  & Imperative `-hu` suffix \\
\deva{गुरु देव} / \telu{గురు దేవ}  &  \deva{गुरूदेवः}  &  Supreme Guru  &   \\
\deva{की नाईं} / \telu{కీ నాయీఁ}  &  \deva{इव / सदृशम्}  &  Like / In the manner of  & Pure Awadhi idiom \\
\bottomrule
\end{tabularx}
\end{table}

\begin{tcolorbox}[poetrybox]
\textbf{The 17-Beat Anomaly \& Bhava-Pradhan:}\\
Both lines mathematically sum up to 17 beats! However, because this is a \textit{Bhava Pradhan} poem, the strict 16-beat limit gracefully bows to devotion. The joyful \deva{जय} naturally compresses into an unbroken triplet (\textit{je-je-jai}) in Line 1, and the short 'u' in \deva{करहु} or \deva{गुरु} slips by silently in Line 2 without interrupting the singing flow. These organic inconsistencies are beautiful mnemonic markers that enhance the chanting experience!
\end{tcolorbox}

\newpage


% ----- CHAUPAI 38 -----

\newpage

\subsection*{\deva{अष्टत्रिंशत्तम चौपाई} (Thirty-Eighth Chaupai)}
\addcontentsline{toc}{subsection}{\deva{अष्टत्रिंशत्तम चौपाई} (Thirty-Eighth Chaupai) :  \deva{जो सत...}}

\begin{minipage}[t]{0.55\textwidth}
\vspace{0pt}

{\large \linespread{1.5}\selectfont
\deva{जो सत बार पाठ कर कोई।}\\
\deva{छूटहि बंदि महा सुख होई॥}
\par}

\vspace{0.5em}

{\large \linespread{1.5}\selectfont
\telu{జో సత బార పాఠ కర కోఈ।}\\
\telu{ఛూటహి బంది మహా సుఖ హోఈ॥}
\par}

\vspace{0.5em}

{\large \linespread{1.3}\selectfont
\textit{jo sata bāra pāṭha kara koī |}\\
\textit{chūṭahi bandi mahā sukha hoī ||}
\par}
\end{minipage}%
\hfill
\begin{minipage}[t]{0.42\textwidth}
\vspace{0pt}
\begin{tcolorbox}[anchorbox]
\textbf{The Power of Persistence:} \textit{Sata Baar} (100 times) represents unwavering consistency. In the *Valmiki Ramayana* (*Sundara Kanda, 12:10*), when Hanuman initially fails to find Sita and feels discouraged, he recovers his resolve by declaring: *Anirvedah shriyo mulam* ("Persistence/cheerfulness is the root of all success and prosperity"). He renews his search, proving that mastery and liberation from failure are only built through relentless, repetitive effort.
\vspace{0.5em}\newline Hanuman's ultimate success was built on his absolute refusal to give up, even when his initial efforts yielded no results.\newline\null\hfill\href{https://www.valmiki.iitk.ac.in/sloka?field_kanda_tid=5&language=dv&field_sarga_value=12}{\textit{Sundara Kanda, 12:10}}
\end{tcolorbox}
\end{minipage}

\vspace{0.5em}
\subsubsection*{\textbf{Visual Rhythm Map:}}
\begin{table}[h!]
\centering
\renewcommand{\arraystretch}{1.2}
\setlength{\tabcolsep}{0pt}
\begin{tabularx}{\textwidth}{|*{16}{>{\centering\arraybackslash\hsize=1.0\hsize}X|}}
\hline
\multicolumn{2}{|>{\centering\arraybackslash\hsize=2.0\hsize}X|}{\textbf{\guru{जो}} \par (2)} &
\multicolumn{2}{>{\centering\arraybackslash\hsize=2.0\hsize}X|}{\textbf{\laghu{स}\laghu{त}} \par (2)} &
\multicolumn{3}{>{\centering\arraybackslash\hsize=3.0\hsize}X|}{\textbf{\guru{बा}\laghu{र}} \par (3)} &
\multicolumn{3}{>{\centering\arraybackslash\hsize=3.0\hsize}X|}{\textbf{\guru{पा}\laghu{ठ}} \par (3)} &
\multicolumn{2}{>{\centering\arraybackslash\hsize=2.0\hsize}X|}{\textbf{\laghu{क}\laghu{र}} \par (2)} &
\multicolumn{4}{>{\centering\arraybackslash\hsize=4.0\hsize}X|}{\textbf{\guru{को}\guru{ई}} \par (4)} \\
\hline
\end{tabularx}
\vspace{-1.5em}
\end{table}

\begin{table}[h!]
\centering
\renewcommand{\arraystretch}{1.2}
\setlength{\tabcolsep}{0pt}
\begin{tabularx}{\textwidth}{|*{16}{>{\centering\arraybackslash\hsize=1.0\hsize}X|}}
\hline
\multicolumn{4}{|>{\centering\arraybackslash\hsize=4.0\hsize}X|}{\textbf{\guru{छू}\laghu{ट}\laghu{हि}} \par (4)} &
\multicolumn{3}{>{\centering\arraybackslash\hsize=3.0\hsize}X|}{\textbf{\guru{बं}\laghu{दि}} \par (3)} &
\multicolumn{3}{>{\centering\arraybackslash\hsize=3.0\hsize}X|}{\textbf{\laghu{म}\guru{हा}} \par (3)} &
\multicolumn{2}{>{\centering\arraybackslash\hsize=2.0\hsize}X|}{\textbf{\laghu{सु}\laghu{ख}} \par (2)} &
\multicolumn{4}{>{\centering\arraybackslash\hsize=4.0\hsize}X|}{\textbf{\guru{हो}\guru{ई}} \par (4)} \\
\hline
\end{tabularx}
\vspace{-1.5em}
\end{table}

\vspace{1em}
\textbf{ \deva{सरल-संस्कृतम्}:}\\
 \deva{यः कश्चित् शत-वारम् अस्य पाठं करोति}\\
 \deva{सः बन्धनात् मुच्यते, तस्य महान् सुखं च भवति॥}\\


Whoever recites this one hundred times is freed from all bondage and experiences immense happiness.


\vspace{1em}
\newpage
\textbf{\deva{प्रतिपदार्थः} (Meaning of each word):}
\begin{table}[h!]
\centering
\renewcommand{\arraystretch}{1.2}
\begin{tabularx}{\textwidth}{@{} l l X X @{}}
\toprule
\textbf{Awadhi} & \textbf{Sanskrit} & \textbf{English} & \textbf{Grammar \& Notes} \\
\midrule
\deva{जो कोई} / \telu{జో కోఈ} / \textit{jo koī}  &  \deva{यः कश्चित्}  &  Whoever  &   \\
\deva{सत बार} \gotchaicon / \telu{సత బార} / \textit{sata bāra}  &  \deva{शत-वारम्}  &  100 Times  &  \\ 
\deva{पाठ कर} / \telu{పాఠ కర} / \textit{pāṭha kara}  &  \deva{पाठं करोति}  &  Recites  &   \\
\deva{छूटहि} / \telu{ఛూటహి} / \textit{chūṭahi}  &  \deva{मुच्यते}  &  Is freed / breaks  & Awadhi passive `-hi` \\
\deva{बंदि} / \telu{బంది} / \textit{bandi}  &  \deva{बन्धनात्}  &  From bondage (prison)  &   \\
\deva{महा सुख} / \telu{మహా సుఖ} / \textit{mahā sukha}  &  \deva{महान् सुखम्}  &  Great happiness  &   \\
\deva{होई} / \telu{హోఈ} / \textit{hoī}  &  \deva{भवति}  &  Happens  &   \\
\bottomrule
\end{tabularx}
\end{table}

\begin{tcolorbox}[gotchabox]
\begin{itemize}
    \item \textbf{The \deva{श/ष} $\rightarrow$ \deva{स} Shift:} Yet another perfect example of palatal softening! The strict Sanskrit \deva{शत} (Shata/100) becomes the pure, flowing Awadhi \deva{सत} (Sata).
\end{itemize}
\end{tcolorbox}

\newpage


% ----- CHAUPAI 39 -----

\newpage

\subsection*{\deva{एकोनचत्वारिंशत्तम चौपाई} (Thirty-Ninth Chaupai)}
\addcontentsline{toc}{subsection}{\deva{एकोनचत्वारिंशत्तम चौपाई} (Thirty-Ninth Chaupai) :  \deva{जो यह...}}

\begin{minipage}[t]{0.55\textwidth}
\vspace{0pt}

{\large \linespread{1.5}\selectfont
\deva{जो यह पढ़ै हनुमान चलीसा।}\\
\deva{होय सिद्धि साखी गौरीसा॥}
\par}

\vspace{0.5em}

{\large \linespread{1.5}\selectfont
\telu{జో యహ పఢై హనుమాన చలీసా।}\\
\telu{హోయ సిద్ధి సాఖీ గౌరీసా॥}
\par}

\vspace{0.5em}

{\large \linespread{1.3}\selectfont
\textit{jo yaha paṛhai hanumāna calīsā |}\\
\textit{hoya siddhi sākhī gaurīsā ||}
\par}
\end{minipage}%
\hfill
\begin{minipage}[t]{0.42\textwidth}
\vspace{0pt}
\begin{tcolorbox}[anchorbox]
\textbf{The Universal Guarantee:} When you fully commit to the path of discipline, hard work, and good character, the universe itself aligns to ensure your ultimate success (\textit{Hoya Siddhi}). Consistent effort always pays off in the end.
\vspace{0.5em}\newline When you commit to character and hard work, the universe aligns the highest authorities to guarantee your success.\newline\null\hfill\href{https://www.valmiki.iitk.ac.in/sloka?field_kanda_tid=6&language=dv&field_sarga_value=119}{\textit{Yuddha Kanda, 119:1-8}}
\end{tcolorbox}
\end{minipage}

\vspace{0.5em}
\subsubsection*{\textbf{Visual Rhythm Map:}}
\begin{table}[h!]
\centering
\renewcommand{\arraystretch}{1.2}
\noindent\makebox[\textwidth][l]{%
  {%
  \setlength{\tabcolsep}{0pt}%
  \begin{tabularx}{1.125\textwidth}{|*{18}{>{\centering\arraybackslash\hsize=1.0\hsize}X|}}
  \hline
  \multicolumn{2}{|>{\centering\arraybackslash\hsize=2.0\hsize}X|}{\textbf{\guru{जो}} \par (2)} &
  \multicolumn{2}{>{\centering\arraybackslash\hsize=2.0\hsize}X|}{\textbf{\laghu{य}\laghu{ह}} \par (2)} &
  \multicolumn{3}{>{\centering\arraybackslash\hsize=3.0\hsize}X|}{\textbf{\laghu{प}\guru{ढै}} \par (3)} &
  \multicolumn{5}{>{\centering\arraybackslash\hsize=5.0\hsize}X|}{\textbf{\laghu{ह}\laghu{नु}\guru{मा}\laghu{न}} \par (5)} &
  \multicolumn{5}{>{\centering\arraybackslash\hsize=5.0\hsize}X|}{\textbf{\laghu{च}\guru{ली}\guru{सा}} \par (5)} \\
  \hline
  \end{tabularx}%
  }%
  \hspace{0.01\textwidth}%
  \begin{minipage}{0.1\textwidth}
  \vspace{-0.5em}
  \centering
  {\Huge \textbf{17}}
  \end{minipage}%
}
\vspace{-1.5em}
\end{table}

\begin{table}[h!]
\centering
\renewcommand{\arraystretch}{1.2}
\setlength{\tabcolsep}{0pt}
\begin{tabularx}{\textwidth}{|*{16}{>{\centering\arraybackslash\hsize=1.0\hsize}X|}}
\hline
\multicolumn{3}{|>{\centering\arraybackslash\hsize=3.0\hsize}X|}{\textbf{\guru{हो}\laghu{य}} \par (3)} &
\multicolumn{3}{>{\centering\arraybackslash\hsize=3.0\hsize}X|}{\textbf{\guru{सि}\laghu{द्धि}} \par (3)} &
\multicolumn{4}{>{\centering\arraybackslash\hsize=4.0\hsize}X|}{\textbf{\guru{सा}\guru{खी}} \par (4)} &
\multicolumn{6}{>{\centering\arraybackslash\hsize=6.0\hsize}X|}{\textbf{\guru{गौ}\guru{री}\guru{सा}} \par (6)} \\
\hline
\end{tabularx}
\vspace{-1.5em}
\end{table}

\vspace{1em}
\textbf{ \deva{सरल-संस्कृतम्}:}\\
 \deva{यः एतत् "हनुमान-चालीसा" स्तोत्रं पठति}\\
 \deva{सः सिद्धिं प्राप्नोति, अत्र गौरीशः (शिवः) एव साक्षी अस्ति॥}\\


Whoever reads this Hanuman Chalisa will attain perfection and success.
Lord Shiva (husband of Gauri) is witness to this!

\vspace{1em}
\newpage
\textbf{\deva{प्रतिपदार्थः} (Meaning of each word):}
\begin{table}[h!]
\centering
\renewcommand{\arraystretch}{1.2}
\begin{tabularx}{\textwidth}{@{} l l X X @{}}
\toprule
\textbf{Awadhi} & \textbf{Sanskrit} & \textbf{English} & \textbf{Grammar \& Notes} \\
\midrule
\deva{जो यह} / \telu{జో యహ}  &  \deva{यः एतत्}  &  Whoever this  &   \\
\deva{पढ़ै} / \telu{పఢై}  &  \deva{पठति}  &  Reads  & Verb ending in `-ai` \\
\deva{हनुमान चलीसा} / \telu{హనుమాన చలీసా}  &  \deva{हनुमान-चालीसा}  &  Forty verses of Hanuman  &   \\
\deva{होय सिद्धि} / \telu{హోయ సిద్ధి}  &  \deva{सिद्धिः भवति}  &  Success/Perfection happens  &   \\
\deva{साखी} \gotchaicon / \telu{సాఖీ}  &  \deva{साक्षी}  &  Witness  & \\ 
\deva{गौरीसा} / \telu{గౌరీసా}  &  \deva{गौरीशः (शिवः)}  &  Lord of Gauri  & \\ 
\bottomrule
\end{tabularx}
\end{table}

\begin{tcolorbox}[gotchabox]
\begin{itemize}
    \item \textbf{The \deva{क्ष} $\rightarrow$ \deva{ख} Consonant Shift:} The majestic Sanskrit word for Witness is \deva{साक्षी} (Saakshi). Awadhi completely shatters the harsh `kshi` conjunct into the smooth and aspirated \deva{साखी} (Saakhee). 
\end{itemize}
\end{tcolorbox}

\newpage


% ----- CHAUPAI 40 -----

\newpage

\subsection*{\deva{चत्वारिंशत्तम चौपाई} (Fortieth Chaupai)}
\addcontentsline{toc}{subsection}{\deva{चत्वारिंशत्तम चौपाई} (Fortieth Chaupai) :  \deva{तुलसीदास सदा...}}

\begin{minipage}[t]{0.55\textwidth}
\vspace{0pt}

{\large \linespread{1.5}\selectfont
\deva{तुलसीदास सदा हरि चेरा।}\\
\deva{कीजै नाथ हृदय महँ डेरा॥}
\par}

\vspace{0.5em}

{\large \linespread{1.5}\selectfont
\telu{తులసీదాస సదా హరి చేరా।}\\
\telu{కీజై నాథ హృదయ మహఁ డేరా॥}
\par}

\vspace{0.5em}

{\large \linespread{1.3}\selectfont
\textit{tulasīdāsa sadā hari cerā |}\\
\textit{kījai nātha hṛdaya mahã ḍerā ||}
\par}
\end{minipage}%
\hfill
\begin{minipage}[t]{0.42\textwidth}
\vspace{0pt}
\begin{tcolorbox}[anchorbox]
\textbf{The Humble Heart:} Instead of demanding massive recognition or building monuments to ego, true greatness and peace begin with an unassuming, hospitable, and humble mind (\textit{Hriday Mahan Dera}).
\vspace{0.5em}\newline True greatness isn't found in massive monuments to the ego, but in building an unassuming, hospitable, and humble mind.\newline\null\hfill\href{https://www.valmiki.iitk.ac.in/sloka?field_kanda_tid=6&language=dv&field_sarga_value=128}{\textit{Yuddha Kanda, 128:83}}
\end{tcolorbox}

\end{minipage}

\vspace{0.5em}
\subsubsection*{\textbf{Visual Rhythm Map:}}
\begin{table}[h!]
\centering
\renewcommand{\arraystretch}{1.2}
\setlength{\tabcolsep}{0pt}
\begin{tabularx}{\textwidth}{|*{16}{>{\centering\arraybackslash\hsize=1.0\hsize}X|}}
\hline
\multicolumn{7}{|>{\centering\arraybackslash\hsize=7.0\hsize}X|}{\textbf{\laghu{तु}\laghu{ल}\guru{सी}\guru{दा}\laghu{स}} \par (7)} &
\multicolumn{3}{>{\centering\arraybackslash\hsize=3.0\hsize}X|}{\textbf{\laghu{स}\guru{दा}} \par (3)} &
\multicolumn{2}{>{\centering\arraybackslash\hsize=2.0\hsize}X|}{\textbf{\laghu{ह}\laghu{रि}} \par (2)} &
\multicolumn{4}{>{\centering\arraybackslash\hsize=4.0\hsize}X|}{\textbf{\guru{चे}\guru{रा}} \par (4)} \\
\hline
\end{tabularx}
\vspace{-1.5em}
\end{table}

\begin{table}[h!]
\centering
\renewcommand{\arraystretch}{1.2}
\setlength{\tabcolsep}{0pt}
\begin{tabularx}{\textwidth}{|*{16}{>{\centering\arraybackslash\hsize=1.0\hsize}X|}}
\hline
\multicolumn{4}{|>{\centering\arraybackslash\hsize=4.0\hsize}X|}{\textbf{\guru{की}\guru{जै}} \par (4)} &
\multicolumn{3}{>{\centering\arraybackslash\hsize=3.0\hsize}X|}{\textbf{\guru{ना}\laghu{थ}} \par (3)} &
\multicolumn{3}{>{\centering\arraybackslash\hsize=3.0\hsize}X|}{\textbf{\laghu{हृ}\laghu{द}\laghu{य}} \par (3)} &
\multicolumn{2}{>{\centering\arraybackslash\hsize=2.0\hsize}X|}{\textbf{\laghu{म}\laghu{हँ}} \par (2)} &
\multicolumn{4}{>{\centering\arraybackslash\hsize=4.0\hsize}X|}{\textbf{\guru{डे}\guru{रा}} \par (4)} \\
\hline
\end{tabularx}
\vspace{-1.5em}
\end{table}

\vspace{1em}
\textbf{ \deva{सरल-संस्कृतम्}:}\\
 \deva{हे नाथ! तुलसीदासः सर्वदा भगवतः हरिः (श्रीरामस्य) परिचारकः (सेवकः) अस्ति।}\\
 \deva{अतः हे स्वामिन्! भवान् मम हृदये एव स्वकीयम् आवासं (डेरा) करोतु॥}\\


Tulsidas is, and will always remain, the servant of Lord Hari (God).\\
Therefore, O Lord! Please make my heart your permanent dwelling.


\vspace{1em}
\newpage
\textbf{\deva{प्रतिपदार्थः} (Meaning of each word):}
\begin{table}[h!]
\centering
\renewcommand{\arraystretch}{1.2}
\begin{tabularx}{\textwidth}{@{} l l X X @{}}
\toprule
\textbf{Awadhi} & \textbf{Sanskrit} & \textbf{English} & \textbf{Grammar \& Notes} \\
\midrule
\deva{तुलसीदास} / \telu{తులసీదాస}  &  \deva{तुलसीदासः}  &  Tulsidas (The Author)  &   \\
\deva{सदा} / \telu{సదా}  &  \deva{सर्वदा}  &  Always  &   \\
\deva{हरि चेरा} / \telu{హరి చేరా}  &  \deva{हरेः सेवकः}  &  Servant of God  & `Chera` is pure Awadhi for servant \\
\deva{कीजै} / \telu{కీజై}  &  \deva{कुर्यात् / करोतु}  &  Please do / Make  & Respectful Imperative \\
\deva{नाथ} / \telu{నాథ}  &  \deva{नाथ}  &  O Lord  &   \\
\deva{हृदय महँ} / \telu{హృదయ మహఁ}  &  \deva{हृदय-मध्ये}  &  In the heart  & Locative postposition `mahan` \\
\deva{डेरा} / \telu{డేరా}  &  \deva{आवासम्}  &  Camp / Dwelling  &   \\
\bottomrule
\end{tabularx}
\end{table}
\vspace{0.5em}
\begin{tcolorbox}[poetrybox]
\textbf{Praasa (The Signature Verse):}\\
It is customary in Indian classical poetry for the author to gracefully 'sign' their work in the final verse. Here, the great classical poet Tulsidas identifies himself simply as a proud servant (\deva{चेरा}) of Hari. 
\end{tcolorbox}
\newpage



% ==========================================
\section*{Group 8: The Final Prayer}
\addcontentsline{toc}{section}{Group 8: The Final Prayer}
% ==========================================

% ----- CONCLUDING DOHA -----
\subsection*{\deva{अन्तिम दोहा} (Concluding Doha)}
\addcontentsline{toc}{subsection}{\deva{अन्तिम दोहा} (Concluding Doha) — \deva{पवन तनय...}}

\begin{minipage}[t]{0.55\textwidth}
\vspace{0pt}

{\large \linespread{1.5}\selectfont
\makebox[\linewidth][l]{\deva{पवन तनय संकट हरन}}\\
\makebox[\linewidth][c]{\deva{मंगल मूरति रूप।}}\\
\makebox[\linewidth][l]{\deva{राम लखन सीता सहित }}\\
\makebox[\linewidth][c]{\deva{हृदय बसहु सुर भूप॥}}
\par}

\vspace{0.5em}

{\large \linespread{1.5}\selectfont
\makebox[\linewidth][l]{\telu{పవన తనయ సంకట హరన} }\\
\makebox[\linewidth][c]{\telu{మంగల మూరతి రూప।}}\\
\makebox[\linewidth][l]{\telu{రామ లఖన సీతా సహిత}}\\
\makebox[\linewidth][c]{\telu{హృదయ బసహు సుర భూప॥}}
\par}

\vspace{0.5em}

{\large \linespread{1.3}\selectfont
\textit{pavana tanaya saṅkaṭa harana, maṅgala mūrati rūpa |}\\
\textit{rāma lakhana sītā sahita, hṛdaya basahu sura bhūpa ||}
\par}
\end{minipage}%
\hfill
\begin{minipage}[t]{0.42\textwidth}
\vspace{0pt}
\begin{tcolorbox}[anchorbox]
\textbf{The Crisis Preventer (Strategic Leadership):} Before returning to Ayodhya, Lord Rama sent Hanuman ahead to meet Bharata and gauge his true feelings. If Bharata had grown accustomed to power, Rama's sudden return could cause a political crisis or emotional distress. Hanuman acted with supreme diplomacy, emotional intelligence, and leadership, ensuring a seamless, joyful transition of power and preventing any potential crisis.
\vspace{0.5em}\newline Hanuman's strategic foresight and crisis leadership resolved a highly delicate situation before it could escalate.\newline\null\hfill\href{https://www.valmiki.iitk.ac.in/sloka?field_kanda_tid=6&language=dv&field_sarga_value=125}{\textit{Yuddha Kanda, 125}}
\end{tcolorbox}
\end{minipage}

\vspace{0.5em}
\subsubsection*{\textbf{Visual Rhythm Map:}}
\begin{table}[h!]
\centering
\renewcommand{\arraystretch}{1.2}
\begin{tabularx}{\textwidth}{|>{\centering\arraybackslash\hsize=0.8750\hsize}X|>{\centering\arraybackslash\hsize=0.8750\hsize}X|>{\centering\arraybackslash\hsize=1.1667\hsize}X|>{\centering\arraybackslash\hsize=0.8750\hsize}X|>{\centering\arraybackslash\hsize=1.1667\hsize}X|>{\centering\arraybackslash\hsize=1.1667\hsize}X|>{\centering\arraybackslash\hsize=0.8750\hsize}X|}
\hline
\textbf{\laghu{प}\laghu{व}\laghu{न}} \par (3) & \textbf{\laghu{त}\laghu{न}\laghu{य}} \par (3) & \textbf{\guru{सं}\laghu{क}\laghu{ट}} \par (4) & \textbf{\laghu{ह}\laghu{र}\laghu{न}} \par (3) & \textbf{\guru{मं}\laghu{ग}\laghu{ल}} \par (4) & \textbf{\guru{मू}\laghu{र}\laghu{ति}} \par (4) & \textbf{\guru{रू}\laghu{प}} \par (3) \\
\hline
\end{tabularx}
\vspace{-1.5em}
\end{table}

\begin{table}[h!]
\centering
\renewcommand{\arraystretch}{1.2}
\begin{tabularx}{\textwidth}{|>{\centering\arraybackslash\hsize=1.0000\hsize}X|>{\centering\arraybackslash\hsize=1.0000\hsize}X|>{\centering\arraybackslash\hsize=1.3333\hsize}X|>{\centering\arraybackslash\hsize=1.0000\hsize}X|>{\centering\arraybackslash\hsize=1.0000\hsize}X|>{\centering\arraybackslash\hsize=1.0000\hsize}X|>{\centering\arraybackslash\hsize=0.6667\hsize}X|>{\centering\arraybackslash\hsize=1.0000\hsize}X|}
\hline
\textbf{\guru{रा}\laghu{म}} \par (3) & \textbf{\laghu{ल}\laghu{ख}\laghu{न}} \par (3) & \textbf{\guru{सी}\guru{ता}} \par (4) & \textbf{\laghu{स}\laghu{हि}\laghu{त}} \par (3) & \textbf{\laghu{हृ}\laghu{द}\laghu{य}} \par (3) & \textbf{\laghu{ब}\laghu{स}\laghu{हु}} \par (3) & \textbf{\laghu{सु}\laghu{र}} \par (2) & \textbf{\guru{भू}\laghu{प}} \par (3) \\
\hline
\end{tabularx}
\vspace{-1.5em}
\end{table}

\vspace{1em}
\textbf{ \deva{सरल-संस्कृतम्}:}\\
 \deva{हे पवनपुत्र! हे सङ्कट-हारक! भवान् मङ्गल-मूर्तेः साक्षात् रूपम् अस्ति।}\\
 \deva{हे देवानां राजन्! भवान् राम-लक्ष्मण-सीताभिः सह सर्वदा मम हृदये वसतु॥}\\


O Son of the Wind, O Destroyer of Crisis! You are the very embodiment of auspiciousness.\\
O King of Gods! Please reside forever in my heart, along with Rama, Lakshmana, and Sita.


\vspace{1em}
\newpage
\textbf{\deva{प्रतिपदार्थः} (Meaning of each word):}
\begin{table}[h!]
\centering
\renewcommand{\arraystretch}{1.2}
\begin{tabularx}{\textwidth}{@{} l l X X @{}}
\toprule
\textbf{Awadhi} & \textbf{Sanskrit} & \textbf{English} & \textbf{Grammar \& Notes} \\
\midrule
\deva{पवन तनय} / \telu{పవన తనయ} / \textit{pavana tanaya}  &  \deva{पवन-तनय}  &  Son of the Wind-God  &   \\
\deva{संकट हरन} / \telu{సంకట హరన} / \textit{saṅkaṭa harana}  &  \deva{सङ्कट-हरण}  &  Destroyer of Crisis  &   \\
\deva{मंगल मूरति} / \telu{మంగల మూరతి} / \textit{maṅgala mūrati}  &  \deva{मङ्गल-मूर्तिः}  &  Auspicious idol/form  &   \\
\deva{रूप} / \telu{రూప} / \textit{rūpa}  &  \deva{रूपम्}  &  Embodiment  &   \\
\deva{राम लखन} / \telu{రామ లఖన} / \textit{rāma lakhana}  &  \deva{राम-लक्ष्मण}  &  Rama and Lakshmana  & \\ 
\deva{सीता सहित} / \telu{సీతా సహిత} / \textit{sītā sahita}  &  \deva{सीतया सह}  &  Along with Sita  &   \\
\deva{हृदय बसहु} / \telu{హృదయ బసహు} / \textit{hṛdaya basahu}  &  \deva{हृदये वसतु}  &  Reside in the heart  & Imperative `-hu` suffix \\
\deva{सुर भूप} / \telu{సుర భూప} / \textit{sura bhūpa}  &  \deva{देवानां राजा / भूपतिः}  &  King of Gods  &   \\
\bottomrule
\end{tabularx}
\end{table}

\vspace{0.5em}
\begin{tcolorbox}[poetrybox]
\textbf{Structural Symmetry:}\\
The Chalisa ends exactly how it began. 
By returning flawlessly to the beautiful 13-11 metric cadence of the \textbf{Doha} format, perfectly bookmarking the 40 marching Chaupais with reverence and praise.
\end{tcolorbox}

\newpage



\newpage
\section*{\deva{परिशिष्ट} (Appendix): Glossary of Divine Terms}
\addcontentsline{toc}{section}{\deva{परिशिष्ट} (Appendix) — Glossary of Divine Terms}

Throughout the Hanuman Chalisa, Tulsidas references several deep cosmological and esoteric lists from Hindu theology. To avoid interrupting the flow of recitation, the detailed definitions of these profound terms are provided below.

\vspace{1em}
\subsection*{\textbf{1. The Eight Siddhis (Ashta-Siddhi)}}
In the 31st Chaupai, Mother Sita blesses Hanuman with the ability to grant the \textit{Ashta-Siddhi} (the eight supernatural Yogic powers). A master of these Siddhis can manipulate the fundamental building blocks of matter and space.

\begin{table}[h!]
\centering
\renewcommand{\arraystretch}{1.3}
\begin{tabularx}{\textwidth}{@{} l l X @{}}
\toprule
\textbf{Siddhi} & \textbf{Sanskrit} & \textbf{Supernatural Ability} \\
\midrule
\textbf{Anima} & \deva{अणिमा} & The power to shrink one's body to the size of an atom or sub-atomic particle (used to secretly enter Lanka). \\
\textbf{Mahima} & \deva{महिमा} & The power to expand one's body to an infinitely large, gargantuan size (used to terrify the demoness Surasa). \\
\textbf{Garima} & \deva{गरिमा} & The power to become infinitely heavy (used when Bhima could not lift even Hanuman's tail). \\
\textbf{Laghima} & \deva{लघिमा} & The power to become almost weightless or lighter than air (used for levitation and swift flight across the ocean). \\
\textbf{Prapti} & \deva{प्राप्ति} & The power to instantly acquire or reach anything anywhere, regardless of physical barriers. \\
\textbf{Prakamya} & \deva{प्राकाम्य} & The power to realize any desire, achieve irresistible willpower, and instantly materialize objects. \\
\textbf{Ishitva} & \deva{ईशित्व} & The power of absolute lordship and supreme dominance over nature and living beings. \\
\textbf{Vashitva} & \deva{वशित्व} & The ultimate power to conquer, subjugate, or mesmerize any entity to follow one's will. \\
\bottomrule
\end{tabularx}
\end{table}

\vspace{1em}
\subsection*{\textbf{2. The Nine Nidhis (Nava-Nidhi)}}
Alongside the 8 Siddhis, Sita also granted Hanuman dominion over the \textit{Nava-Nidhi} (the nine divine, inexhaustible treasures of Kubera, the god of wealth). These are not merely financial riches, but profound spiritual and cosmic abundances.

\begin{table}[h!]
\centering
\renewcommand{\arraystretch}{1.3}
\begin{tabularx}{\textwidth}{@{} l l X @{}}
\toprule
\textbf{Nidhi} & \textbf{Sanskrit} & \textbf{Divine Treasure} \\
\midrule
\textbf{Mahapadma} & \deva{महापद्म} & The ''Great Lotus Casket'' containing immense wealth, generating generations of prosperity and philanthropy. \\
\textbf{Padma} & \deva{पद्म} & The ''Lotus Casket'' offering abundant material riches, sattvic tendencies, and untainted gifts. \\
\textbf{Shankha} & \deva{शङ्ख} & The ''Conch Shell'' bringing self-sustaining wealth, typically enjoyed and isolated by the owner. \\
\textbf{Makara} & \deva{मकर} & The ''Crocodile'' granting martial proficiency, strategic intellect, and dominance in warfare/arms. \\
\textbf{Kacchapa} & \deva{कच्छप} & The ''Tortoise'' establishing wealth accumulated slowly, hidden secretly, and passed down carefully. \\
\textbf{Mukunda} & \deva{मुकुन्द} & The ''Quicksilver'' or ''Kettle'' bringing an aesthetic inclination towards fine arts, music, and royal patronage. \\
\textbf{Kunda} & \deva{कुन्द} & The ''Jasmine'' representing wealth related to morality, spiritual upliftment, and charitable endowments. \\
\textbf{Nila} & \deva{नील} & The ''Sapphire'' giving massive resources related to trade, commerce, and material exchange. \\
\textbf{Kharva} & \deva{खर्व} & The ''Dwarf'' representing wealth intertwined with mixed (both auspicious and complex) desires. \\
\bottomrule
\end{tabularx}
\end{table}

\vspace{1em}
\subsection*{\textbf{3. The Four Yugas (Chatur-Yuga)}}
In the 29th Chaupai, Tulsidas proclaims that Hanuman's glory is famous across all four ages (\textit{Yugas}). According to Vedic cosmology, time is cyclical, divided into four grand planetary epochs.

\begin{table}[h!]
\centering
\renewcommand{\arraystretch}{1.3}
\begin{tabularx}{\textwidth}{@{} l l X @{}}
\toprule
\textbf{Yuga} & \textbf{Sanskrit} & \textbf{Era Characteristics} \\
\midrule
\textbf{Satya Yuga} & \deva{सत्ययुग} & The Age of Truth and absolute Purity. Humanity is deeply spiritual, peaceful, and perfectly completely aligned with Dharma. \\
\textbf{Treta Yuga} & \deva{त्रेतायुग} & The Silver Age. Dharma slightly declines. This is the era in which Lord Rama and Hanuman incarnated to establish righteousness. \\
\textbf{Dvapara Yuga} & \deva{द्वापरयुग} & The Bronze Age. Morality drops to 50\%. This era featured the Mahabharata war and the incarnation of Lord Krishna. \\
\textbf{Kali Yuga} & \deva{कलियुग} & The Iron Age (our current era). An age of deep ignorance, conflict, and spiritual amnesia. Hanuman remains immortal (\textit{Chiranjeevi}) physically present on Earth specifically to guide seekers through this dark age. \\
\bottomrule
\end{tabularx}
\end{table}

\vspace{1em}
\subsection*{\textbf{4. The Four Fruits of Life (Phala Chaari)}}
In the very first opening Doha, Tulsidas prays for the ''four fruits'' (\textit{Phala Chaari}). In Vedic philosophy, these are the \textit{Purusharthas}—the four fundamental grand goals or aims of human life that lead to complete fulfillment.

\begin{table}[h!]
\centering
\renewcommand{\arraystretch}{1.3}
\begin{tabularx}{\textwidth}{@{} l l X @{}}
\toprule
\textbf{Fruit} & \textbf{Sanskrit} & \textbf{Life Goal} \\
\midrule
\textbf{Dharma} & \deva{धर्म} & Righteousness, moral duty, and living in cosmic law and harmony. The foundational aim of a noble life. \\
\textbf{Artha} & \deva{अर्थ} & Material prosperity, wealth, economic security, and physical well-being required to sustain life and family. \\
\textbf{Kama} & \deva{काम} & Desire, pleasure, emotional fulfillment, aesthetics, and the enjoyment of the world in a balanced way. \\
\textbf{Moksha} & \deva{मोक्ष} & The ultimate liberation. Freedom from the cycle of birth and death (\textit{Samsara}) and union with the divine. \\
\bottomrule
\end{tabularx}
\end{table}

\newpage
\phantomsection
\addcontentsline{toc}{section}{References \& Academic Sources}
\section*{References \& Academic Sources}

The educational material, grammatical rules, and narrative anecdotes provided in this guide are deeply rooted in rigorous, authoritative scholarship.

\subsection*{The Valmiki Ramayana}
All character-building anecdotes and contextual stories are strictly sourced from the first six Kandas of the \textbf{Valmiki Ramayana}, intentionally avoiding contested texts or later additions (such as the Uttara Kanda).
\begin{itemize}
    \item \href{https://www.valmiki.iitk.ac.in}{\textbf{IIT Kanpur Valmiki Ramayana Portal}}: Our primary digital repository for exact Shloka verification and English translations.
\end{itemize}

\subsection*{The Ramcharitmanas \& Awadhi Text}
The core Devanagari text and orthography follow the traditional, universally accepted manuscripts.
\begin{itemize}
    \item \href{https://gitapress.org/result?search=Hanuman%20Chalisa&category=}{\textbf{Gita Press, Gorakhpur}}: The gold standard for the source Devanagari and Awadhi text of the Hanuman Chalisa. \href{https://dn710202.ca.archive.org/0/items/xqZf_shri-hanuman-chalisha-gita-press/Shri%20Hanuman%20Chalisha%20-%20Gita%20Press.pdf}{\textbf{PDF from ArXiv.org}}.


    \item \href{https://www.ramcharitmanas.iitk.ac.in/content/tulsidas}{\textbf{IIT Kanpur Ramcharitmanas Portal}}: Used for cross-referencing Tulsidas's broader poetic patterns and lexical choices.
\end{itemize}

\subsection*{Awadhi Grammar \& Linguistics}
The phonetic shift explanations (like the \(v \rightarrow b\) shift) and grammatical notes found in the ''Gotcha'' boxes are derived from foundational academic research in Indo-Aryan linguistics.
\begin{itemize}
    \item \textbf{Baburam Saksena, \href{https://archive.org/details/in.ernet.dli.2015.238311/mode/2up}{\textit{Evolution of Awadhi}}} (1937): The seminal academic work detailing the historical development, phonology, and morphology of the Awadhi language.
    \item \textbf{Rupert Snell, \href{https://ia600500.us.archive.org/16/items/Pushtimarg/BrajBhashaReaderenglish.pdf}{\textit{The Hindi Classical Tradition: A Braj Bhas\={a} Reader}}} (1991): The foundational text for understanding the metrical pulse (Matra/Beats) and the treatment of syllable weights in medieval North Indian poetry.
    \item \textbf{George A. Grierson, \href{https://dsal.uchicago.edu/books/lsi/lsi.php?volume=6&pages=286}{\textit{Linguistic Survey of India}}} (Vol VI): The historical baseline for dialectal mapping and grammatical structures of Eastern Hindi and Awadhi.
\end{itemize}


\end{document}
